% Production et traduction de textes liés à la vie économique — Allemand 1ère A — Cours signé OnBuch+
% Programme officiel MINESEC. Compiler avec : tectonic ALA09-production-et-traduction-de-textes-lies-a-la-vie-economique.tex
\newcommand{\DOCMATIERE}{Allemand}
\newcommand{\DOCNIVEAU}{1ère A}
\newcommand{\DOCMODULE}{Chapitre 2 : Vie économique}
\newcommand{\DOCLECON}{ALA09}
\newcommand{\DOCTITRE}{Production et traduction de textes liés à la vie économique}
% OnBuch+ — préambule « cours signés » ENRICHI (compilé avec tectonic / XeLaTeX)
% Système visuel « Pop » d'OnBuch+ : fond crème (jamais blanc), bordures encre
% épaisses, ombres « dures » décalées, titres Archivo Black en capitales.
% Version MATHÉMATIQUES : graphes (pgfplots/tikz), boîtes pédagogiques
% (définition, propriété, méthode, attention, le savais-tu, exercice résolu,
% exercice, TP, bilan), page de garde et signature OnBuch+ ancrée en bas.
%
% Macros attendues dans main.tex AVANT \input{preamble} :
%   \DOCMATIERE  \DOCNIVEAU  \DOCMODULE  \DOCLECON (numéro)  \DOCTITRE
\documentclass[11pt]{article}

\usepackage{fontspec}
\usepackage[ngerman,french]{babel} % français = langue principale ; l'allemand via \all{..} / textebox
\usepackage[a4paper,margin=2cm,top=3.4cm,bottom=2.8cm,headheight=1.4cm,footskip=1.3cm]{geometry}
\usepackage{amsmath,amssymb,amsfonts}
\usepackage{mathtools} % \xmapsto, \coloneqq... souvent utilisés par les modèles
\usepackage{cancel} % simplifications barrées (bilans de forces)
% \abs{z} (module, valeur absolue) et \norm{u} (norme), utilisés par les modèles
\providecommand{\abs}{}\let\abs\relax\DeclarePairedDelimiter{\abs}{\lvert}{\rvert}
\providecommand{\norm}{}\let\norm\relax\DeclarePairedDelimiter{\norm}{\lVert}{\rVert}
\usepackage[table]{xcolor} % [table] : fournit aussi \rowcolors (alternance de lignes)
\usepackage{pagecolor}
\usepackage{tikz}
\usetikzlibrary{arrows.meta,positioning,calc,shapes.geometric,shapes.misc,shapes.arrows,decorations.pathmorphing,decorations.pathreplacing,decorations.markings,patterns,shadows,mindmap,trees,fit,backgrounds,matrix,intersections,angles,quotes}
\usepackage{pgfplots}
\usepackage[version=4]{mhchem} % équations nucléaires \ce{^{235}_{92}U}
\usepackage[european,siunitx]{circuitikz} % schémas électriques
\pgfplotsset{compat=1.18}
\usepgfplotslibrary{fillbetween,groupplots} % aires entre courbes (calcul intégral)
\usepackage{enumitem}
\usepackage{fancyhdr}
% [most] charge toutes les bibliothèques tcolorbox (listings, minted, raster,
% documentation...) et gonfle inutilement la mémoire TeX sur un document long
% (~300 boîtes) : on ne charge que ce qui est réellement utilisé.
\usepackage[skins,breakable]{tcolorbox}
\usepackage{graphicx}
\usepackage{colortbl}
\usepackage{booktabs}
\usepackage{tabularx}
\usepackage{diagbox} % case d'en-tête diagonale des tableaux à double entrée
% Colonnes extensibles centrée / gauche / droite (C, L, R), souvent utilisées par les modèles
\newcolumntype{C}{>{\centering\arraybackslash}X}
\newcolumntype{L}{>{\raggedright\arraybackslash}X}
\newcolumntype{R}{>{\raggedleft\arraybackslash}X}
\usepackage{array}
\usepackage{multicol}
\usepackage{multirow}
\usepackage{siunitx}
% parse-numbers=false : accepte aussi \SI{2,5 \times 10^{-3}}{...}
\sisetup{locale=FR,output-decimal-marker={,},group-minimum-digits=4,parse-numbers=false}
\DeclareSIUnit\ohm{\ensuremath{\symup{\Omega}}} % U+2126 absent de la police
\DeclareSIUnit\division{div} % réglages d'oscilloscope (ms/div, V/div)
\DeclareSIUnit\tr{tr} % tours (vitesse de rotation, ex. \SI{1500}{\tr\per\minute})
\DeclareSIUnit\molar{M} % concentration molaire (ex. \SI{3}{\molar}, \SI{1}{\micro\molar})
\DeclareSIUnit\an{an} % année (le modèle confond parfois avec \year, un compteur TeX) — ex. \SI{5}{\kilo\meter\per\an}
\DeclareSIUnit\FCFA{FCFA} % franc CFA (ex. \SI{500}{\FCFA\per\kilo\gram})
\DeclareSIUnit\degreeBrix{\textdegree Brix} % teneur en sucre d'un sirop/jus (réfractométrie)
\DeclareSIUnit\month{mois} % le modèle utilise parfois \month, un compteur TeX, comme unité de durée
\usepackage{qrcode}
% \degree : commande fréquemment utilisée par les modèles pour un angle ou une
% température en dehors de \SI{}{} (non fournie par les paquets chargés).
\providecommand{\degree}{^{\circ}}
% \qed (fin de démonstration), utilisé par les modèles sans amsthm
\providecommand{\qed}{\hfill$\square$}
% \barre : double barre des valeurs interdites dans un tableau de variations (tkz-tab)
\providecommand{\barre}{\ensuremath{\|}}
% environnement proof (amsthm non chargé) utilisé par les modèles
\ifdefined\proof\else\newenvironment{proof}[1][Démonstration]{\par\noindent\textit{#1.}\ }{\qed\par}\fi
\usepackage{lastpage}
\usepackage{hyperref}
\hypersetup{hidelinks}

% ============================================================
% Palette « Pop » OnBuch+ — mêmes hex que la classe P (plus_ui.dart)
% ============================================================
\definecolor{poporange}{HTML}{F59321}
\definecolor{popdark}{HTML}{E07A0C}
\definecolor{popcream}{HTML}{FFF6E8}
\definecolor{popink}{HTML}{1C1714}
\definecolor{poppurple}{HTML}{7A5AE0}
\definecolor{popgreen}{HTML}{2E9E6A}
\definecolor{poppink}{HTML}{E0457B}
\definecolor{popblue}{HTML}{2F7DE1}
\definecolor{popgold}{HTML}{C98A12}
\definecolor{popmuted}{HTML}{8A7F76}
\definecolor{popline}{HTML}{EADFD2}
\definecolor{popgray}{HTML}{8A7F76}
\colorlet{popgrey}{popgray}
% Alias tolérants (le modèle nomme parfois les couleurs différemment)
\colorlet{popwhite}{white}
\colorlet{popblack}{popink}
\colorlet{popred}{poppink}
\colorlet{popyellow}{popgold}
% Teintes claires (fonds de tableaux/encadrés) — le modèle nomme parfois
% les couleurs avec un suffixe L (ex. popblueL) sans les avoir définies.
\colorlet{poporangeL}{poporange!20}
\colorlet{popdarkL}{popdark!20}
\colorlet{poppurpleL}{poppurple!20}
\colorlet{popgreenL}{popgreen!20}
\colorlet{poppinkL}{poppink!20}
\colorlet{popblueL}{popblue!20}
\colorlet{popgoldL}{popgold!20}
\colorlet{popredL}{popred!20}
\colorlet{popgrayL}{popgray!20}
% Teintes claires pour fonds de tableaux / zones
\colorlet{popblueL}{popblue!10!white}
\colorlet{poporangeL}{poporange!14!white}
\colorlet{popgreenL}{popgreen!12!white}
\colorlet{poppurpleL}{poppurple!12!white}
\colorlet{poppinkL}{poppink!12!white}
\colorlet{popgoldL}{popgold!14!white}

\pagecolor{popcream}

% ============================================================
% Typographie
% ============================================================
\defaultfontfeatures{Ligatures=TeX}
\setmainfont{PlusJakartaSans-Medium.ttf}[Path=../../../fonts/,BoldFont=PlusJakartaSans-Bold.ttf]
% Maths en sans-serif (Fira Math) avec chiffres et lettres droites de Plus
% Jakarta, pour que les formules se fondent dans le texte.
\usepackage[math-style=ISO,bold-style=ISO]{unicode-math}
\setmathfont{FiraMath-Regular.otf}[Path=../../../fonts/]
\setmathfont{PlusJakartaSans-Medium.ttf}[Path=../../../fonts/,range={up/num,up/{Latin,latin},\mathup}]
\usepackage{icomma} % virgule décimale sans espace en mode math
% Caractères mathématiques Unicode absents des polices de texte (Plus Jakarta,
% Archivo Black) : rendus par leur équivalent mathématique, en texte comme en
% formule — sinon ils s'affichent en carré vide (ex. « Arithmétique dans ℤ »).
\usepackage{newunicodechar}
\newunicodechar{ℝ}{\ensuremath{\mathbb{R}}}
\newunicodechar{ℤ}{\ensuremath{\mathbb{Z}}}
\newunicodechar{ℂ}{\ensuremath{\mathbb{C}}}
\newunicodechar{ℕ}{\ensuremath{\mathbb{N}}}
\newunicodechar{ℚ}{\ensuremath{\mathbb{Q}}}
\newunicodechar{𝔻}{\ensuremath{\mathbb{D}}}
\newunicodechar{α}{\ensuremath{\alpha}}
\newunicodechar{β}{\ensuremath{\beta}}
\newunicodechar{γ}{\ensuremath{\gamma}}
\newunicodechar{δ}{\ensuremath{\delta}}
\newunicodechar{ε}{\ensuremath{\varepsilon}}
\newunicodechar{θ}{\ensuremath{\theta}}
\newunicodechar{λ}{\ensuremath{\lambda}}
\newunicodechar{μ}{\ensuremath{\mu}}
\newunicodechar{σ}{\ensuremath{\sigma}}
\newunicodechar{τ}{\ensuremath{\tau}}
\newunicodechar{φ}{\ensuremath{\varphi}}
\newunicodechar{ψ}{\ensuremath{\psi}}
\newunicodechar{ω}{\ensuremath{\omega}}
\newunicodechar{Σ}{\ensuremath{\Sigma}}
% majuscules produites par \MakeUppercase dans les titres (α-aminés -> Α)
\newunicodechar{Α}{\ensuremath{\alpha}}
\newunicodechar{Β}{\ensuremath{\beta}}
\newunicodechar{Γ}{\ensuremath{\Gamma}}
\newunicodechar{Δ}{\ensuremath{\Delta}}
\newunicodechar{Ε}{\ensuremath{\varepsilon}}
\newunicodechar{Θ}{\ensuremath{\Theta}}
\newunicodechar{Λ}{\ensuremath{\Lambda}}
\newunicodechar{Μ}{\ensuremath{\mu}}
\newunicodechar{Τ}{\ensuremath{\tau}}
\newunicodechar{Φ}{\ensuremath{\Phi}}
\newunicodechar{Ψ}{\ensuremath{\Psi}}
\newunicodechar{Ω}{\ensuremath{\Omega}}
\newunicodechar{↦}{\ensuremath{\mapsto}}
\newunicodechar{∈}{\ensuremath{\in}}
\newunicodechar{∉}{\ensuremath{\notin}}
\newunicodechar{∧}{\ensuremath{\wedge}}
\newunicodechar{⇔}{\ensuremath{\Leftrightarrow}}
\newunicodechar{⟺}{\ensuremath{\Longleftrightarrow}}
\newunicodechar{⇒}{\ensuremath{\Rightarrow}}
\newunicodechar{⟹}{\ensuremath{\Longrightarrow}}
\newunicodechar{∘}{\ensuremath{\circ}}
\newunicodechar{∪}{\ensuremath{\cup}}
\newunicodechar{∩}{\ensuremath{\cap}}
\newunicodechar{≡}{\ensuremath{\equiv}}
\newunicodechar{⊂}{\ensuremath{\subset}}
\newunicodechar{⊥}{\ensuremath{\perp}}
\newunicodechar{∃}{\ensuremath{\exists}}
\newunicodechar{∀}{\ensuremath{\forall}}
\newunicodechar{∛}{\ensuremath{\sqrt[3]{\,}}}
\newunicodechar{∜}{\ensuremath{\sqrt[4]{\,}}}
\newunicodechar{⃗}{\ensuremath{{}^{\to}}}
\newunicodechar{ⁿ}{\ifmmode^{n}\else\textsuperscript{\ensuremath{n}}\fi}
\newunicodechar{ˣ}{\ifmmode^{x}\else\textsuperscript{\ensuremath{x}}\fi}
\newunicodechar{ᵇ}{\ifmmode^{b}\else\textsuperscript{\ensuremath{b}}\fi}
\newunicodechar{ʳ}{\ifmmode^{r}\else\textsuperscript{\ensuremath{r}}\fi}
\newunicodechar{ᵘ}{\ifmmode^{u}\else\textsuperscript{\ensuremath{u}}\fi}
\newunicodechar{ʸ}{\ifmmode^{y}\else\textsuperscript{\ensuremath{y}}\fi}
\newunicodechar{ᵃ}{\ifmmode^{a}\else\textsuperscript{\ensuremath{a}}\fi}
\newunicodechar{ᵉ}{\ifmmode^{e}\else\textsuperscript{\ensuremath{e}}\fi}
\newunicodechar{ᵏ}{\ifmmode^{k}\else\textsuperscript{\ensuremath{k}}\fi}
\newunicodechar{ᵖ}{\ifmmode^{p}\else\textsuperscript{\ensuremath{p}}\fi}
\newunicodechar{ᴺ}{\ifmmode^{N}\else\textsuperscript{\ensuremath{N}}\fi}
\newunicodechar{ᶜ}{\ifmmode^{c}\else\textsuperscript{\ensuremath{c}}\fi}
\newunicodechar{ʰ}{\ifmmode^{h}\else\textsuperscript{\ensuremath{h}}\fi}
\newunicodechar{ᵠ}{\ifmmode^{q}\else\textsuperscript{\ensuremath{q}}\fi}
\newunicodechar{ₙ}{\ifmmode_{n}\else\textsubscript{\ensuremath{n}}\fi}
\newunicodechar{ₐ}{\ifmmode_{a}\else\textsubscript{\ensuremath{a}}\fi}
\newunicodechar{ᵢ}{\ifmmode_{i}\else\textsubscript{\ensuremath{i}}\fi}
\newunicodechar{ᵣ}{\ifmmode_{r}\else\textsubscript{\ensuremath{r}}\fi}
\newunicodechar{ₖ}{\ifmmode_{k}\else\textsubscript{\ensuremath{k}}\fi}
\newunicodechar{ᵦ}{\ifmmode_{\beta}\else\textsubscript{\ensuremath{\beta}}\fi}
\newunicodechar{ₓ}{\ifmmode_{x}\else\textsubscript{\ensuremath{x}}\fi}
\newfontfamily\popdisplay{ArchivoBlack-Regular.ttf}[Path=../../../fonts/]
\newfontfamily\poplogo{SpaceGrotesk-Bold.ttf}[Path=../../../fonts/]
\newfontfamily\popsemi{PlusJakartaSans-SemiBold.ttf}[Path=../../../fonts/]
\newfontfamily\popxbold{PlusJakartaSans-ExtraBold.ttf}[Path=../../../fonts/]
\newfontfamily\popxboldls{PlusJakartaSans-ExtraBold.ttf}[Path=../../../fonts/,LetterSpace=3.0]

\setlist[enumerate]{leftmargin=1.7em,itemsep=5pt,topsep=4pt}
\setlist[itemize]{leftmargin=1.7em,itemsep=4pt,topsep=4pt,label={\color{poporange}\textbullet}}
\setlength{\parindent}{0pt}
\setlength{\parskip}{5pt}
\renewcommand{\arraystretch}{1.25}

% pgfplots : style Pop par défaut
\pgfplotsset{
  every axis/.append style={axis line style={popink,line width=1pt},tick style={popink},
    grid style={popline,line width=0.5pt},label style={font=\small},tick label style={font=\footnotesize}},
  popcurve/.style={poporange,line width=1.8pt,no markers,smooth},
}
% Même style disponible dans un \draw TikZ simple (hors pgfplots)
\tikzset{popcurve/.style={poporange,line width=1.8pt}}
\tikzset{popgrid/.style={popline,very thin}}
\colorlet{popcurve}{poporange} % le nom du style est parfois utilisé comme couleur

% ============================================================
% Pill / squiggle
% ============================================================
\newcommand{\poppill}[3]{%
  \tikz[baseline=-0.55ex]\node[draw=popink,line width=1.2pt,rounded corners=2.6pt,%
    fill=#2,inner xsep=6.5pt,inner ysep=3pt]{%
    \popxboldls\fontsize{7}{7}\selectfont\color{#3}\MakeUppercase{#1}};%
}
\newcommand{\popsquiggle}[1]{%
  \tikz[baseline=-0.4ex]{
    \draw[#1,line width=2.6pt,line cap=round]
      plot [smooth] coordinates {(0,0.09) (0.34,0.01) (0.68,0.21) (1.02,0.01) (1.36,0.10)};
  }%
}
% Mot-clé surligné (marqueur orange sous le texte)
\newcommand{\cle}[1]{{\popxbold\color{popdark}#1}}

% ============================================================
% En-tête / pied
% ============================================================
\pagestyle{fancy}
\fancyhf{}
\renewcommand{\headrulewidth}{0pt}
\renewcommand{\footrulewidth}{0pt}
\lhead{\raisebox{-0.15ex}{\poplogo\fontsize{13}{13}\selectfont\color{popink}OnBuch\color{poporange}+}}
\rhead{\poppill{\DOCMATIERE\ \textbullet\ \DOCNIVEAU\ \textbullet\ Leçon \DOCLECON}{white}{popink}}
\lfoot{\popxbold\fontsize{7.6}{7.6}\selectfont\color{popmuted}\MakeUppercase{Cours signé OnBuch+}\ \textbullet\ {\popsemi\DOCTITRE}}
\rfoot{\tikz[baseline=-0.6ex]\node[rounded corners=8.5pt,fill=popink,minimum height=17pt,minimum width=17pt,inner xsep=5pt,inner sep=0]{\popxbold\fontsize{8}{8}\selectfont\color{popcream}\thepage\,/\,\pageref*{LastPage}};}

% ============================================================
% Titres
% ============================================================
\newcommand{\chaptitre}[2]{%
  \begin{tcolorbox}[enhanced,colback=poporange,colframe=popink,boxrule=2.6pt,arc=11pt,
      drop shadow=popink,left=15pt,right=15pt,top=12pt,bottom=11pt,before skip=3pt,after skip=18pt]
    \poppill{#2}{popink}{popcream}\par\vspace{9pt}
    {\raggedright\hyphenpenalty=10000\exhyphenpenalty=10000\popdisplay\fontsize{22}{24}\selectfont\color{popcream}\MakeUppercase{#1}\par}
  \end{tcolorbox}
}
% Section numérotée : pastille orange avec le numéro + titre Archivo
\newcounter{popsec}
\newcommand{\coursec}[1]{%
  \stepcounter{popsec}\setcounter{popsub}{0}%
  \par\vspace{16pt}\noindent
  \begin{minipage}{\linewidth}
  \tikz[baseline=-0.7ex]\node[draw=popink,line width=1.6pt,fill=poporange,rounded corners=4pt,minimum size=22pt,inner sep=0]{\popdisplay\fontsize{12}{12}\selectfont\color{popcream}\thepopsec};%
  \hspace{8pt}{\popdisplay\fontsize{15}{17}\selectfont\color{popink}\MakeUppercase{#1}}\par
  \vspace{4pt}\noindent{\color{poporange}\rule{36pt}{3.6pt}}
  \end{minipage}\par\nopagebreak\vspace{8pt}
}
\newcounter{popsub}
\newcommand{\courssub}[1]{%
  \stepcounter{popsub}%
  \par\vspace{9pt}\noindent{\popxbold\fontsize{11.5}{13}\selectfont\color{popink}{\color{poporange}\thepopsec.\thepopsub}\hspace{6pt}#1}\par\nopagebreak\vspace{4pt}
}

% ============================================================
% Boîtes pédagogiques — même recette : fond blanc, bordure épaisse colorée,
% ombre dure encre, badge plein. Titre optionnel [#1].
% ============================================================
\tcbset{popbase/.style={breakable,enhanced,colback=white,boxrule=2.3pt,arc=9pt,
  shadow={3pt}{-3pt}{0pt}{popink},left=14pt,right=14pt,top=11pt,bottom=11pt,before skip=11pt,after skip=15pt,
  fontupper=\popsemi\fontsize{10.6}{14.5}\selectfont\color{popink},
  fontlower=\popsemi\fontsize{10.6}{14.5}\selectfont\color{popink}}}
% Coche dessinée (le glyphe ✓ n'existe pas dans les polices du document)
\newcommand{\popcheck}[1]{\tikz[baseline=-0.5ex]\draw[#1,line width=1.6pt,line cap=round,line join=round](-0.13,0.0)--(-0.04,-0.09)--(0.13,0.1);}
\providecommand{\coche}{\popcheck{popgreen}}
\providecommand{\resultat}[1]{\par\smallskip\noindent\fcolorbox{popgreen}{popgreenL}{\parbox{\dimexpr\linewidth-2\fboxsep-2\fboxrule\relax}{\textbf{Résultat.} #1}}\par\smallskip}
\newcommand{\popboxhead}[3]{\poppill{#1}{#2}{white}\ifx\relax#3\relax\else\hspace{6pt}{\popxbold\fontsize{10.6}{12}\selectfont\color{popink}#3}\fi\par\vspace{7pt}}

\newtcolorbox{aretenir}[1][]{popbase,colframe=popgreen,before upper={\popboxhead{À retenir}{popgreen}{#1}}}
% Texte en allemand (document de lecture, dialogue, extrait) : cadre
% d'étude. Un titre optionnel (source, auteur) s'affiche dans l'en-tête.
\newtcolorbox{textebox}[1][]{popbase,colframe=popblue,before upper={\popboxhead{Text}{popblue}{#1}\begin{otherlanguage}{ngerman}},after upper={\end{otherlanguage}}}
% Marques ✓ / ✗ (forme correcte / fautive) souvent utilisées par les modèles
\providecommand{\cross}{\textcolor{poppink}{\ensuremath{\times}}}
\providecommand{\xmark}{\cross}\providecommand{\crossmark}{\cross}\providecommand{\cmark}{\ensuremath{\checkmark}}
% Ligne de réponse / justification à compléter (inventée par les modèles)
\providecommand{\justification}[1]{\par\noindent\textit{Justification~:} \dotfill\par}
% \textipa (package tipa non chargé) : la police ne couvre pas l'API -> simple passage
\providecommand{\textipa}[1]{#1}
% Soulignement (ulem non chargé) : \uline, \ul
\providecommand{\uline}[1]{\underline{#1}}\providecommand{\ul}[1]{\underline{#1}}
% \alert (beamer) inventée par les modèles : texte mis en évidence
\providecommand{\alert}[1]{\textbf{\textcolor{poppink}{#1}}}
% variantes ASCII d'ordinaux écrites en macro
\providecommand{\iere}{\textsuperscript{re}}\providecommand{\ieme}{\textsuperscript{e}}\providecommand{\ere}{\textsuperscript{re}}\providecommand{\eme}{\textsuperscript{e}}\providecommand{\er}{\textsuperscript{er}}\providecommand{\ie}{\textsuperscript{e}}
% Ordinaux français écrits en macro par les modèles : 1\ière, 3\ième
\newcommand{\ière}{\textsuperscript{re}}\newcommand{\ième}{\textsuperscript{e}}
% \exemple (inventée par les modèles) : étiquette « Exemple : »
\providecommand{\exemple}{\textbf{Exemple~:}\ }
% \square utilisable aussi hors mode math (cases à cocher)
\AtBeginDocument{\let\squareMath\square\renewcommand{\square}{\ensuremath{\squareMath}}}
% Mot ou phrase en allemand à l'intérieur d'un paragraphe français
\newcommand{\all}[1]{\ifmmode\text{\textit{#1}}\else\begingroup\selectlanguage{ngerman}\itshape #1\endgroup\fi}
\let\esp\all % alias toléré
\newtcolorbox{exemplebox}[1][]{popbase,colframe=poporange,before upper={\popboxhead{Exemple}{poporange}{#1}}}
\newtcolorbox{definition}[1][]{popbase,colframe=popblue,before upper={\popboxhead{Définition}{popblue}{#1}}}
\newtcolorbox{propriete}[1][]{popbase,colframe=poppurple,before upper={\popboxhead{Propriété}{poppurple}{#1}}}
\newtcolorbox{methode}[1][]{popbase,colframe=popgold,before upper={\popboxhead{Méthode}{popgold}{#1}}}
\newtcolorbox{attention}[1][]{popbase,colframe=poppink,before upper={\popboxhead{Attention}{poppink}{#1}}}
\newtcolorbox{experience}[1][]{popbase,colframe=popblue,colback=popblueL,before upper={\popboxhead{Expérience}{popblue}{#1}}}
% « Le savais-tu ? » : carte violette pleine (contexte, culture, Cameroun)
\newtcolorbox{savaistu}[1][]{popbase,colback=poppurple,colframe=popink,
  fontupper=\popsemi\fontsize{10.4}{14.2}\selectfont\color{white},
  before upper={\poppill{Le savais-tu\,?}{white}{poppurple}\ifx\relax#1\relax\else\hspace{6pt}{\popxbold\color{white}#1}\fi\par\vspace{7pt}}}
% Exercice résolu : énoncé en haut, \tcblower puis la solution sur fond crème
\newcounter{popexo}\newcounter{popexr}
\newtcolorbox{exoresolu}[1][]{popbase,colframe=poporange,colbacklower=popcream,
  segmentation style={popink,line width=1.2pt,dashed},
  before upper={\stepcounter{popexr}\popboxhead{Exercice résolu \thepopexr}{poporange}{#1}},
  before lower={\poppill{Solution}{popink}{popcream}\par\vspace{6pt}}}
% Exercice (non résolu en place) — la correction est dans la partie Corrigés
\newtcolorbox{exercice}[1][]{popbase,colframe=popink,
  before upper={\stepcounter{popexo}\popboxhead{Exercice \thepopexo}{popink}{#1}}}
% Corrigé
\newtcolorbox{corrige}[1][]{popbase,colframe=popgreen,colback=popgreenL,
  before upper={\popboxhead{Corrigé}{popgreen}{#1}}}
% Objectifs / prérequis / bilan
\newtcolorbox{objectifs}{popbase,colframe=popink,colback=poporangeL,
  before upper={\popboxhead{Objectifs de la leçon}{popink}{}}}
\newtcolorbox{prerequis}{popbase,colframe=popink,colback=popblueL,
  before upper={\popboxhead{Prérequis}{popblue}{}}}
\newtcolorbox{bilan}{popbase,colframe=popink,colback=white,boxrule=2.8pt,
  before upper={\popboxhead{Fiche bilan}{poporange}{L'essentiel à connaître pour l'examen}}}
% Figure encadrée avec légende
\newtcolorbox{popfigure}[1][]{enhanced,breakable=false,colback=white,colframe=popline,boxrule=1.4pt,arc=8pt,
  left=10pt,right=10pt,top=10pt,bottom=8pt,before skip=10pt,after skip=12pt,halign=center,
  lower separated=false,halign lower=center,
  fontlower=\popsemi\fontsize{9}{11.5}\selectfont\color{popmuted},
  #1}
\newcommand{\legende}[1]{\par\vspace{6pt}{\popsemi\fontsize{9}{11.5}\selectfont\color{popmuted}#1}}

% ============================================================
% Signature OnBuch+
% ============================================================
\newcommand{\poplogoinline}[1]{{\poplogo\fontsize{#1}{#1}\selectfont\color{popink}OnBuch\color{poporange}+}}
\newcommand{\signatureonbuch}{%
  \begin{tcolorbox}[enhanced,colback=popink,colframe=popink,boxrule=2.2pt,arc=10pt,
      drop shadow=poporange,left=14pt,right=14pt,top=10pt,bottom=10pt,before skip=0pt,after skip=0pt]
    \tikz[baseline=-0.7ex]\node[circle,fill=popgreen,draw=popcream,line width=1.3pt,minimum size=22pt,inner sep=0]{\popcheck{white}};%
    \hspace{9pt}\begin{minipage}[c]{0.62\linewidth}
      {\popxbold\fontsize{12}{13}\selectfont\color{popcream}Signé\ \textbullet\ {\poplogo OnBuch\color{poporange}+}}\par\vspace{2pt}
      {\raggedright\popsemi\fontsize{8}{10.5}\selectfont\color{popline}\MakeUppercase{Équipe pédagogique OnBuch+}\\\MakeUppercase{Conforme au programme officiel MINESEC}\par}
    \end{minipage}\hfill
    {\poplogo\fontsize{22}{22}\selectfont\color{popcream}OnBuch\color{poporange}+}
  \end{tcolorbox}%
}

% ============================================================
% Page de garde
% #1 titre · #2 matière · #3 niveau · #4 module · #5 n° leçon · #6 durée ·
% #7 description / objectifs (texte ou itemize)
% ============================================================
\newcommand{\pagegarde}[7]{%
  \thispagestyle{empty}
  \newgeometry{a4paper,margin=1.7cm,top=1.6cm,bottom=1.4cm}%
  \begin{tikzpicture}[remember picture,overlay]
    \fill[poporange!18] ([xshift=-3.2cm,yshift=-3.6cm]current page.north east) circle (4.6cm);
    \fill[poppurple!14] ([xshift=2.4cm,yshift=7.6cm]current page.south west) circle (3.8cm);
  \end{tikzpicture}%
  {\poplogo\fontsize{30}{30}\selectfont\color{popink}OnBuch\color{poporange}+}\hfill
  \raisebox{0.6ex}{\poppill{Cours signé \textbullet\ Premium}{popink}{popcream}}\par
  \vspace{1pt}\popsquiggle{poppurple}\par
  \vspace{8pt}
  \begin{tcolorbox}[enhanced,colback=poporange,colframe=popink,boxrule=2.8pt,arc=13pt,
      drop shadow=popink,left=18pt,right=18pt,top=12pt,bottom=13pt,after skip=10pt]
    \poppill{#2 \textbullet\ #3}{popink}{popcream}\hspace{5pt}\poppill{#4}{white}{popink}\par\vspace{8pt}
    {\popdisplay\fontsize{14}{14}\selectfont\color{popink}LEÇON #5}\par\vspace{4pt}
    {\raggedright\hyphenpenalty=10000\exhyphenpenalty=10000\popdisplay\fontsize{25}{27}\selectfont\color{popcream}\MakeUppercase{#1}\par}
    \vspace{8pt}
    {\popxbold\fontsize{9.5}{9.5}\selectfont\color{popink}\MakeUppercase{Durée conseillée : #6}}
  \end{tcolorbox}
  \begin{tcolorbox}[enhanced,colback=white,colframe=popink,boxrule=2.2pt,arc=10pt,
      drop shadow=popink,left=16pt,right=16pt,top=10pt,bottom=10pt,after skip=8pt]
    \poppill{Dans cette leçon}{poppurple}{white}\par\vspace{5pt}
    \popsemi\fontsize{9.3}{12.6}\selectfont\color{popink}#7
  \end{tcolorbox}
  \noindent\begin{minipage}[t]{0.487\linewidth}
    \begin{tcolorbox}[enhanced,colback=popgreen,colframe=popink,boxrule=2.2pt,arc=10pt,
        drop shadow=popink,left=11pt,right=11pt,top=10pt,bottom=10pt,height=3.1cm,valign=center]
      \tcbox[colback=white,colframe=popink,boxrule=1.3pt,arc=5pt,boxsep=0pt,nobeforeafter,
        left=3pt,right=3pt,top=3pt,bottom=3pt,valign=center]{\qrcode[height=1.9cm]{https://play.google.com/store/apps/details?id=com.luvvix.onbuch&hl=fr}}%
      \hspace{8pt}\begin{minipage}[c]{0.5\linewidth}
        {\popdisplay\fontsize{11}{12.5}\selectfont\color{white}\MakeUppercase{Télécharge OnBuch}}\par\vspace{4pt}
        {\raggedright\popsemi\fontsize{8.4}{11}\selectfont\color{white}Gratuit \textbullet\ Google Play\\Tuteur IA, annales, concours, résultats\par}
      \end{minipage}
    \end{tcolorbox}
  \end{minipage}\hfill
  \begin{minipage}[t]{0.487\linewidth}
    \begin{tcolorbox}[enhanced,colback=poppurple,colframe=popink,boxrule=2.2pt,arc=10pt,
        drop shadow=popink,left=11pt,right=11pt,top=10pt,bottom=10pt,height=3.1cm,valign=center]
      \tikz[baseline=-0.7ex]\node[circle,fill=white,draw=popink,line width=1.3pt,minimum size=1.6cm,inner sep=0]{\popdisplay\fontsize{24}{24}\selectfont\color{poppurple}+};%
      \hspace{8pt}\begin{minipage}[c]{0.6\linewidth}
        {\popdisplay\fontsize{11}{12.5}\selectfont\color{white}\MakeUppercase{Passe à OnBuch+}}\par\vspace{4pt}
        {\raggedright\popsemi\fontsize{8.4}{11}\selectfont\color{white}Tous les cours signés, Tuteur IA illimité, fiches PDF — onglet OnBuch+ de l'app\par}
      \end{minipage}
    \end{tcolorbox}
  \end{minipage}
  \vspace{8pt}
  \popcontenu
  \vfill
  \signatureonbuch
  \restoregeometry
  \newpage
}

% Rangée « Ce cours contient » (page de garde)
\newcommand{\poptile}[3]{% #1 couleur #2 gros texte #3 légende
  \begin{tcolorbox}[enhanced,colback=#1,colframe=popink,boxrule=2pt,arc=9pt,shadow={2.5pt}{-2.5pt}{0pt}{popink},
      left=6pt,right=6pt,top=9pt,bottom=9pt,halign=center,valign=center,height=1.75cm,width=\linewidth,nobeforeafter]
    {\popdisplay\fontsize{12.5}{13}\selectfont\color{white}\MakeUppercase{#2}}\par\vspace{3pt}
    {\centering\popsemi\fontsize{7.8}{9.5}\selectfont\color{white}#3\par}
  \end{tcolorbox}}
\newcommand{\popcontenu}{%
  \par\noindent{\popxboldls\fontsize{7.5}{7.5}\selectfont\color{popmuted}CE COURS SIGNÉ CONTIENT}\par\vspace{7pt}
  \noindent\begin{minipage}[t]{0.235\linewidth}\poptile{popblue}{Cours}{complet, illustré, pas à pas}\end{minipage}\hfill
  \begin{minipage}[t]{0.235\linewidth}\poptile{poporange}{Méthodes}{et exercices résolus}\end{minipage}\hfill
  \begin{minipage}[t]{0.235\linewidth}\poptile{poppink}{Exercices}{type examen + corrigés}\end{minipage}\hfill
  \begin{minipage}[t]{0.235\linewidth}\poptile{popgreen}{Fiche bilan}{l'essentiel à retenir}\end{minipage}\par}

% Sommaire
\newtcolorbox{sommaire}{enhanced,colback=white,colframe=popink,boxrule=2.2pt,arc=10pt,
  drop shadow=popink,left=18pt,right=18pt,top=13pt,bottom=13pt,before skip=6pt,after skip=20pt,
  before upper={\poppill{Sommaire}{poppurple}{white}\par\vspace{9pt}\popsemi\fontsize{10.6}{14.5}\selectfont\color{popink}}}

% Signature de fin de cours — poussée tout en bas de la dernière page
\newcommand{\findecours}{%
  \par\vfill\nopagebreak
  \begin{center}\popsquiggle{poporange}\end{center}\vspace{4pt}
  \signatureonbuch
}
\AtBeginDocument{\def\llbracket{\mathopen{[\mkern-3.5mu[}}\def\rrbracket{\mathclose{]\mkern-3.5mu]}}} % intervalles d'entiers (glyphe ⟦ absent de la police)
% Glyphes absents de Fira Math : redéfinis à partir de glyphes disponibles
\AtBeginDocument{%
  \def\setminus{\mathbin{\backslash}}\let\smallsetminus\setminus
  \def\vdots{\mathord{\vbox{\baselineskip3.2pt\lineskiplimit0pt\kern4pt\hbox{.}\hbox{.}\hbox{.}}}}%
  \def\ddots{\mathinner{\mkern1mu\raise6.5pt\hbox{.}\mkern2mu\raise3.6pt\hbox{.}\mkern2mu\raise0.7pt\hbox{.}\mkern1mu}}%
  \def\triangle{\mathord{\scalebox{0.9}{$\symup{\Delta}$}}}%
  \def\nabla{\mathord{\raisebox{\depth}{\scalebox{1}[-1]{$\symup{\Delta}$}}}}%
  \def\checkmark{\popcheck{popdark}}%
  \def\star{\mathbin{\ast}}\def\bigstar{\mathord{\ast}}%
  \def\diamond{\mathbin{\scalebox{0.6}{$\Diamond$}}}%
  \def\bigcup{\mathop{\mathchoice{\raisebox{-0.3em}{\scalebox{1.7}{$\cup$}}}{\scalebox{1.25}{$\cup$}}{\cup}{\cup}}\displaylimits}%
  \def\bigcap{\mathop{\mathchoice{\raisebox{-0.3em}{\scalebox{1.7}{$\cap$}}}{\scalebox{1.25}{$\cap$}}{\cap}{\cap}}\displaylimits}%
  \def\lesseqqgtr{\mathrel{\lessgtr}}\def\bigoplus{\mathop{\oplus}\displaylimits}%
}
\providecommand{\ssi}{si et seulement si} % abréviation fréquente
\AtBeginDocument{\providecommand{\wideparen}{\overparen}} % arc de cercle

%% FIGFIX-BEGIN v3 — correctifs globaux des figures (docs/onbuch-generation/tools/figfix)
\usepackage{adjustbox}\usepackage{etoolbox}
% toute figure TikZ/pgfplots plus large que la ligne est réduite à la largeur disponible
\BeforeBeginEnvironment{tikzpicture}{\begin{adjustbox}{max width=\linewidth}}
\AfterEndEnvironment{tikzpicture}{\end{adjustbox}}
% légendes pgfplots lisibles par-dessus les courbes
\pgfplotsset{every axis/.append style={legend style={fill=white,fill opacity=0.92,text opacity=1,draw=black!15,inner sep=3pt}}}
% étiquettes d'arêtes (syntaxe edge["..."]) avec fond blanc
\tikzset{every edge quotes/.append style={fill=white,inner sep=1.2pt}}
% bulles de cartes mentales : texte à la ligne dans la bulle, bulle un peu plus grande
\tikzset{every concept/.append style={text width=2.0cm,align=center,minimum size=2.3cm,inner sep=2pt}}
%% FIGFIX-END
\begin{document}

\pagegarde{Production et traduction de textes liés à la vie économique}{Allemand}{1ère A}{Chapitre 2 : Vie économique}{ALA09}{2 h}{Cette leçon mixte du chapitre « Vie économique » entraîne les élèves de Première A à produire et à traduire des textes liés au monde économique : petites annonces, lettres et comptes rendus. Elle réinvestit le comparatif, les wenn-Sätze et le futur dans des tâches authentiques de thème et de version.
\vspace{4pt}
\begin{multicols}{2}\small
\begin{itemize}
\item À la fin de cette leçon, je sais identifier les caractéristiques d'une petite annonce, d'une lettre et d'un compte rendu en allemand.
\item À la fin de cette leçon, je sais utiliser le vocabulaire de la vie économique (travail, commerce, argent, entreprise) dans mes productions.
\item À la fin de cette leçon, je sais former et employer correctement le comparatif et le superlatif des adjectifs.
\item À la fin de cette leçon, je sais construire des wenn-Sätze (condition et temps) avec le verbe en fin de subordonnée.
\item À la fin de cette leçon, je sais exprimer le futur avec werden + infinitif pour parler de projets économiques.
\item À la fin de cette leçon, je sais rédiger une petite annonce, une lettre courte et un compte rendu sur un sujet économique.
\end{itemize}\end{multicols}}

\begin{sommaire}
\begin{multicols}{2}
\begin{enumerate}
\item Texte support et compréhension
\item Lexique thématique : la vie économique
\item Grammaire 1 : révision du comparatif et du superlatif
\item Grammaire 2 : wenn-Sätze et futur
\item Types de textes : petite annonce, lettre, compte rendu
\item Traduction : thème et version
\item Activité d'intégration
\item Méthodes et exercices résolus
\item Exercices
\item Corrigés des exercices
\item Fiche bilan
\end{enumerate}
\end{multicols}
\end{sommaire}

\begin{objectifs}
\begin{itemize}
\item À la fin de cette leçon, je sais identifier les caractéristiques d'une petite annonce, d'une lettre et d'un compte rendu en allemand.
\item À la fin de cette leçon, je sais utiliser le vocabulaire de la vie économique (travail, commerce, argent, entreprise) dans mes productions.
\item À la fin de cette leçon, je sais former et employer correctement le comparatif et le superlatif des adjectifs.
\item À la fin de cette leçon, je sais construire des wenn-Sätze (condition et temps) avec le verbe en fin de subordonnée.
\item À la fin de cette leçon, je sais exprimer le futur avec werden + infinitif pour parler de projets économiques.
\item À la fin de cette leçon, je sais rédiger une petite annonce, une lettre courte et un compte rendu sur un sujet économique.
\item À la fin de cette leçon, je sais traduire de courtes phrases et un texte simple du français vers l'allemand (thème).
\item À la fin de cette leçon, je sais traduire un texte économique court de l'allemand vers le français (version).
\item À la fin de cette leçon, je sais relire et corriger ma production (orthographe, majuscules, ordre des mots, cas).
\end{itemize}
\end{objectifs}

\begin{prerequis}
\begin{itemize}
\item La conjugaison des verbes au présent et au Präteritum (Seconde)
\item L'ordre des mots dans la phrase principale et la subordonnée (Seconde)
\item Les déclinaisons des articles et des adjectifs (Seconde)
\item Le vocabulaire de base des métiers et des achats (Seconde)
\item Les connecteurs temporels et causaux simples : weil, denn, als, wenn (début de Première)
\end{itemize}
\end{prerequis}

\newpage

\coursec{Texte support et compréhension}

\courssub{Situation d'accroche et problématique}

Aminatou, élève de Première A au lycée de Yaoundé, a un rêve : faire un stage dans une entreprise allemande implantée à Douala. Pour postuler, elle doit rédiger une lettre de motivation en allemand, répondre à une petite annonce publiée dans un journal et traduire un court compte rendu d'activité économique. En Allemagne, en Autriche et en Suisse, la maîtrise des textes économiques — petites annonces, lettres commerciales, compte rendus — est indispensable dans la vie professionnelle quotidienne. Cette leçon te donne tous les outils pour produire et traduire ces textes comme un vrai professionnel !

\textbf{Problématique :} Comment lire, comprendre et exploiter un texte sur la vie économique pour ensuite produire et traduire des textes professionnels en allemand ?

\begin{methode}[Lire un texte économique en 3 étapes]
\begin{enumerate}
\item \textbf{Survoler} le texte (titre, première et dernière phrase) pour identifier le thème général : de quoi parle-t-on ?
\item \textbf{Lire en détail} avec le glossaire : souligner les mots inconnus, les chercher dans le glossaire, puis relire la phrase entière.
\item \textbf{Repérer les chiffres et les mots-clés} : noms de pays, monnaies, secteurs d'activité, verbes importants. Ce sont eux qui portent l'information économique.
\end{enumerate}
\end{methode}

\courssub{Lecture du texte support}

\begin{textebox}[Die Wirtschaft in Kamerun und Deutschland]
In Kamerun arbeiten viele Menschen in der Landwirtschaft. Sie produzieren Kakao, Kaffee und Bananen. In den Städten wie Douala und Yaoundé gibt es auch viele kleine Geschäfte und Märkte. In Deutschland gibt es mehr Industrie als in Kamerun. Viele Menschen arbeiten in großen Betrieben, zum Beispiel in der Autoindustrie. Die Arbeitslosigkeit ist in Deutschland niedriger als in vielen anderen Ländern Europas.

Wenn die Wirtschaft wächst, gibt es mehr Arbeitsplätze. Wenn junge Menschen eine gute Ausbildung haben, finden sie schneller einen Job. In Deutschland lernen viele Jugendliche einen Beruf in der dualen Ausbildung : sie arbeiten in einem Betrieb und gehen gleichzeitig zur Berufsschule.

Viele junge Kameruner werden in Zukunft in modernen Betrieben arbeiten. Sie werden mehr Geld verdienen, wenn sie gut ausgebildet sind. Der Euro ist die Währung in Deutschland ; er ist stärker als der CFA-Franc, aber das Leben in Deutschland ist auch teurer. Ein Brot kostet in München mehr als in Yaoundé, aber die Gehälter sind auch höher.

Der Handel zwischen Kamerun und Deutschland ist wichtig für beide Länder. Kamerun exportiert Kakao und Holz nach Europa, und Deutschland verkauft Maschinen und Autos nach Afrika. Wenn die beiden Länder zusammenarbeiten, profitieren beide Seiten.
\end{textebox}

\begin{definition}[Die Wirtschaft]
\all{Die Wirtschaft} (féminin, sans pluriel usuel) désigne l'ensemble des activités de production, d'échange et de consommation d'un pays : c'est \textbf{l'économie}. On dit : \all{die deutsche Wirtschaft} « l'économie allemande », \all{die Wirtschaft wächst} « l'économie croît / est en croissance ».
\end{definition}

\courssub{Glossaire bilingue des mots difficiles}

\begin{center}
\begin{tabularx}{\linewidth}{|X|X|}
\hline
\rowcolor{popblueL} \textbf{Deutsch} & \textbf{Français} \\
\hline
die Wirtschaft, -en & l'économie ; l'entreprise (sens 2) \\
\hline
die Landwirtschaft, -en & l'agriculture \\
\hline
der Betrieb, -e & l'entreprise, l'exploitation \\
\hline
die Arbeitslosigkeit (Sg.) & le chômage \\
\hline
der Arbeitsplatz, "-e & l'emploi, le poste de travail \\
\hline
verdienen (verdient, verdiente, hat verdient) & gagner (de l'argent) \\
\hline
die Währung, -en & la monnaie \\
\hline
das Gehalt, "-er & le salaire \\
\hline
die Ausbildung, -en & la formation (professionnelle) \\
\hline
die Berufsschule, -n & l'école professionnelle, le CFA \\
\hline
der Handel (Sg.) & le commerce, les échanges \\
\hline
exportieren / verkaufen & exporter / vendre \\
\hline
wachsen (wächst, wuchs, ist gewachsen) & croître, augmenter \\
\hline
niedrig / hoch — höher & bas / haut — plus haut \\
\hline
teuer — teurer & cher — plus cher \\
\hline
profitieren (von + Dat.) & profiter (de) \\
\hline
\end{tabularx}
\end{center}

\begin{attention}[Faux amis et pièges de vocabulaire]
\begin{itemize}
\item \all{der Betrieb} ne veut pas dire « le bétria » ni « le bâtiment » : c'est \textbf{l'entreprise}. « Le bâtiment » se dit \all{das Gebäude}.
\item \all{die Währung} = la monnaie (l'euro, le franc CFA) ; « la monnaie » au sens de « la petite monnaie, les pièces » se dit \all{das Kleingeld}.
\item \all{Geld verdienen} = gagner de l'argent ; ne dis jamais \all{Geld gewinnen} dans ce sens : \all{gewinnen} signifie « gagner » un match, un prix, à la loterie.
\item N'oublie pas la majuscule à tous les noms : \all{die Wirtschaft}, \all{der Handel}, \all{das Gehalt}.
\end{itemize}
\end{attention}

\courssub{Carte mentale du champ lexical}

\begin{popfigure}
\begin{center}
\begin{tikzpicture}[mindmap, grow cyclic, every node/.style={concept}, concept color=popblue!40, text=popdark, root concept/.append style={font=\Large\bfseries}, level 1/.append style={level distance=3.2cm, sibling angle=90}]
\node[root concept]{Wirtschaft}
  child[concept color=popgreen!40] { node {Arbeit} }
  child[concept color=popgold!50] { node {Geld} }
  child[concept color=poporange!40] { node {Handel} }
  child[concept color=poppurple!35] { node {Unternehmen} };
\end{tikzpicture}
\end{center}
\legende{Carte mentale : les quatre grandes branches du champ lexical de \all{die Wirtschaft} : le travail, l'argent, le commerce et les entreprises.}
\end{popfigure}

Complète mentalement chaque branche avec les mots du glossaire : \all{der Arbeitsplatz} va avec \all{Arbeit}, \all{die Währung} avec \all{Geld}, \all{exportieren} avec \all{Handel}, \all{der Betrieb} avec \all{Unternehmen}.

\courssub{Compréhension globale : les W-Fragen}

Pour vérifier la compréhension globale d'un texte, on pose les questions en \all{W-} : \all{wer ?} (qui ?), \all{was ?} (quoi ?), \all{wo ?} (où ?), \all{wann ?} (quand ?), \all{warum ?} (pourquoi ?).

\begin{center}
\begin{tabularx}{\linewidth}{|X|X|X|}
\hline
\rowcolor{popblueL} \textbf{W-Frage} & \textbf{Français} & \textbf{Réponse dans le texte} \\
\hline
Wer ? & Qui ? & viele Menschen, junge Kameruner, Jugendliche \\
\hline
Was ? & Quoi ? & die Wirtschaft in Kamerun und Deutschland \\
\hline
Wo ? & Où ? & in Kamerun (Douala, Yaoundé) und in Deutschland \\
\hline
Wann ? & Quand ? & heute und in Zukunft \\
\hline
Warum ? & Pourquoi ? & weil der Handel für beide Länder wichtig ist \\
\hline
\end{tabularx}
\end{center}

\begin{exoresolu}[Questions de compréhension globale]
Réponds en allemand par une phrase complète.
\begin{enumerate}
\item Worum geht es in diesem Text ?
\item Wo arbeiten viele Menschen in Kamerun ?
\end{enumerate}
\tcblower
\textbf{Solution détaillée :}
\begin{enumerate}
\item \all{Der Text handelt von der Wirtschaft in Kamerun und Deutschland.} « Le texte traite de l'économie au Cameroun et en Allemagne. » On utilise le verbe \all{handeln von} + datif : « traiter de ». Attention à la place du verbe en deuxième position.
\item \all{In Kamerun arbeiten viele Menschen in der Landwirtschaft.} « Au Cameroun, beaucoup de gens travaillent dans l'agriculture. » Le complément de lieu \all{in Kamerun} est en tête, donc le verbe \all{arbeiten} vient juste après, puis le sujet \all{viele Menschen} : c'est l'inversion obligatoire en allemand.
\end{enumerate}
\end{exoresolu}

\courssub{Compréhension détaillée}

\begin{exercice}[Richtig oder falsch ?]
Lis le texte et dis si les phrases sont vraies (\all{richtig}) ou fausses (\all{falsch}). Corrige les phrases fausses.
\begin{enumerate}
\item In Kamerun gibt es mehr Industrie als in Deutschland.
\item Die Arbeitslosigkeit ist in Deutschland niedriger als in vielen anderen Ländern Europas.
\item Der CFA-Franc ist stärker als der Euro.
\item Kamerun exportiert Kakao und Holz nach Europa.
\item Junge Kameruner werden in Zukunft weniger Geld verdienen.
\end{enumerate}
\end{exercice}

\begin{corrige}[Exercice « Richtig oder falsch ? »]
\begin{enumerate}
\item \textbf{Falsch.} \all{In Deutschland gibt es mehr Industrie als in Kamerun.} C'est l'inverse : l'Allemagne est plus industrialisée.
\item \textbf{Richtig.} Le texte dit : \all{Die Arbeitslosigkeit ist in Deutschland niedriger als in vielen anderen Ländern Europas.}
\item \textbf{Falsch.} \all{Der Euro ist stärker als der CFA-Franc.} Attention au sens du comparatif \all{stärker als} : « plus fort que ».
\item \textbf{Richtig.} Le texte dit : \all{Kamerun exportiert Kakao und Holz nach Europa.}
\item \textbf{Falsch.} \all{Sie werden mehr Geld verdienen, wenn sie gut ausgebildet sind.} Ils gagneront \emph{plus} d'argent, pas moins : \all{mehr} = « plus », \all{weniger} = « moins ».
\end{enumerate}
\end{corrige}

\begin{exercice}[Questions de compréhension détaillée]
Réponds par des phrases complètes en allemand.
\begin{enumerate}
\item Was produzieren die Menschen in Kamerun ?
\item Was lernen viele Jugendliche in Deutschland in der dualen Ausbildung ?
\item Was verkauft Deutschland nach Afrika ?
\item Warum profitieren beide Seiten ?
\end{enumerate}
\end{exercice}

\begin{corrige}[Exercice « Questions de compréhension détaillée »]
\begin{enumerate}
\item \all{Sie produzieren Kakao, Kaffee und Bananen.} « Ils produisent du cacao, du café et des bananes. »
\item \all{Sie lernen einen Beruf : sie arbeiten in einem Betrieb und gehen gleichzeitig zur Berufsschule.} « Ils apprennent un métier : ils travaillent dans une entreprise et vont en même temps à l'école professionnelle. »
\item \all{Deutschland verkauft Maschinen und Autos nach Afrika.} « L'Allemagne vend des machines et des voitures vers l'Afrique. »
\item \all{Beide Seiten profitieren, wenn die beiden Länder zusammenarbeiten.} « Les deux parties profitent quand les deux pays coopèrent. »
\end{enumerate}
\end{corrige}

\courssub{Observation guidée : la grammaire du texte}

Relisons le texte avec des yeux de grammairien. Trois structures reviennent souvent dans les textes économiques ; elles seront révisées en détail dans les sections 3 et 4.

\textbf{1. Le comparatif.} Cherche dans le texte les formes en \all{-er als} :

\all{In Deutschland gibt es \textbf{mehr} Industrie \textbf{als} in Kamerun.} « Il y a plus d'industrie en Allemagne qu'au Cameroun. »

\all{Der Euro ist \textbf{stärker als} der CFA-Franc.} « L'euro est plus fort que le franc CFA. »

\all{Das Leben in Deutschland ist \textbf{teurer}.} « La vie en Allemagne est plus chère. »

\all{Die Gehälter sind \textbf{höher}.} « Les salaires sont plus élevés. »

On observe : adjectif + \all{-er} + \all{als} = adjectif + « -er/plus » + « que » en français. Note l'Umlaut sur les adjectifs courts : \all{stark} devient \all{stärker}, comme \all{alt} → \all{älter}.

\textbf{2. Les wenn-Sätze.} Cherche les phrases qui commencent par \all{wenn} :

\all{\textbf{Wenn} die Wirtschaft \textbf{wächst}, gibt es mehr Arbeitsplätze.} « Quand / si l'économie croît, il y a plus d'emplois. »

\all{\textbf{Wenn} die beiden Länder \textbf{zusammenarbeiten}, profitieren beide Seiten.} « Si les deux pays coopèrent, les deux parties en profitent. »

On observe : dans la subordonnée introduite par \all{wenn}, le verbe conjugué va \textbf{à la fin} ; quand la subordonnée est en tête, le verbe de la principale vient juste après la virgule (\all{gibt es}, \all{profitieren beide Seiten}).

\textbf{3. Le futur.} Cherche les phrases avec \all{werden} + infinitif :

\all{Viele junge Kameruner \textbf{werden} in Zukunft in modernen Betrieben \textbf{arbeiten}.} « Beaucoup de jeunes Camerounais travailleront à l'avenir dans des entreprises modernes. »

\all{Sie \textbf{werden} mehr Geld \textbf{verdienen}.} « Ils gagneront plus d'argent. »

On observe : \all{werden} conjugué en deuxième position + infinitif \textbf{à la fin} de la phrase = le futur français.

\begin{attention}[Pièges des francophones]
\begin{itemize}
\item Ne traduis jamais « que » de comparaison par \all{wie} : on dit \all{stärker \textbf{als}}, jamais \all{stärker wie}. \all{Wie} sert pour l'égalité : \all{so stark wie} « aussi fort que ».
\item \all{Wenn} n'est pas \all{wann} : \all{wenn} = « si / quand » (subordonnée), \all{wann} = « quand ? » (question directe ou indirecte).
\item Au futur, l'infinitif se place \textbf{à la toute fin} : \all{Sie werden mehr Geld \textbf{verdienen}}, et non \all{Sie werden verdienen mehr Geld}.
\end{itemize}
\end{attention}

\courssub{Tableau comparatif : Kamerun — Deutschland}

\begin{popfigure}
\begin{center}
\begin{tabularx}{\linewidth}{|X|X|X|}
\hline
\rowcolor{popgreenL} \textbf{Aspekt} & \textbf{Kamerun} & \textbf{Deutschland} \\
\hline
Währung & der CFA-Franc & der Euro \\
\hline
Wichtige Sektoren & Landwirtschaft (Kakao, Kaffee, Bananen), kleine Geschäfte & Industrie (Autos, Maschinen), große Betriebe \\
\hline
Sprachen bei der Arbeit & Französisch, Englisch, Sprachen Kameruns & Deutsch \\
\hline
Ausbildung & Schulen und Universitäten & die duale Ausbildung (Betrieb + Berufsschule) \\
\hline
Exporte & Kakao und Holz nach Europa & Maschinen und Autos nach Afrika \\
\hline
\end{tabularx}
\end{center}
\legende{Tableau comparatif : quelques aspects de la vie économique au Cameroun et en Allemagne.}
\end{popfigure}

\courssub{Discussion : similitudes et différences}

\begin{experience}[Débat en classe : Kamerun und Deutschland im Vergleich]
Par groupes de quatre, discutez en allemand à partir du texte et du tableau. Chaque groupe prépare deux similitudes (\all{Gemeinsamkeiten}) et deux différences (\all{Unterschiede}), puis les présente à la classe.

Phrases d'aide pour s'exprimer :
\begin{itemize}
\item \all{In Kamerun gibt es ..., aber in Deutschland ...} « Au Cameroun il y a ..., mais en Allemagne ... »
\item \all{In beiden Ländern ...} « Dans les deux pays ... »
\item \all{Ich denke, dass ...} « Je pense que ... » (verbe à la fin !)
\item \all{Meiner Meinung nach ...} « À mon avis ... »
\end{itemize}
Exemple de production attendue : \all{In beiden Ländern ist die Ausbildung wichtig. In Deutschland ist das Leben teurer als in Kamerun, aber die Gehälter sind höher.} « Dans les deux pays, la formation est importante. En Allemagne, la vie est plus chère qu'au Cameroun, mais les salaires sont plus élevés. »
\end{experience}

\begin{savaistu}[L'Allemagne, première économie d'Europe]
L'Allemagne est la première économie d'Europe et l'un des plus grands exportateurs du monde. Son système de formation professionnelle en alternance, \all{die duale Ausbildung}, est un modèle mondial : les jeunes partagent leur semaine entre l'entreprise (\all{der Betrieb}) et l'école professionnelle (\all{die Berufsschule}), et ils reçoivent déjà un petit salaire pendant leur formation. Ce modèle inspire aujourd'hui de nombreux pays, y compris en Afrique, pour lutter contre le chômage des jeunes.
\end{savaistu}

\begin{aretenir}[Ce qu'il faut retenir de cette section]
\begin{itemize}
\item Un texte économique se lit en trois étapes : survol, lecture détaillée avec glossaire, repérage des mots-clés et des chiffres.
\item Les \all{W-Fragen} (\all{wer, was, wo, wann, warum}) permettent de vérifier la compréhension globale.
\item Trois structures grammaticales dominent ces textes : le comparatif (\all{stärker als}), les subordonnées en \all{wenn} (verbe à la fin) et le futur (\all{werden} + infinitif final).
\item Le vocabulaire clé : \all{die Wirtschaft}, \all{der Betrieb}, \all{die Arbeitslosigkeit}, \all{verdienen}, \all{die Währung}, \all{der Handel}.
\end{itemize}
\end{aretenir}

\coursec{Lexique thématique : la vie économique}

Dans la section précédente, nous avons lu et commenté un texte sur la vie économique. Vous avez remarqué que ce texte contenait beaucoup de mots nouveaux : \all{der Betrieb}, \all{der Lohn}, \all{der Kunde}... Pour bien comprendre, parler et traduire des textes économiques, il faut d'abord maîtriser ce vocabulaire. C'est l'objectif de cette section : construire pas à pas un lexique solide, organisé en quatre grands champs lexicaux : le \cle{travail}, l'\cle{argent}, le \cle{commerce} et l'\cle{entreprise}.

\begin{methode}[Comment apprendre un mot allemand ?]
\begin{enumerate}
\item Apprenez toujours un nom avec son \textbf{article} (\all{der}, \all{die}, \all{das}) et son \textbf{pluriel} : \all{der Markt, die Märkte}.
\item Apprenez les verbes avec leur \textbf{construction} (cas et préposition) : \all{sich bewerben um + Akkusativ}.
\item Classez les mots par \textbf{thème} (travail, argent, commerce...) dans votre cahier de vocabulaire.
\item Faites des phrases complètes avec chaque mot nouveau, par exemple sur le contexte camerounais : \all{Der Markt in Mokolo ist sehr groß.} « Le marché de Mokolo est très grand. »
\end{enumerate}
\end{methode}

\courssub{Le champ lexical du travail}

\begin{definition}[Les mots du travail]
\begin{itemize}
\item \all{der Beruf, -e} : le métier, la profession \quad \all{Was sind Sie von Beruf?} « Quel est votre métier ? »
\item \all{die Arbeit, -en} : le travail \quad \all{Die Arbeit beginnt um 8 Uhr.} « Le travail commence à 8 heures. »
\item \all{der Arbeitsplatz, die Arbeitsplätze} : le poste de travail, l'emploi
\item \all{der Arbeitgeber, -} : l'employeur \quad \all{die Arbeitnehmerin, -nen} : l'employée (féminin ; masculin : \all{der Arbeitnehmer, -})
\item \all{die Bewerbung, -en} : la candidature
\item \all{das Vorstellungsgespräch, -e} : l'entretien d'embauche
\item \all{die Stelle, -n} : le poste, l'emploi (offre d'emploi)
\end{itemize}
\end{definition}

Observez la famille de mots autour de \all{arbeiten} : \all{die Arbeit} (le travail), \all{der Arbeiter} (l'ouvrier), \all{der Arbeitgeber} (celui qui \emph{donne} le travail), \all{der Arbeitnehmer} (celui qui \emph{prend} le travail). L'allemand forme des mots composés très logiques : \all{der Geber} vient de \all{geben} (donner) et \all{der Nehmer} de \all{nehmen} (prendre). Retenez cette logique, elle vous aidera à deviner le sens de nombreux mots !

\begin{exemplebox}[Exemples traduits]
\all{Sie arbeitet bei einer Bank.} « Elle travaille dans une banque. »

\all{Er bewirbt sich um eine Stelle.} « Il postule pour un emploi. »

\all{Nach der Bewerbung hat sie ein Vorstellungsgespräch.} « Après sa candidature, elle a un entretien d'embauche. »

\all{Mein Onkel ist Arbeitgeber in einer Fabrik in Douala.} « Mon oncle est employeur dans une usine à Douala. »
\end{exemplebox}

\courssub{Le champ lexical de l'argent}

\begin{definition}[Les mots de l'argent]
\begin{itemize}
\item \all{das Geld, -er} : l'argent \quad \all{der Preis, -e} : le prix
\item \all{der Lohn, die Löhne} : le salaire (horaire, des ouvriers)
\item \all{das Gehalt, die Gehälter} : le salaire (mensuel, des employés et cadres)
\item \all{die Rechnung, -en} : la facture, l'addition (au restaurant)
\item \all{die Währung, -en} : la monnaie (d'un pays)
\item \all{das Konto, die Konten} : le compte (en banque)
\item \all{sparen} : économiser, épargner \quad \all{ausgeben (gibt aus, gab aus, hat ausgegeben)} : dépenser
\end{itemize}
\end{definition}

\begin{definition}[der Lohn ou das Gehalt ?]
Les deux mots signifient « salaire », mais attention à la différence :
\begin{itemize}
\item \all{der Lohn} : salaire payé à l'\textbf{heure} ou à la semaine, traditionnellement pour les ouvriers et les travaux manuels. \all{Der Bauarbeiter bekommt seinen Lohn jede Woche.} « L'ouvrier du bâtiment reçoit son salaire chaque semaine. »
\item \all{das Gehalt} : salaire \textbf{mensuel} fixe, pour les employés de bureau et les cadres. \all{Die Lehrerin bekommt ihr Gehalt am Ende des Monats.} « La professeure reçoit son salaire à la fin du mois. »
\end{itemize}
\end{definition}

\begin{exemplebox}[Exemples traduits]
\all{Das Auto kostet 15 000 Euro.} « La voiture coûte 15 000 euros. »

\all{Ich spare Geld für ein neues Handy.} « J'économise de l'argent pour un nouveau portable. »

\all{Sie gibt viel Geld für Kleider aus.} « Elle dépense beaucoup d'argent pour les vêtements. »

\all{Die Rechnung, bitte!} « L'addition, s'il vous plaît ! »

\all{Er hat ein Konto bei einer Bank in Yaoundé.} « Il a un compte dans une banque à Yaoundé. »
\end{exemplebox}

\begin{attention}[ausgeben : verbe à particule séparable]
\all{Ausgeben} est un verbe à \textbf{particule séparable} : à la forme conjuguée, la particule \all{aus} va à la fin de la phrase : \all{Sie \textbf{gibt} viel Geld \textbf{aus}.} Au Perfekt, le \all{ge-} s'insère entre la particule et le radical : \all{Sie hat viel Geld \textbf{ausgegeben}.} « Elle a beaucoup dépensé d'argent. »
\end{attention}

\begin{savaistu}[Les monnaies des pays germanophones]
L'Allemagne et l'Autriche utilisent \all{der Euro} (l'euro). Mais la Suisse n'est pas membre de l'Union européenne : on y paie avec \all{der Schweizer Franken} (le franc suisse). Au Cameroun, la monnaie est le franc CFA : \all{der CFA-Franken}. Le mot allemand \all{die Währung} désigne la monnaie d'un pays : \all{Der Euro ist die Währung in Deutschland.} « L'euro est la monnaie en Allemagne. »
\end{savaistu}

\courssub{Le champ lexical du commerce}

\begin{definition}[Les mots du commerce]
\begin{itemize}
\item \all{der Handel} : le commerce (l'activité) \quad \all{das Geschäft, -e} : le magasin, le commerce, l'affaire
\item \all{der Markt, die Märkte} : le marché \quad \all{der Supermarkt, die Supermärkte} : le supermarché
\item \all{der Kunde, -n / die Kundin, -nen} : le client / la cliente
\item \all{kaufen} : acheter \quad \all{verkaufen} : vendre
\item \all{der Verkäufer, - / die Verkäuferin, -nen} : le vendeur / la vendeuse
\item \all{der Export, -e} : l'exportation \quad \all{der Import, -e} : l'importation
\item \all{exportieren / importieren} : exporter / importer
\end{itemize}
\end{definition}

\begin{exemplebox}[Exemples traduits]
\all{Der Verkäufer verkauft Obst auf dem Markt.} « Le vendeur vend des fruits au marché. »

\all{Die Kundin kauft Gemüse im Geschäft.} « La cliente achète des légumes dans le magasin. »

\all{Kamerun exportiert Kakao und importiert Autos.} « Le Cameroun exporte du cacao et importe des voitures. »

\all{Der Handel zwischen Deutschland und Kamerun ist wichtig.} « Le commerce entre l'Allemagne et le Cameroun est important. »
\end{exemplebox}

\begin{attention}[Attention : der Kunde est un nom faible !]
\all{Der Kunde} appartient à la déclinaison faible (n-Deklination) : il prend un \textbf{-n} à tous les cas sauf au nominatif singulier : \all{Ich helfe dem Kunden.} (datif) « J'aide le client. » ; \all{Der Verkäufer fragt den Kunden.} (accusatif) « Le vendeur interroge le client. » Ne dites donc jamais \all{dem Kunde} ! Autres noms faibles déjà connus : \all{der Junge, -n}, \all{der Student, -en}, \all{der Tourist, -en}.
\end{attention}

\courssub{Le champ lexical de l'entreprise}

\begin{definition}[Les mots de l'entreprise]
\begin{itemize}
\item \all{die Firma, die Firmen} : la firme, l'entreprise (société commerciale)
\item \all{der Betrieb, -e} : l'entreprise, l'exploitation (le lieu de production)
\item \all{das Unternehmen, -} : l'entreprise (terme général et économique)
\item \all{der Chef, -s / die Chefin, -nen} : le chef, la cheffe, le patron
\item \all{der Mitarbeiter, - / die Mitarbeiterin, -nen} : l'employé, le collaborateur / l'employée
\item \all{die Produktion, -en} : la production
\item \all{gründen} : fonder, créer (une entreprise) \quad \all{produzieren} : produire
\end{itemize}
\end{definition}

\begin{exemplebox}[Exemples traduits]
\all{Die Firma hat 50 Mitarbeiter.} « L'entreprise compte 50 employés. »

\all{Sein Vater hat ein kleines Unternehmen gegründet.} « Son père a créé une petite entreprise. »

\all{Der Chef ist mit der Produktion zufrieden.} « Le chef est satisfait de la production. »

\all{In diesem Betrieb arbeiten viele Jugendliche.} « Beaucoup de jeunes travaillent dans cette entreprise. »
\end{exemplebox}

\courssub{Adjectifs utiles et verbes à construction particulière}

\begin{definition}[Adjectifs du domaine économique]
\begin{itemize}
\item \all{billig} : bon marché $\leftrightarrow$ \all{teuer} : cher
\item \all{reich} : riche $\leftrightarrow$ \all{arm} : pauvre
\item \all{erfolgreich} : qui réussit, prospère \quad \all{modern} : moderne
\item \all{wirtschaftlich} : économique (relatif à l'économie)
\end{itemize}
\all{Der Markt ist billig, das Geschäft ist teuer.} « Le marché est bon marché, le magasin est cher. »
\end{definition}

\begin{propriete}[Verbes à construction particulière : mémorisez verbe + cas + préposition]
\begin{itemize}
\item \all{verdienen + Akkusativ} : gagner (de l'argent). \all{Er verdient 100 000 Franken im Monat.} « Il gagne 100 000 francs par mois. »
\item \all{kosten + Akkusativ} : coûter. \all{Das Brot kostet 200 Franken.} « Le pain coûte 200 francs. »
\item \all{arbeiten bei + Dativ} : travailler dans (une entreprise). \all{Sie arbeitet bei einer Bank.} « Elle travaille dans une banque. »
\item \all{sich bewerben um + Akkusativ} : postuler pour, candidater à. Verbe réfléchi et fort : \all{ich bewerbe mich, er bewarb sich, er hat sich beworben}. \all{Wir bewerben uns um einen Arbeitsplatz.} « Nous postulons pour un emploi. »
\end{itemize}
\end{propriete}

\begin{attention}[Faux amis et pièges de traduction]
\begin{itemize}
\item « actuel » se traduit par \all{aktuell} : \all{die aktuelle Wirtschaftslage} « la situation économique actuelle ».
\item « une fabrique » = \all{die Fabrik}, mais « fabriquer » = \all{herstellen} ou \all{produzieren}, jamais un verbe « fabrizieren » : \all{Die Fabrik produziert Schokolade.} « L'usine fabrique du chocolat. »
\item « dépenser » (de l'argent) = \all{Geld ausgeben} ; ne confondez pas avec \all{sparen} (économiser).
\item « gagner de l'argent » = \all{Geld verdienen} ; \all{gewinnen} signifie « gagner » un match, un prix : \all{Kamerun gewinnt das Spiel.} « Le Cameroun gagne le match. »
\item N'oubliez pas la majuscule à tous les noms : \all{das Geld}, \all{der Preis}, \all{die Rechnung}.
\end{itemize}
\end{attention}

\courssub{Tableau récapitulatif du lexique}

\begin{center}
\begin{tabularx}{\linewidth}{|X|X|X|X|}
\hline
\rowcolor{popblueL} \textbf{Arbeit (travail)} & \textbf{Geld (argent)} & \textbf{Handel (commerce)} & \textbf{Unternehmen (entreprise)} \\
\hline
der Beruf, -e : le métier & das Geld : l'argent & der Handel : le commerce & die Firma, Firmen : la firme \\
\hline
die Arbeit, -en : le travail & der Preis, -e : le prix & das Geschäft, -e : le magasin & der Betrieb, -e : l'entreprise \\
\hline
der Arbeitsplatz : l'emploi & der Lohn / das Gehalt : le salaire & der Markt, Märkte : le marché & das Unternehmen, - : l'entreprise \\
\hline
der Arbeitgeber : l'employeur & die Rechnung, -en : la facture & der Kunde, -n : le client & der Chef, -s : le patron \\
\hline
die Bewerbung, -en : la candidature & die Währung, -en : la monnaie & kaufen / verkaufen : acheter / vendre & der Mitarbeiter, - : l'employé \\
\hline
das Vorstellungsgespräch : l'entretien & sparen / ausgeben : épargner / dépenser & der Export / der Import & die Produktion : la production \\
\hline
sich bewerben um + Akk. & das Konto, Konten : le compte & der Verkäufer, - : le vendeur & gründen : fonder \\
\hline
\end{tabularx}
\end{center}

\courssub{Le circuit économique : comprendre les relations entre les mots}

Le vocabulaire économique n'est pas une liste isolée : les mots sont liés entre eux dans un circuit. L'entreprise produit, le client achète, l'argent circule et revient à l'entreprise, qui paie ses employés. Observez le schéma suivant :

\begin{popfigure}
\begin{center}
\begin{tikzpicture}[node distance=2.6cm, >=Stealth]
\node[draw=popblue, fill=popblueL, rounded corners, thick, align=center, minimum width=3.2cm, minimum height=1.2cm] (unt) {\textbf{das Unternehmen}\\ l'entreprise};
\node[draw=popgreen, fill=popgreenL, rounded corners, thick, align=center, minimum width=3.2cm, minimum height=1.2cm, right=4.5cm of unt] (kunde) {\textbf{der Kunde}\\ le client};
\node[draw=poporange, fill=poporangeL, rounded corners, thick, align=center, minimum width=3.2cm, minimum height=1.2cm, below=2.6cm of kunde] (geld) {\textbf{das Geld}\\ l'argent};
\node[draw=poppurple, fill=poppurpleL, rounded corners, thick, align=center, minimum width=3.2cm, minimum height=1.2cm, left=4.5cm of geld] (mit) {\textbf{der Mitarbeiter}\\ l'employé};
\draw[->, very thick, popblue] (unt) -- node[above, align=center]{produziert und\\ verkauft} (kunde);
\draw[->, very thick, popgreen] (kunde) -- node[right, align=center]{kauft und\\ bezahlt} (geld);
\draw[->, very thick, poporange] (geld) -- node[below, align=center]{der Lohn /\\ das Gehalt} (mit);
\draw[->, very thick, poppurple] (mit) -- node[left, align=center]{arbeitet\\ für} (unt);
\end{tikzpicture}
\end{center}
\legende{Le circuit économique simple : l'entreprise produit, le client achète, l'argent paie les employés qui travaillent pour l'entreprise.}
\end{popfigure}

\courssub{Texte de réinvestissement : une journée au marché}

Lisez ce petit texte qui réemploie le vocabulaire de la section, puis répondez aux questions.

\begin{textebox}[Ein Tag auf dem Markt in Yaoundé]
Frau Mbarga arbeitet bei einer kleinen Firma in Yaoundé. Die Firma produziert Saft und verkauft ihn auf dem Markt. Frau Mbarga ist Verkäuferin. Sie beginnt ihre Arbeit um 7 Uhr morgens. Viele Kunden kommen: Sie kaufen Orangensaft und Mangosaft. Der Saft ist nicht teuer, eine Flasche kostet nur 500 Franken. „Unser Saft ist billig und modern produziert“, sagt Frau Mbarga. „Die Kunden sind sehr zufrieden.“

Am Ende des Monats bekommt Frau Mbarga ihr Gehalt. Sie verdient nicht viel, aber sie spart jeden Monat ein bisschen Geld auf ihrem Konto. Ihr Traum: Sie will später ein eigenes Unternehmen gründen. „Ich möchte mein eigener Chef sein“, sagt sie. Ihre Tochter Amina ist 18 Jahre alt. Sie bewirbt sich jetzt um eine Stelle bei einer Bank. Nächste Woche hat sie ein Vorstellungsgespräch. Die ganze Familie hofft: Amina wird erfolgreich sein!

Kamerun exportiert Kakao, Kaffee und Bananen nach Europa. Der Handel mit Deutschland ist wichtig für die wirtschaftliche Entwicklung des Landes.
\end{textebox}

\textbf{Glossaire :} \all{der Saft, die Säfte} : le jus ; \all{die Flasche, -n} : la bouteille ; \all{zufrieden} : satisfait ; \all{der Traum, die Träume} : le rêve ; \all{eigen} : propre ; \all{die Entwicklung, -en} : le développement ; \all{das Land, die Länder} : le pays.

\textbf{Questions de compréhension :}
\begin{enumerate}
\item \all{Wo arbeitet Frau Mbarga?} (Où travaille Mme Mbarga ?)
\item \all{Was kostet eine Flasche Saft?} (Combien coûte une bouteille de jus ?)
\item \all{Was ist Frau Mbargas Traum?} (Quel est le rêve de Mme Mbarga ?)
\item \all{Was exportiert Kamerun?} (Qu'est-ce que le Cameroun exporte ?)
\end{enumerate}

\courssub{Exercice résolu}

\begin{exoresolu}[Complétez avec le mot qui convient]
Complétez les phrases avec : \all{Gehalt}, \all{bewirbt}, \all{kostet}, \all{Kunden}, \all{gegründet}, \all{bei}.
\begin{enumerate}
\item Mein Bruder \ldots\ sich um eine Stelle in Douala.
\item Das Brot \ldots\ 250 Franken.
\item Meine Tante arbeitet \ldots\ einer Bank.
\item Der Verkäufer hilft den \ldots\ .
\item Am Ende des Monats bekommt sie ihr \ldots\ .
\item Er hat eine kleine Firma \ldots\ .
\end{enumerate}
\tcblower
\textbf{Solution détaillée :}
\begin{enumerate}
\item \all{Mein Bruder \textbf{bewirbt} sich um eine Stelle in Douala.} « Mon frère postule pour un emploi à Douala. » — construction : \all{sich bewerben um + Akk.}, ici 3\up{e} personne du singulier.
\item \all{Das Brot \textbf{kostet} 250 Franken.} « Le pain coûte 250 francs. » — \all{kosten + Akk.}, sujet au singulier : \all{kostet}.
\item \all{Meine Tante arbeitet \textbf{bei} einer Bank.} « Ma tante travaille dans une banque. » — \all{arbeiten bei + Dativ} : \all{bei einer Bank}.
\item \all{Der Verkäufer hilft den \textbf{Kunden}.} « Le vendeur aide les clients. » — au datif pluriel, \all{Kunde} (nom faible) prend -n : \all{den Kunden}.
\item \all{Am Ende des Monats bekommt sie ihr \textbf{Gehalt}.} « À la fin du mois, elle reçoit son salaire. » — salaire mensuel d'une employée : \all{das Gehalt} (et non \all{der Lohn}).
\item \all{Er hat eine kleine Firma \textbf{gegründet}.} « Il a créé une petite entreprise. » — Perfekt de \all{gründen} : participe \all{gegründet}.
\end{enumerate}
\end{exoresolu}

\begin{exercice}[À vous de jouer : mini-thème lexical]
Traduisez en allemand les phrases suivantes, en réutilisant le vocabulaire de cette section :
\begin{enumerate}
\item Mon père travaille dans une entreprise à Douala.
\item La vendeuse aide la cliente au marché.
\item Je postule pour un emploi dans une banque.
\item Le pain coûte 300 francs.
\item Le Cameroun exporte du cacao et du café.
\end{enumerate}
\end{exercice}

\begin{corrige}[Exercice : mini-thème lexical]
\begin{enumerate}
\item \all{Mein Vater arbeitet in einem Betrieb in Douala.} (ou : \all{bei einer Firma})
\item \all{Die Verkäuferin hilft der Kundin auf dem Markt.} — \all{helfen + Dativ} : \all{der Kundin}.
\item \all{Ich bewerbe mich um eine Stelle bei einer Bank.}
\item \all{Das Brot kostet 300 Franken.}
\item \all{Kamerun exportiert Kakao und Kaffee.}
\end{enumerate}
\end{corrige}

\begin{aretenir}[L'essentiel du lexique économique]
\begin{itemize}
\item Quatre champs lexicaux : \cle{travail} (\all{der Beruf}, \all{die Bewerbung}), \cle{argent} (\all{das Geld}, \all{der Preis}), \cle{commerce} (\all{der Markt}, \all{kaufen/verkaufen}), \cle{entreprise} (\all{die Firma}, \all{gründen}).
\item \all{der Lohn} = salaire horaire (ouvriers) ; \all{das Gehalt} = salaire mensuel (employés, cadres).
\item Constructions à mémoriser : \all{verdienen + Akk.}, \all{kosten + Akk.}, \all{arbeiten bei + Dat.}, \all{sich bewerben um + Akk.}
\item \all{der Kunde} est un nom faible : \all{dem Kunden}, \all{den Kunden}.
\item Faux amis : « actuel » = \all{aktuell} ; « fabriquer » = \all{herstellen/produzieren} ; « gagner de l'argent » = \all{Geld verdienen}.
\item Apprenez chaque nom avec article et pluriel : \all{der Markt, die Märkte}.
\end{itemize}
\end{aretenir}

Maintenant que ce lexique est bien installé, nous pourrons l'utiliser dans les sections suivantes pour réviser le comparatif, les phrases avec \all{wenn} et le futur, puis pour rédiger et traduire des textes économiques.

\coursec{Grammaire 1 : révision du comparatif et du superlatif}

Dans la section précédente, nous avons découvert le lexique de la vie économique : \all{der Preis, -e} (le prix), \all{der Lohn, \"-Löhne} (le salaire), \all{die Ware, -n} (la marchandise), \all{die Firma, Firmen} (l'entreprise). Or, dès qu'on parle d'économie, on compare sans cesse : un produit est \emph{moins cher} qu'un autre, une entreprise est \emph{plus grande}, un salaire est \emph{le plus élevé}. En allemand comme en français, ces comparaisons passent par le \cle{comparatif} et le \cle{superlatif} de l'adjectif. Révisons-les en profondeur, car ils sont indispensables pour rédiger et traduire des textes économiques — et ils tombent régulièrement aux épreuves de grammaire et de traduction.

\courssub{Observation : le comparatif dans un texte économique}

Lisons d'abord ce court article de presse économique (texte original) et soulignons mentalement toutes les formes de comparaison.

\begin{textebox}[„Zwei Märkte in Yaoundé“ — article original]
Wer in Yaoundé einkaufen will, hat die Wahl zwischen vielen Märkten. Der Mokolo-Markt ist größer als der Marché Central, und die Preise sind dort oft billiger. Eine Händlerin erklärt: „Bei uns sind die Tomaten frischer, und der Reis ist billiger als im Supermarkt.“ Doch der Supermarkt hat auch Vorteile: die Waren sind dort besser verpackt, und die Qualität ist oft höher. Am teuersten sind die importierten Produkte aus Europa. Viele Kunden sagen: „Das beste Angebot findet man auf dem Markt, aber die bequemste Atmosphäre gibt es im Supermarkt.“ Für die kamerunische Wirtschaft ist der lokale Markt so wichtig wie der moderne Handel. Der Kaffee aus der Westregion wird stärker exportiert als früher, und der Kakao bleibt das wichtigste Exportprodukt des Landes. Die jüngeren Händler nutzen jetzt das Handy, um die Preise schneller zu vergleichen. So wird der Handel immer moderner.
\end{textebox}

\textbf{Glossaire :} \all{die Wahl, -en} (le choix) ; \all{einkaufen} (faire les courses) ; \all{die Händlerin, -nen} (la vendeuse) ; \all{frisch} (frais, fraîche) ; \all{verpackt} (emballé) ; \all{importiert} (importé) ; \all{das Angebot, -e} (l'offre) ; \all{bequem} (confortable) ; \all{der Handel} (le commerce) ; \all{exportieren} (exporter) ; \all{vergleichen} (comparer) ; \all{die Atmosphäre, -n} (l'atmosphère).

\textbf{Questions de compréhension :}
\begin{enumerate}
\item Quel marché est plus grand que le Marché Central ?
\item Où le riz est-il moins cher ?
\item Quels produits sont les plus chers ?
\item Quel est le produit d'exportation le plus important du Cameroun ?
\end{enumerate}

Relevons les formes de comparaison du texte : \all{größer als}, \all{billiger als}, \all{frischer}, \all{höher}, \all{am teuersten}, \all{das beste Angebot}, \all{die bequemste Atmosphäre}, \all{so wichtig wie}, \all{stärker exportiert als}, \all{das wichtigste Exportprodukt}, \all{jüngere Händler}, \all{schneller}, \all{immer moderner}. Trois degrés apparaissent : le \cle{positif} (\all{billig}), le \cle{comparatif} (\all{billiger}) et le \cle{superlatif} (\all{am billigsten}). Voyons maintenant les règles de formation, degré par degré.

\courssub{Formation du comparatif : adjectif + -er}

\begin{propriete}[Formation du comparatif]
En allemand, le comparatif se forme en ajoutant la terminaison \cle{-er} à l'adjectif :
\begin{center}
\all{billig} $\rightarrow$ \all{billiger} ; \all{teuer} $\rightarrow$ \all{teurer} ; \all{schnell} $\rightarrow$ \all{schneller}
\end{center}
La comparaison d'inégalité s'exprime avec \cle{als} (que) :
\all{Diese Ware ist billiger als jene.} (Cette marchandise est moins chère que celle-là.)
\end{propriete}

\begin{attention}[Jamais « mehr » + adjectif !]
En français, on dit « plus cher », « moins cher », « plus important » : le comparatif se forme avec \emph{plus} ou \emph{moins}. En allemand, c'est \textbf{impossible} : le comparatif se forme \textbf{toujours} avec la terminaison \all{-er}, jamais avec \all{mehr} + adjectif.
\begin{itemize}
\item \textbf{Faux :} \all{*mehr billig}, \all{*mehr teuer}
\item \textbf{Correct :} \all{billiger}, \all{teurer}
\end{itemize}
Attention au faux ami : \all{mehr} signifie « plus (de quantité) » : \all{mehr Geld} (plus d'argent), \all{mehr Kunden} (plus de clients), mais jamais devant un adjectif pour former un comparatif. Pour dire « moins + adjectif », l'allemand utilise \all{weniger} + adjectif au positif : \all{weniger teuer} (moins cher), ou bien le comparatif de l'adjectif opposé : \all{billiger}.
\end{attention}

\begin{propriete}[Petites modifications orthographiques au comparatif]
\begin{itemize}
\item Les adjectifs en \all{-el} perdent le \all{e} devant \all{-er} : \all{dunkel} $\rightarrow$ \all{dunkler} (plus sombre).
\item Les adjectifs en \all{-er} perdent souvent le \all{e} : \all{teuer} $\rightarrow$ \all{teurer}.
\item Les adjectifs monosyllabiques en \all{a}, \all{o}, \all{u} prennent l'umlaut (voir plus bas) : \all{alt} $\rightarrow$ \all{älter}.
\end{itemize}
\end{propriete}

\courssub{Formation du superlatif : am ... -sten ou der/die/das ... -ste}

\begin{propriete}[Formation du superlatif]
Le superlatif allemand a deux formes selon l'emploi :
\begin{enumerate}
\item \textbf{Emploi prédicatif} (après le verbe \all{sein}) : \cle{am} + adjectif + \cle{-sten} :
\all{Dieses Produkt ist am billigsten.} (Ce produit est le moins cher.)
\item \textbf{Emploi attributif} (devant un nom) : \cle{der/die/das} + adjectif + \cle{-ste} + nom, avec déclinaison de l'adjectif :
\all{Deutschland hat die größte Wirtschaft Europas.} (L'Allemagne a la plus grande économie d'Europe.)
\end{enumerate}
Les adjectifs se terminant par \all{-d}, \all{-t}, \all{-s}, \all{-ß}, \all{-sch}, \all{-z} prennent \all{-esten} au superlatif pour faciliter la prononciation : \all{alt} $\rightarrow$ \all{am ältesten}, \all{heiß} $\rightarrow$ \all{am heißesten}, \all{kurz} $\rightarrow$ \all{am kürzesten}.
\end{propriete}

\begin{exemplebox}[Comparatif et superlatif dans le contexte économique]
\all{Der Euro ist stärker als der CFA-Franc.} (L'euro est plus fort que le franc CFA.)\\
\all{Im Januar sind die Preise am höchsten.} (En janvier, les prix sont les plus élevés.)\\
\all{Douala ist die größte Wirtschaftsstadt Kameruns.} (Douala est la plus grande ville économique du Cameroun.)\\
\all{Die kleinste Firma verdient am wenigsten.} (La plus petite entreprise gagne le moins.)
\end{exemplebox}

\courssub{Comparaison d'égalité : so ... wie ; d'inégalité : als}

\begin{propriete}[Égalité et inégalité]
\begin{itemize}
\item \textbf{Égalité} (aussi ... que) : \cle{so} + adjectif au positif + \cle{wie} :
\all{Kaffee ist so wichtig wie Kakao für Kamerun.} (Le café est aussi important que le cacao pour le Cameroun.)
\item \textbf{Négation de l'égalité} (pas aussi ... que) : \all{nicht so ... wie} :
\all{Der Markt ist nicht so teuer wie der Supermarkt.} (Le marché n'est pas aussi cher que le supermarché.)
\item \textbf{Inégalité} (plus/moins ... que) : comparatif + \cle{als} :
\all{Der Supermarkt ist teurer als der Markt.} (Le supermarché est plus cher que le marché.)
\end{itemize}
\end{propriete}

\begin{attention}[Après « als » : pas de virgule, attention au cas]
Deux pièges pour les francophones :
\begin{enumerate}
\item En français on écrit « plus cher \emph{que} lui » ; en allemand, \all{als} n'est \textbf{jamais} précédé d'une virgule : \all{billiger als jene} (et non \all{*billiger, als jene}).
\item Après \all{als} de comparaison, le cas suit la fonction logique : quand on compare deux sujets, le second élément est au \textbf{nominatif} :
\all{Er verdient mehr als ich.} (Il gagne plus que moi.) — et non \all{*als mich}.
Comparez : \all{Ich mag ihn mehr als sie.} (Je l'aime plus qu'elle — \all{sie} au nominatif, sujet sous-entendu « qu'elle ne l'aime ».)
\end{enumerate}
\end{attention}

\courssub{Adjectifs irréguliers et umlaut}

Certains adjectifs très fréquents ont des formes irrégulières qu'il faut connaître par cœur :

\begin{center}
\begin{tabular}{|l|l|l|}
\hline
\rowcolor{popblueL} \textbf{Positif} & \textbf{Comparatif} & \textbf{Superlatif} \\
\hline
\all{gut} (bon) & \all{besser} & \all{am besten} \\
\hline
\all{viel} (beaucoup) & \all{mehr} & \all{am meisten} \\
\hline
\all{gern} (volontiers) & \all{lieber} & \all{am liebsten} \\
\hline
\all{hoch} (haut, élevé) & \all{höher} & \all{am höchsten} \\
\hline
\all{nah} (proche) & \all{näher} & \all{am nächsten} \\
\hline
\end{tabular}
\end{center}

\begin{propriete}[Umlaut au comparatif et au superlatif]
Les adjectifs \textbf{monosyllabiques} dont la voyelle radicale est \cle{a}, \cle{o} ou \cle{u} prennent généralement l'\textbf{umlaut} au comparatif et au superlatif :
\all{alt} $\rightarrow$ \all{älter} $\rightarrow$ \all{am ältesten} ; \all{jung} $\rightarrow$ \all{jünger} $\rightarrow$ \all{am jüngsten} ; \all{stark} $\rightarrow$ \all{stärker} $\rightarrow$ \all{am stärksten} ; \all{lang} $\rightarrow$ \all{länger} $\rightarrow$ \all{am längsten} ; \all{groß} $\rightarrow$ \all{größer} $\rightarrow$ \all{am größten} ; \all{kurz} $\rightarrow$ \all{kürzer} $\rightarrow$ \all{am kürzesten}.
\end{propriete}

\begin{exemplebox}[Exemples traduits]
\all{Die jüngeren Arbeiter verdienen weniger.} (Les travailleurs plus jeunes gagnent moins.)\\
\all{Der Lohn ist höher als früher.} (Le salaire est plus élevé qu'avant.)\\
\all{Ich kaufe lieber auf dem Markt.} (J'achète de préférence au marché.)\\
\all{Die stärkste Wirtschaft Europas ist Deutschland.} (L'économie la plus forte d'Europe, c'est l'Allemagne.)
\end{exemplebox}

\begin{propriete}[« De plus en plus » : immer + comparatif]
Pour exprimer une progression (« de plus en plus »), l'allemand utilise \cle{immer} + comparatif :
\all{Die Preise werden immer höher.} (Les prix deviennent de plus en plus élevés.)
\all{Der Handel wird immer moderner.} (Le commerce devient de plus en plus moderne.)
Ne traduisez jamais « de plus en plus cher » par \all{*mehr und mehr teuer}.
\end{propriete}

\courssub{Emploi attributif : la déclinaison du comparatif}

Quand le comparatif ou le superlatif est placé \textbf{devant un nom} (emploi attributif), il se \textbf{décline} comme tout adjectif allemand, selon le genre, le cas et l'article :

\begin{center}
\begin{tabular}{|l|l|l|}
\hline
\rowcolor{popblueL} \textbf{Français} & \textbf{Deutsch} & \textbf{Remarque} \\
\hline
une meilleure offre & \all{ein besseres Angebot} & neutre, nominatif, article indéfini \\
\hline
la marchandise plus chère & \all{die teurere Ware} & féminin, nominatif, article défini \\
\hline
le prix le plus bas & \all{der niedrigste Preis} & masculin, nominatif \\
\hline
avec un salaire plus élevé & \all{mit einem höheren Lohn} & masculin, datif \\
\hline
les plus grandes entreprises & \all{die größten Firmen} & pluriel, nominatif \\
\hline
\end{tabular}
\end{center}

\begin{attention}[La déclinaison s'ajoute à la terminaison du degré]
N'oubliez jamais que \all{-er} (comparatif) et \all{-ste} (superlatif) ne dispensent \textbf{pas} de la déclinaison adjectivale : on ajoute la désinence après le degré. Ainsi : \all{ein besser-\textbf{es} Angebot}, \all{der teurer-\textbf{e} Preis}, \all{mit dem best-\textbf{en} Angebot}. C'est une des erreurs les plus fréquentes des francophones, car en français « meilleur » est invariable en genre au masculin.
\end{attention}

\courssub{Tableau récapitulatif des formes}

\begin{center}
\begin{tabular}{|l|l|l|l|}
\hline
\rowcolor{popblueL} \textbf{Positif} & \textbf{Comparatif} & \textbf{Superlatif} & \textbf{Particularité} \\
\hline
\all{billig} & \all{billiger} & \all{am billigsten} & régulier \\
\hline
\all{teuer} & \all{teurer} & \all{am teuersten} & régulier (perte du \all{e}) \\
\hline
\all{gut} & \all{besser} & \all{am besten} & irrégulier \\
\hline
\all{viel} & \all{mehr} & \all{am meisten} & irrégulier \\
\hline
\all{hoch} & \all{höher} & \all{am höchsten} & irrégulier (perte du \all{c}) \\
\hline
\all{alt} & \all{älter} & \all{am ältesten} & umlaut \\
\hline
\all{gern} & \all{lieber} & \all{am liebsten} & irrégulier \\
\hline
\end{tabular}
\end{center}

\begin{popfigure}
\begin{center}
\begin{tikzpicture}[scale=1.0]
\draw[thick, popline] (0,0) -- (2,0) -- (2,1) -- (4,1) -- (4,2) -- (6,2);
\node[draw=popblue, fill=popblueL, rounded corners, align=center, minimum width=1.9cm, minimum height=0.8cm] at (1,-0.75) {\all{billig}\\ \footnotesize positif};
\node[draw=poporange, fill=poporangeL, rounded corners, align=center, minimum width=1.9cm, minimum height=0.8cm] at (3,0.25) {\all{billiger}\\ \footnotesize comparatif};
\node[draw=poppurple, fill=poppurpleL, rounded corners, align=center, minimum width=2.3cm, minimum height=0.8cm] at (5,1.25) {\all{am billigsten}\\ \footnotesize superlatif};
\node[align=center, font=\footnotesize, popmuted] at (1,-1.5) {bon marché};
\node[align=center, font=\footnotesize, popmuted] at (3,-0.5) {moins cher};
\node[align=center, font=\footnotesize, popmuted] at (5,0.5) {le moins cher};
\end{tikzpicture}
\end{center}
\legende{Les trois degrés de l'adjectif : un escalier montant, du positif au superlatif.}
\end{popfigure}

\courssub{Exercice résolu}

\begin{exoresolu}[Comparer des prix au marché]
\textbf{Énoncé.} Complétez avec le comparatif ou le superlatif de l'adjectif entre parenthèses, puis traduisez en français.
\begin{enumerate}
\item Auf dem Markt sind die Tomaten \_\_\_\_\_\_\_\_ als im Supermarkt. (billig)
\item Der Mokolo-Markt ist \_\_\_\_\_\_\_\_ Markt der Stadt. (groß)
\item Diese Firma zahlt einen \_\_\_\_\_\_\_\_ Lohn als jene. (hoch)
\item Kaffee ist so wichtig \_\_\_\_\_\_\_\_ Kakao.
\item Das ist das \_\_\_\_\_\_\_\_ Angebot. (gut)
\end{enumerate}
\tcblower
\textbf{Solution détaillée.}
\begin{enumerate}
\item \all{billiger} : comparatif régulier (billig + -er) suivi de \all{als}. Traduction : « Au marché, les tomates sont moins chères qu'au supermarché. »
\item \all{der größte} : superlatif attributif de \all{groß} (irrégulier avec umlaut : groß $\rightarrow$ größer $\rightarrow$ am größten), décliné au masculin nominatif après l'article défini. Traduction : « Le marché Mokolo est le plus grand marché de la ville. »
\item \all{höheren} : comparatif de \all{hoch} (irrégulier : perte du \all{c} et umlaut), décliné au masculin accusatif après \all{einen}. Traduction : « Cette entreprise paie un salaire plus élevé que celle-là. »
\item \all{wie} : comparaison d'égalité \all{so ... wie}. Traduction : « Le café est aussi important que le cacao. »
\item \all{beste} : superlatif attributif de \all{gut} (irrégulier : gut $\rightarrow$ besser $\rightarrow$ am besten), décliné au neutre accusatif après \all{das}. Traduction : « C'est la meilleure offre. »
\end{enumerate}
\end{exoresolu}

\begin{exercice}[À vous de jouer]
Traduisez en allemand :
\begin{enumerate}
\item Le marché est moins cher que le supermarché.
\item L'Allemagne a la plus grande économie d'Europe.
\item Il gagne plus que moi.
\item Le café est aussi important que le cacao.
\item Cette entreprise offre un meilleur salaire.
\end{enumerate}
\end{exercice}

\begin{corrige}[Exercice « À vous de jouer »]
\begin{enumerate}
\item \all{Der Markt ist billiger als der Supermarkt.} (comparatif régulier + \all{als} sans virgule)
\item \all{Deutschland hat die größte Wirtschaft Europas.} (superlatif attributif de \all{groß}, féminin accusatif)
\item \all{Er verdient mehr als ich.} (\all{ich} au nominatif, fonction logique de sujet)
\item \all{Kaffee ist so wichtig wie Kakao.} (égalité : \all{so ... wie})
\item \all{Diese Firma bietet einen besseren Lohn.} (comparatif de \all{gut} décliné : masculin accusatif)
\end{enumerate}
\end{corrige}

\begin{aretenir}[L'essentiel du comparatif et du superlatif]
\begin{itemize}
\item Comparatif : adjectif + \all{-er}, jamais \all{mehr} + adjectif ; inégalité avec \all{als} (sans virgule).
\item Superlatif : \all{am} + adjectif + \all{-sten} (prédicatif) ou \all{der/die/das} + adjectif + \all{-ste} (attributif, avec déclinaison).
\item Égalité : \all{so ... wie} ; négation : \all{nicht so ... wie}.
\item Progression : \all{immer} + comparatif (\all{immer teurer} = de plus en plus cher).
\item Irréguliers à mémoriser : \all{gut – besser – am besten} ; \all{viel – mehr – am meisten} ; \all{gern – lieber – am liebsten} ; \all{hoch – höher – am höchsten}.
\item Umlaut pour les monosyllabes en a/o/u : \all{alt – älter}, \all{jung – jünger}, \all{stark – stärker}.
\item Devant un nom, le comparatif et le superlatif se déclinent : \all{ein besseres Angebot}, \all{die größte Wirtschaft}.
\end{itemize}
\end{aretenir}

Ces outils de comparaison nous seront immédiatement utiles dans la section suivante, où nous réviserons les phrases conditionnelles avec \all{wenn} et le futur — deux structures fréquentes dès qu'on parle de projets économiques et d'avenir professionnel.

\coursec{Grammaire 2 : wenn-Sätze et futur}

Dans la section précédente, nous avons révisé le comparatif et le superlatif, qui nous permettent de comparer des prix, des salaires ou des produits (\all{billiger als}, \all{am teuersten}). Mais la vie économique, ce n'est pas seulement comparer : c'est aussi \cle{envisager l'avenir} et \cle{formuler des conditions}. Un commerçant de Douala dit : « Si les prix baissent, j'achèterai plus de marchandises. » Un élève de Première déclare : « L'année prochaine, je créerai ma petite entreprise. » Pour exprimer cela en allemand, il faut maîtriser deux outils essentiels : la \cle{subordonnée en wenn} et le \cle{futur I}. C'est l'objet de cette section.

\courssub{Observation : un texte pour découvrir}

Lisons d'abord ce court texte économique. Observez les phrases soulignées mentalement : où se trouve le verbe conjugué ? Quel mot introduit la condition ?

\begin{textebox}[Zukunftspläne in Yaoundé — Projets d’avenir à Yaoundé]
Marie-Claire ist 17 Jahre alt und besucht ein Gymnasium in Yaoundé. Sie hat große Pläne. „Wenn ich das Abitur habe, werde ich Wirtschaft studieren“, sagt sie. „Und wenn ich genug Geld gespart habe, werde ich eine kleine Firma gründen.“ Ihr Bruder Paul denkt anders: „Wenn die Wirtschaft in Kamerun wächst, gibt es mehr Arbeitsplätze für junge Leute. Dann werde ich vielleicht Ingenieur.“ Die Mutter der beiden ist Händlerin auf dem Markt. Sie erklärt: „Wenn die Preise steigen, kaufen die Kunden weniger. Aber wenn die Qualität gut ist, kommen die Kunden immer zurück.“ Am Abend spricht die Familie über die Zukunft. „Als ich jung war“, sagt der Vater, „war das Leben einfacher. Aber heute gibt es mehr Chancen. Wenn ihr fleißig arbeitet, werdet ihr Erfolg haben.“ Marie-Claire lächelt: „Nächstes Jahr werde ich schon anfangen, Geld zu sparen. Wenn ich 25 bin, werde ich meine eigene Boutique haben!“
\end{textebox}

\textbf{Glossaire :} das Abitur (baccalauréat) ; sparen (épargner) ; die Firma, die Firmen (l'entreprise) ; gründen (fonder, créer) ; die Ware, die Waren (la marchandise) ; die Wirtschaft (l'économie) ; wachsen (croître) ; der Arbeitsplatz, die Arbeitsplätze (l'emploi, le poste) ; die Händlerin (la commerçante) ; steigen (augmenter) ; die Qualität (la qualité) ; die Chance, die Chancen (l'opportunité) ; fleißig (travailleur, appliqué) ; der Erfolg (le succès) ; die Boutique (la boutique).

\textbf{Questions de compréhension :}
\begin{enumerate}
\item Que veut faire Marie-Claire après le baccalauréat ?
\item Que dit la mère sur les prix et les clients ?
\item Quelle phrase du texte parle d'un événement passé unique ? Quel mot l'introduit ?
\end{enumerate}

\textbf{Observations grammaticales.} Relevons trois phrases :
\begin{itemize}
\item \all{Wenn ich das Abitur habe, werde ich Wirtschaft studieren.} « Quand j'aurai le bac, j'étudierai l'économie. »
\item \all{Wenn die Preise steigen, kaufen die Kunden weniger.} « Si les prix augmentent, les clients achètent moins. »
\item \all{Als ich jung war, war das Leben einfacher.} « Quand j'étais jeune, la vie était plus simple. »
\end{itemize}
On remarque : (1) dans la subordonnée introduite par \all{wenn} ou \all{als}, le verbe conjugué (\all{habe}, \all{steigen}, \all{war}) se trouve \textbf{à la fin} ; (2) quand la subordonnée est en tête, le verbe de la principale vient \textbf{immédiatement après la virgule} (\all{..., werde ich...}) ; (3) le futur se construit avec \all{werden} + un infinitif placé en fin de phrase (\all{... Wirtschaft studieren}).

\courssub{Le wenn-Satz temporel : wenn = quand, lorsque}

\begin{propriete}[La subordonnée temporelle en wenn]
\all{Wenn} signifie « quand, lorsque » pour :
\begin{itemize}
\item un événement \textbf{futur} : \all{Wenn ich 25 bin, werde ich eine Boutique haben.} « Quand j'aurai 25 ans, j'aurai ma propre boutique. »
\item une \textbf{habitude} ou une répétition (présent ou passé) : \all{Wenn die Qualität gut ist, kommen die Kunden zurück.} « Quand la qualité est bonne, les clients reviennent. »
\end{itemize}
\textbf{Forme :} \all{wenn} + sujet + compléments + \textbf{verbe conjugué en dernière position}. Une virgule sépare toujours la subordonnée de la principale.
\end{propriete}

\begin{attention}[wenn ou als ? Le piège n° 1 des francophones]
En français, « quand » s'emploie partout. En allemand, il faut choisir :
\begin{itemize}
\item \all{wenn} : présent, futur, ou \textbf{habitude au passé} (« quand... chaque fois ») : \all{Wenn ich Zeit habe, gehe ich zum Markt.} « Quand j'ai le temps, je vais au marché. » / \all{Wenn wir in den Ferien nach Bafoussam fuhren, ...} « Quand nous allions à Bafoussam pendant les vacances... » (répétition).
\item \all{als} : \textbf{événement unique au passé} (récit) : \all{Als ich Kind war, wohnten wir in Douala.} « Quand j'étais enfant, nous habitions à Douala. »
\end{itemize}
On ne peut donc JAMAIS dire \all{als} pour un fait futur, ni \all{wenn} pour un événement passé unique (sauf habitude répétée).
\end{attention}

\courssub{Le wenn-Satz conditionnel : wenn = si}

\begin{propriete}[La condition réelle avec wenn]
Pour exprimer une condition réelle (« si + présent » en français), l'allemand utilise \all{wenn} + présent :
\all{Wenn ich Geld habe, kaufe ich ein Auto.} « Si j'ai de l'argent, j'achète une voiture. »
La principale peut être au présent (conséquence habituelle) ou au futur I (conséquence envisagée) :
\all{Wenn die Wirtschaft wächst, wird es mehr Arbeitsplätze geben.} « Si l'économie croît, il y aura plus d'emplois. »
\end{propriete}

\begin{exemplebox}[Exemples traduits]
\begin{itemize}
\item \all{Wenn die Preise steigen, kaufen die Kunden weniger.} « Si les prix augmentent, les clients achètent moins. »
\item \all{Wenn du fleißig arbeitest, bekommst du ein besseres Gehalt.} « Si tu travailles dur, tu obtiendras un meilleur salaire. »
\item \all{Wenn es morgen regnet, bleibt der Markt fast leer.} « S'il pleut demain, le marché restera presque vide. »
\item \all{Ich spare Geld, wenn ich nicht zu viel ausgebe.} « J'épargne de l'argent si je ne dépense pas trop. » (Subordonnée après la principale : la virgule reste obligatoire avant \all{wenn}.)
\end{itemize}
\end{exemplebox}

\begin{propriete}[Ordre des mots : la règle d'or]
Deux positions possibles :
\begin{enumerate}
\item \textbf{Principale d'abord :} verbe en position 2, puis \all{wenn}-Satz avec verbe final. \all{Ich kaufe ein Auto, wenn ich Geld habe.}
\item \textbf{Subordonnée d'abord :} la subordonnée entière occupe la position 1 ; le verbe conjugué de la principale vient \textbf{immédiatement après la virgule} (inversion sujet-verbe) : \all{Wenn ich Geld habe, \textbf{kaufe ich} ein Auto.} (et non \all{ich kaufe}).
\end{enumerate}
La \textbf{virgule} entre les deux propositions est \textbf{toujours obligatoire} en allemand, contrairement au français.
\end{propriete}

\begin{popfigure}
\begin{tikzpicture}[font=\small, >=Stealth]
\node[draw=popblue, fill=popblueL, rounded corners, align=center, minimum width=5.5cm] (sub) at (0,0) {\textbf{Subordonnée (position 1)}\\ Wenn die Wirtschaft \textbf{wächst},};
\node[draw=poporange, fill=poporangeL, rounded corners, align=center, minimum width=5.5cm] (main) at (6.5,0) {\textbf{Principale}\\ \textbf{gibt} es mehr Arbeitsplätze.};
\draw[->, poppurple, thick] (sub.east) -- node[above, align=center]{Inversion : le verbe\\ de la principale en\\ position 2 (après la virgule)} (main.west);
\node[draw=popgreen, fill=popgreenL, rounded corners, align=center] at (3.25,-1.8) {Position 1 = toute la subordonnée \quad Position 2 = verbe conjugué};
\end{tikzpicture}
\legende{Structure de la phrase : \all{Wenn die Wirtschaft wächst, gibt es mehr Arbeitsplätze.}}
\end{popfigure}

\courssub{Le futur I : werden + infinitif}

\begin{propriete}[Formation du futur I]
Le futur I se forme avec l'auxiliaire \all{werden} \textbf{conjugué} (position 2) + l'\textbf{infinitif} du verbe principal placé \textbf{en toute fin de phrase} :
\all{Ich werde nächstes Jahr eine Firma gründen.} « Je créerai une entreprise l'année prochaine. »
\end{propriete}

\begin{center}
\begin{tabular}{|l|l|l|}
\hline
\rowcolor{popblueL} Personne & werden & Exemple (Infinitiv am Ende) \\
\hline
ich & werde & Ich \textbf{werde} Wirtschaft \textbf{studieren}. \\
\hline
du & wirst & Du \textbf{wirst} ein gutes Gehalt \textbf{bekommen}. \\
\hline
er/sie/es & wird & Sie \textbf{wird} eine Boutique \textbf{eröffnen}. \\
\hline
wir & werden & Wir \textbf{werden} nach Deutschland \textbf{reisen}. \\
\hline
ihr & werdet & Ihr \textbf{werdet} Erfolg \textbf{haben}. \\
\hline
sie/Sie & werden & Sie \textbf{werden} mehr Kunden \textbf{finden}. \\
\hline
\end{tabular}
\end{center}

\textbf{Emplois du futur I :}
\begin{itemize}
\item \textbf{Prévision} : \all{Die Wirtschaft wird wachsen.} « L'économie croîtra. »
\item \textbf{Projet, intention} : \all{Ich werde eine Firma gründen.} « Je vais créer / je créerai une entreprise. »
\item \textbf{Promesse} : \all{Ich werde dich morgen anrufen.} « Je t'appellerai demain. »
\item \textbf{Supposition} (souvent avec \all{wohl}, \all{wahrscheinlich}) : \all{Er wird wohl krank sein.} « Il est probablement malade. »
\end{itemize}

\begin{attention}[Deux pièges du futur]
\begin{enumerate}
\item \all{werden} se conjugue (jamais \all{ich werden} !) et l'infinitif va \textbf{en toute fin}, même après de longs compléments : \all{Wir werden nächstes Jahr in Yaoundé ein neues Geschäft eröffnen.}
\item En allemand, le \textbf{présent} suffit souvent pour exprimer le futur quand un complément de temps est présent : \all{Morgen fahre ich nach Berlin.} « Demain, je vais à Berlin. » Le futur I n'est donc pas obligatoire ; il insiste sur le projet ou la prévision.
\end{enumerate}
\end{attention}

\begin{center}
\begin{tabularx}{\linewidth}{|X|X|X|}
\hline
\rowcolor{popblueL} Français & Deutsch & Remarque \\
\hline
Si j'ai de l'argent, j'achète une voiture. & Wenn ich Geld habe, kaufe ich ein Auto. & wenn + présent ; verbe final ; inversion \\
\hline
Quand j'étais enfant, ... & Als ich Kind war, ... & passé unique = als \\
\hline
Quand j'aurai le bac, j'étudierai. & Wenn ich das Abitur habe, werde ich studieren. & futur des deux côtés en français ; présent + futur en allemand \\
\hline
Demain, je vais au marché. & Morgen gehe ich zum Markt. & présent allemand + complément de temps \\
\hline
Je créerai une entreprise. & Ich werde eine Firma gründen. & werden conjugué + infinitif final \\
\hline
\end{tabularx}
\end{center}

\begin{popfigure}
\begin{tikzpicture}[font=\small, >=Stealth]
\draw[->, thick, popdark] (0,0) -- (11,0);
\foreach \x/\lab in {1/{Vergangenheit}, 5.5/{Gegenwart}, 10/{Zukunft}} {
  \draw[thick, popdark] (\x,0.15) -- (\x,-0.15);
  \node[below, align=center] at (\x,-0.2) {\lab};
}
\node[draw=popblue, fill=popblueL, rounded corners, align=center] at (5.5,1.1) {Präsens\\ \all{Ich arbeite heute.}};
\node[draw=poporange, fill=poporangeL, rounded corners, align=center] at (10,1.1) {Futur I\\ \all{Ich werde arbeiten.}};
\draw[->, poppurple, thick] (6.6,1.1) -- node[above]{\all{werden} + Infinitiv} (8.6,1.1);
\end{tikzpicture}
\legende{Frise des temps : du présent (\all{Präsens}) vers l'avenir (\all{Futur I}).}
\end{popfigure}

\begin{methode}[Traduire « si » et « quand » du français vers l'allemand]
\begin{enumerate}
\item Le mot est-il « si » (condition) ? → \all{wenn} + présent allemand.
\item Le mot est-il « quand » ? Posez la question : événement futur ou habituel → \all{wenn} ; événement \textbf{unique au passé} → \all{als} + Präteritum/Perfekt.
\item Placez le verbe conjugué à la fin de la subordonnée.
\item Mettez la virgule, puis vérifiez l'inversion si la subordonnée est en tête.
\item Pour le futur : conjuguez \all{werden} et envoyez l'infinitif en fin de phrase.
\end{enumerate}
\end{methode}

\begin{exoresolu}[Thème : conditions et projets]
Traduisez en allemand : 1) « Si les prix baissent, j'achèterai plus de marchandises. » 2) « Quand j'étais jeune, je travaillais au marché de Douala. » 3) « L'année prochaine, nous ouvrirons un magasin à Yaoundé. »
\tcblower
1) Condition au présent → \all{wenn} ; verbe \all{fallen} en fin de subordonnée ; futur I dans la principale avec inversion : \all{Wenn die Preise fallen, werde ich mehr Waren kaufen.} 2) Passé unique → \all{als} ; Präteritum : \all{Als ich jung war, arbeitete ich auf dem Markt in Douala.} 3) Futur I, infinitif \all{eröffnen} en fin : \all{Nächstes Jahr werden wir ein Geschäft in Yaoundé eröffnen.}
\end{exoresolu}

\begin{exercice}[À vous]
Traduisez : 1) « Si tu économises ton argent, tu pourras voyager. » 2) « Quand il pleut, le marché est vide. » 3) « Quand nous avons visité l'Allemagne, nous avons vu beaucoup d'usines. » 4) « Dans dix ans, je serai (Futur : \all{sein}) un homme d'affaires. »
\end{exercice}

\begin{corrige}[Exercice : conditions et futur]
1) \all{Wenn du dein Geld sparst, wirst du reisen können.} (ou \all{..., kannst du reisen.} pour une capacité présente/future). 2) Habitude → \all{Wenn es regnet, ist der Markt leer.} 3) Passé unique → \all{Als wir Deutschland besuchten, sahen wir viele Fabriken.} 4) \all{In zehn Jahren werde ich Geschäftsmann sein.}
\end{corrige}

\begin{aretenir}[L'essentiel]
\begin{itemize}
\item \all{wenn} = « si » (condition) ou « quand » (futur, habitude) ; \all{als} = « quand » au passé unique.
\item Subordonnée : verbe conjugué \textbf{en fin} ; virgule obligatoire ; si la subordonnée ouvre la phrase, \textbf{inversion} dans la principale.
\item Futur I : \all{werden} conjugué + infinitif \textbf{en toute fin} ; le présent + complément de temps peut le remplacer.
\end{itemize}
\end{aretenir}

\coursec{Types de textes : petite annonce, lettre, compte rendu}

Dans la section précédente, nous avons révisé les phrases conditionnelles avec \all{wenn} et le futur : ces outils grammaticaux nous serviront directement ici, car une lettre de motivation contient souvent des phrases comme \all{Wenn ich die Stelle bekomme, werde ich sehr motiviert arbeiten.} (« Si j'obtiens le poste, je travaillerai avec beaucoup de motivation. »). Nous allons maintenant apprendre à produire trois types de textes très fréquents dans la vie économique : la petite annonce, la lettre formelle (dont la lettre de motivation) et le compte rendu. Chacun a ses propres règles de structure et de style.

\courssub{La petite annonce (die Kleinanzeige)}

\begin{definition}[Die Kleinanzeige]
La \cle{petite annonce} (\all{die Kleinanzeige, -n}) est un texte très court publié dans un journal, sur un panneau d'affichage ou sur Internet pour vendre, acheter, louer ou proposer quelque chose. Son style est \cle{télégraphique} : on supprime les articles, les pronoms et les verbes inutiles pour aller à l'essentiel.
\end{definition}

Observons d'abord un modèle authentique, adapté au contexte camerounais :

\begin{textebox}[Modèle de petite annonce]
\textbf{Zu verkaufen: Fahrrad}

Fahrrad, fast neu, sehr guter Zustand, nur 50 000 FCFA. Ideal für Schüler. Tel.: 6XX XX XX XX. Abholung in Yaoundé, Quartier Nkolndongo.

\textbf{Zimmer zu vermieten}

Ruhiges Zimmer, 25 000 FCFA/Monat, VB. Zentrum von Douala, 5 Min. vom Markt. Bei Interesse: Frau Mbarga, Tel.: 6XX XX XX XX.
\end{textebox}

Glossaire : \all{der Zustand} (l'état), \all{die Abholung} (l'enlèvement, le retrait), \all{zu vermieten} (à louer), \all{ruhig} (calme), \all{bei Interesse} (en cas d'intérêt).

\begin{propriete}[Structure et style de la petite annonce]
Une petite annonce allemande contient cinq éléments, dans cet ordre :
\begin{enumerate}
\item un \cle{titre} qui annonce l'objet : \all{Zu verkaufen: ...} (à vendre), \all{Zu vermieten: ...} (à louer), \all{Suche ...} (je cherche) ;
\item l'\cle{objet} et sa \cle{description} brève : \all{fast neu, sehr guter Zustand} ;
\item le \cle{prix} : \all{nur 50 000 FCFA} ;
\item le \cle{contact} : \all{Tel.: ...} ;
\item parfois le \cle{lieu} : \all{Abholung in Yaoundé}.
\end{enumerate}
Le style télégraphique supprime les mots de liaison : au lieu de \all{Ich verkaufe ein Fahrrad, das fast neu ist}, on écrit simplement \all{Fahrrad, fast neu}.
\end{propriete}

\begin{aretenir}[Abréviations typiques des petites annonces]
\begin{center}
\begin{tabularx}{\linewidth}{|l|X|X|}\hline
\rowcolor{popblueL} \textbf{Abréviation} & \textbf{Forme complète} & \textbf{Français} \\ \hline
z.B. & zum Beispiel & par exemple \\ \hline
Tel. & Telefon & téléphone \\ \hline
usw. & und so weiter & et cetera \\ \hline
VB & Verhandlungsbasis & prix à débattre \\ \hline
ca. & circa & environ \\ \hline
Min. & Minuten & minutes \\ \hline
zzgl. & zuzüglich & en plus de \\ \hline
\end{tabularx}
\end{center}
L'abréviation \all{VB} (\all{Verhandlungsbasis}) est très importante : elle signifie que le prix peut être discuté, comme au marché de Mokolo ou de Sandaga ! La préposition \all{ab} (à partir de) figure aussi souvent : \all{ab 1. Mai} (à partir du 1er mai).
\end{aretenir}

\begin{attention}[Piège : le style télégraphique a des limites]
Même dans une annonce, les \cle{noms} gardent leur majuscule et les \cle{adjectifs attributs} restent déclinés : on écrit \all{sehr guter Zustand} et non \all{sehr gut Zustand}. Et attention au faux ami : \all{die Anzeige} signifie « l'annonce » mais aussi « la dénonciation » (\all{Anzeige erstatten} = porter plainte) ; le mot français « annonce » au sens publicitaire se dit aussi \all{die Werbung}.
\end{attention}

\courssub{La lettre formelle (der Brief)}

Observons le début d'une lettre de candidature rédigée par un élève de Première :

\begin{textebox}[Modèle de début de lettre]
Amina Fotso \\
Rue 1.234, Quartier Bastos \\
Yaoundé, Kamerun

Kamda GmbH \\
Herrn Jean Kamda \\
Avenue de la Liberté \\
Douala

Yaoundé, den 15. März 2025

\textbf{Bewerbung um die Stelle als Verkäuferin}

Sehr geehrte Damen und Herren,

ich schreibe Ihnen, weil ich mich um die Stelle als Verkäuferin bewerben möchte. Ich habe Ihre Anzeige in der Zeitung gelesen und ich interessiere mich sehr für diese Arbeit. Ich spreche gut Deutsch und Französisch, und ich habe schon ein Jahr im Geschäft meiner Tante gearbeitet.

Wenn Sie mir die Stelle geben, werde ich pünktlich und motiviert arbeiten. Ich stehe ab sofort zur Verfügung.

Mit freundlichen Grüßen

Amina Fotso
\end{textebox}

Glossaire : \all{die Stelle, -n} (le poste), \all{sich bewerben um + Akk.} (postuler à), \all{das Geschäft, -e} (le magasin), \all{pünktlich} (ponctuel), \all{zur Verfügung stehen} (être disponible).

\begin{propriete}[La structure de la lettre formelle allemande]
Une lettre formelle allemande suit un plan fixe, qu'il faut respecter scrupuleusement :
\begin{enumerate}
\item \all{der Absender} (l'expéditeur) : nom et adresse, en haut à gauche ;
\item \all{der Empfänger} (le destinataire) : nom et adresse, en dessous ;
\item \all{der Ort und das Datum} (le lieu et la date) : à droite, sous la forme \all{Yaoundé, den 15. März 2025} ;
\item \all{die Betreffzeile} (la ligne d'objet), en gras, sans verbe : \all{Bewerbung um die Stelle als Verkäuferin} ;
\item \all{die Anrede} (la formule d'appel) : \all{Sehr geehrte Damen und Herren,} ;
\item \all{der Text} (le corps de la lettre), qui commence par une \cle{minuscule} après la virgule ;
\item \all{der Gruß} (la formule de politesse) : \all{Mit freundlichen Grüßen} ;
\item \all{die Unterschrift} (la signature).
\end{enumerate}
\end{propriete}

\begin{popfigure}
\begin{center}
\begin{tikzpicture}[
block/.style={rectangle, draw=popblue, fill=popblueL, rounded corners=2pt, minimum width=6.5cm, minimum height=0.8cm, align=center, font=\small},
arr/.style={-{Stealth[length=2.5mm]}, thick, popdark}]
\node[block] (a) {Absender (expéditeur)};
\node[block, below=0.35cm of a] (b) {Empfänger (destinataire)};
\node[block, below=0.35cm of b] (c) {Ort und Datum : Yaoundé, den 15. März 2025};
\node[block, below=0.35cm of c] (d) {Betreff + Anrede : Sehr geehrte Damen und Herren,};
\node[block, below=0.35cm of d] (e) {Text (corps de la lettre, minuscule après la virgule)};
\node[block, below=0.35cm of e] (f) {Gruß : Mit freundlichen Grüßen};
\node[block, below=0.35cm of f] (g) {Unterschrift (signature)};
\draw[arr] (a) -- (b);
\draw[arr] (b) -- (c);
\draw[arr] (c) -- (d);
\draw[arr] (d) -- (e);
\draw[arr] (e) -- (f);
\draw[arr] (f) -- (g);
\end{tikzpicture}
\end{center}
\legende{Les sept blocs de la lettre formelle allemande, dans l'ordre obligatoire.}
\end{popfigure}

\begin{attention}[Date et formule de politesse : deux pièges classiques]
En allemand, la date s'écrit avec l'article accusatif : \all{Yaoundé, den 15. März 2025} (littéralement « le 15 mars »), avec un point après le chiffre car c'est un ordinal (\all{den 15.} = \all{den fünfzehnten}). De plus, la formule \all{Mit freundlichen Grüßen} ne prend \cle{pas de virgule} après elle, contrairement au français « Cordialement, ». Enfin, après \all{Sehr geehrte Damen und Herren,} la phrase suivante commence par une \cle{minuscule} : \all{..., ich schreibe Ihnen...}.
\end{attention}

\begin{savaistu}[MfG : la politesse en trois lettres]
Dans les e-mails professionnels allemands, la formule \all{Mit freundlichen Grüßen} est très souvent abrégée en \all{MfG}. On la voit partout dans la correspondance de bureau à Berlin, Vienne ou Zurich. Mais attention : dans une lettre formelle imprimée ou dans un examen, il faut écrire la formule complète, sans abréviation !
\end{savaistu}

\begin{savaistu}[Formes juridiques : SARL vs GmbH]
Au Cameroun, la forme juridique courante est la \cle{SARL} (Société à Responsabilité Limitée). En Allemagne, l'équivalent est la \cle{GmbH} (\all{Gesellschaft mit beschränkter Haftung}). Dans une lettre rédigée en allemand, on utilise la dénomination allemande : \all{Kamda GmbH} et non \all{Kamda SARL}.
\end{savaistu}

\courssub{La lettre de motivation (das Bewerbungsschreiben)}

\begin{definition}[Das Bewerbungsschreiben]
La \cle{lettre de motivation} (\all{das Bewerbungsschreiben, -}) est une lettre formelle particulière : elle accompagne un CV pour demander un emploi, un stage (\all{das Praktikum}) ou une formation. Elle suit la structure du \all{Brief}, mais son corps obéit à un plan en quatre parties.
\end{definition}

\begin{propriete}[Le corps de la lettre de motivation en quatre parties]
\begin{enumerate}
\item \cle{Présentation} : qui je suis et pourquoi j'écris — \all{Ich schreibe Ihnen, weil ich mich um die Stelle als ... bewerben möchte.} ;
\item \cle{Compétences} : ce que je sais faire — \all{Ich spreche gut Deutsch und Französisch. Ich habe schon ein Jahr als Verkäuferin gearbeitet.} ;
\item \cle{Motivation} : pourquoi ce poste m'intéresse — \all{Ich interessiere mich sehr für diese Arbeit, weil ich den Kontakt mit Kunden mag.} ;
\item \cle{Disponibilité} : quand je peux commencer — \all{Ich stehe ab sofort zur Verfügung.} ou \all{Ich kann ab dem 1. Juni arbeiten.}
\end{enumerate}
Le \all{wenn}-Satz et le futur révisés précédemment sont parfaits ici : \all{Wenn Sie mir diese Chance geben, werde ich mein Bestes geben.} (« Si vous me donnez cette chance, je donnerai le meilleur de moi-même. »).
\end{propriete}

\begin{propriete}[Le registre de langue : Sie obligatoire]
Dans tous les textes professionnels (annonce adressée au public, lettre, compte rendu officiel), on utilise le \cle{vouvoiement} \all{Sie}, toujours écrit avec une \cle{majuscule}, ainsi que ses formes dérivées : \all{Ihnen} (datif), \all{Ihre} (possessif). Le tutoiement \all{du} est réservé à la famille et aux amis.
\end{propriete}

\begin{center}
\begin{tabularx}{\linewidth}{|X|X|X|}\hline
\rowcolor{popblueL} \textbf{Français} & \textbf{Deutsch} & \textbf{Remarque} \\ \hline
Madame, Monsieur, & Sehr geehrte Damen und Herren, & virgule obligatoire \\ \hline
Monsieur Kamda, & Sehr geehrter Herr Kamda, & adjectif décliné \\ \hline
Je vous écris pour... & Ich schreibe Ihnen, weil... & \all{weil} : verbe à la fin \\ \hline
Cordialement & Mit freundlichen Grüßen & sans virgule après \\ \hline
Je reste à votre disposition & Ich stehe Ihnen zur Verfügung & datif \all{Ihnen} \\ \hline
\end{tabularx}
\end{center}

\courssub{Le compte rendu (der Bericht)}

\begin{definition}[Der Bericht]
Le \cle{compte rendu} (\all{der Bericht, -e}) est un texte qui relate des \cle{faits objectifs} (une réunion, une visite d'entreprise, un événement) sans donner son opinion personnelle. Il répond aux \all{W-Fragen} : \all{wer?} (qui), \all{was?} (quoi), \all{wann?} (quand), \all{wo?} (où), \all{warum?} (pourquoi).
\end{definition}

\begin{propriete}[Temps et ordre dans le compte rendu]
Le compte rendu relate des faits passés : on utilise le \cle{Perfekt} (style oral et courant : \all{Die Klasse hat die Fabrik besucht.}) ou le \cle{Präteritum} (style écrit, presse : \all{Die Klasse besuchte die Fabrik.}). Les faits sont présentés dans l'\cle{ordre chronologique}, reliés par des connecteurs :
\begin{itemize}
\item \all{zuerst} (d'abord), \all{dann} (puis), \all{danach} (ensuite), \all{schließlich} (finalement) ;
\item \all{außerdem} (de plus), \all{deshalb} (c'est pourquoi), \all{trotzdem} (néanmoins).
\end{itemize}
Attention : ces connecteurs occupent la première place de la phrase, donc le verbe suit immédiatement (inversion) : \all{Zuerst hat der Direktor uns begrüßt.} (« D'abord, le directeur nous a accueillis. »).
\end{propriete}

\begin{methode}[Rédiger un compte rendu en 4 étapes]
\begin{enumerate}
\item \cle{Collecter les faits} : répondre aux W-Fragen (\all{wer, was, wann, wo, warum}) et noter les réponses au brouillon.
\item \cle{Ordonner chronologiquement} : classer les faits du premier au dernier.
\item \cle{Rédiger au passé} (Perfekt ou Präteritum) en reliant les phrases avec \all{zuerst, dann, danach, schließlich}, et en restant objectif (pas de \all{ich finde...}).
\item \cle{Relire} : vérifier la place du verbe, les majuscules des noms, les terminaisons et l'orthographe.
\end{enumerate}
\end{methode}

\begin{exemplebox}[Un mini-compte rendu]
\all{Am Freitag, den 14. März, hat die Klasse 1A die Schokoladenfabrik in Douala besucht. Zuerst hat der Direktor uns begrüßt. Dann haben wir die Maschinen gesehen, und ein Ingenieur hat die Produktion erklärt. Danach haben wir Kakao und Schokolade probiert. Schließlich sind wir um 12 Uhr zur Schule zurückgekehrt. Der Besuch war sehr interessant.}

« Vendredi 14 mars, la classe de 1A a visité la chocolaterie de Douala. D'abord, le directeur nous a accueillis. Puis nous avons vu les machines, et un ingénieur a expliqué la production. Ensuite, nous avons goûté du cacao et du chocolat. Finalement, nous sommes retournés à l'école à midi. La visite était très intéressante. »
\end{exemplebox}

\begin{popfigure}
\begin{center}
\begin{tikzpicture}[
typ/.style={rectangle, draw=popdark, fill=poporangeL, rounded corners=3pt, minimum width=3.4cm, minimum height=0.9cm, align=center, font=\small\bfseries},
car/.style={rectangle, draw=popblue, fill=popblueL, rounded corners=2pt, minimum width=3.4cm, minimum height=1.5cm, align=center, font=\footnotesize},
arr/.style={-{Stealth[length=2.5mm]}, thick, popdark}]
\node[typ] (t1) {Kleinanzeige};
\node[typ, right=1.2cm of t1] (t2) {Brief};
\node[typ, right=1.2cm of t2] (t3) {Bericht};
\node[car, below=0.5cm of t1, align=center] (c1){bref, style\\télégraphique\\prix + contact};
\node[car, below=0.5cm of t2, align=center] (c2){formel, vouvoiement\\Sie, structure fixe\\MfG};
\node[car, below=0.5cm of t3, align=center] (c3){factuel, objectif\\passé, ordre\\chronologique};
\draw[arr] (t1) -- (c1);
\draw[arr] (t2) -- (c2);
\draw[arr] (t3) -- (c3);
\end{tikzpicture}
\end{center}
\legende{Trois types de textes économiques et leurs caractéristiques essentielles.}
\end{popfigure}

\begin{exoresolu}[Rédiger une petite annonce]
Énoncé : Marie veut vendre son téléphone portable. Il est presque neuf, en très bon état. Prix : 30 000 FCFA, à débattre. Elle habite à Yaoundé. Rédigez sa petite annonce en allemand (4 lignes maximum).

\tcblower

Solution détaillée : on applique la structure titre → objet + description → prix → contact, en style télégraphique et avec l'abréviation \all{VB} puisque le prix est à débattre :

\all{\textbf{Zu verkaufen: Handy}}

\all{Handy, fast neu, sehr guter Zustand, 30 000 FCFA, VB.}

\all{Abholung in Yaoundé. Tel.: 6XX XX XX XX.}

Points de vigilance : majuscule à \all{Handy} (nom), adjectif décliné \all{guter Zustand}, pas de phrase complète avec sujet et verbe conjugué dans le corps de l'annonce.
\end{exoresolu}

\begin{exercice}[À vous de produire]
\begin{enumerate}
\item Rédigez une petite annonce : vous cherchez (\all{Suche...}) un vélo d'occasion, maximum 40 000 FCFA, à Douala.
\item Complétez cette lettre avec les éléments manquants : \all{......, den 20. April 2025 — Sehr geehrte Damen und Herren, ich schreibe Ihnen, weil... — Mit ......... Grüßen}.
\item Rédigez un compte rendu (5 phrases) de la visite d'un marché à Yaoundé, en utilisant \all{zuerst, dann, danach, schließlich} et le Perfekt.
\end{enumerate}
\end{exercice}

\begin{corrige}[Exercice : éléments de réponse]
1. \all{Suche: Fahrrad. Gebrauchtes Fahrrad, guter Zustand, max. 40 000 FCFA. Douala. Tel.: 6XX XX XX XX.} — 2. \all{Yaoundé, den 20. April 2025} ; \all{Mit freundlichen Grüßen} (sans virgule). — 3. Exemple : \all{Am Samstag hat unsere Klasse den Markt Mokolo besucht. Zuerst haben wir die Obststände gesehen. Dann hat eine Händlerin uns die Preise erklärt. Danach haben wir Früchte gekauft. Schließlich sind wir um 11 Uhr nach Hause gegangen.}
\end{corrige}

\begin{aretenir}[L'essentiel de la section]
La petite annonce est brève et télégraphique (\all{VB} = prix à débattre) ; la lettre formelle suit sept blocs fixes, avec \all{Sie} obligatoire, la date \all{Yaoundé, den 15. März 2025} et \all{Mit freundlichen Grüßen} sans virgule ; le compte rendu relate des faits objectifs au passé, dans l'ordre chronologique, grâce aux connecteurs \all{zuerst, dann, danach, schließlich}. Ces trois outils nous préparent directement à la dernière section : la traduction (thème et version) de textes économiques courts.
\end{aretenir}

\coursec{Traduction : thème et version}

Après avoir étudié les types de textes caractéristiques de la vie économique — petite annonce, lettre formelle, compte rendu — nous abordons maintenant la compétence transversale par excellence : la \cle{traduction}. Passer du français à l'allemand (\cle{Thème}) ou de l'allemand au français (\cle{Version}) ne consiste pas à remplacer mot à mot, mais à \textbf{transposer le sens} en respectant les structures propres à chaque langue. Cette section vous donne la méthode rigoureuse, les pièges à éviter et un entraînement progressif sur le thème du marché et de l'économie.

\courssub{Méthodologie de la version (allemand → français)}

La version est un exercice de \textbf{compréhension fine} suivi d'une \textbf{reformulation naturelle} en français. Ne traduisez jamais mot à mot : l'ordre des mots, les cas et les temps diffèrent trop.

\begin{methode}[Traduire un texte allemand en français — 4 étapes]
\begin{enumerate}
    \item \textbf{Lire globalement :} Lisez le texte une première fois sans dictionnaire pour saisir le sujet, le ton (informatif, narratif, argumentatif) et la structure (paragraphes, connecteurs).
    \item \textbf{Analyser les structures :} Repérez le \cle{verbe} de chaque phrase (souvent en position 2 ou en fin de subordonnée). Identifiez le \cle{sujet} (Nominatif), les \cle{compléments} (Accusatif, Datif, Génitif) et les \cle{subordonnées} (\all{wenn}, \all{weil}, \all{dass}, \all{als}). Soulignez les \cle{comparatifs} (\all{-er als}, \all{weniger ... als}) et les \cle{formes du futur} (\all{werden} + infinitif).
    \item \textbf{Traduire par segments de sens :} Traitez proposition par proposition. Cherchez l'équivalent français \emph{idiomatique} (ex. \all{es gibt} $\to$ « il y a », pas « il donne » ; \all{man} $\to$ « on » ou passif). Respectez la chronologie des temps (Präteritum $\to$ passé simple/imparfait ; Perfekt $\to$ passé composé).
    \item \textbf{Relire et polir :} Relisez votre version française \emph{sans regarder l'allemand}. Est-elle fluide ? Le vocabulaire est-il précis (économique) ? L'orthographe et la ponctuation sont-elles correctes ?
\end{enumerate}
\end{methode}

\begin{exemplebox}[Version guidée : analyse d'une phrase]
\all{Wenn die Wirtschaft wächst, wird es mehr Arbeitsplätze geben.}
\begin{itemize}
    \item \textbf{Structure :} Subordonnée \all{Wenn die Wirtschaft wächst} (verbe \all{wächst} à la fin) + Principale \all{wird es mehr Arbeitsplätze geben} (verbe \all{wird} en position 2, infinitif \all{geben} à la fin).
    \item \textbf{Vocabulaire :} \all{die Wirtschaft} (l'économie), \all{wachsen} (croître), \all{der Arbeitsplatz} (l'emploi, pl. \all{die Arbeitsplätze}), \all{es gibt} (il y a).
    \item \textbf{Traduction :} « \textbf{Si l'économie croît, il y aura plus d'emplois.} » (Futur simple en français pour \all{wird ... geben}).
\end{itemize}
\end{exemplebox}

\courssub{Méthodologie du thème (français → allemand)}

Le thème est un exercice de \textbf{construction grammaticale consciente}. Vous devez « penser en allemand » : déterminer le cas de chaque nom, placer le verbe correctement, choisir le bon temps.

\begin{methode}[Traduire une phrase française en allemand — 4 temps]
\begin{enumerate}
    \item \textbf{Trouver le verbe et son sujet :} Conjuguez le verbe au temps requis (Présent, Perfekt, Präteritum, Futur I). Le sujet dicte la personne (ich, du, er/sie/es, wir, ihr, sie/Sie).
    \item \textbf{Identifier les compléments et leurs CAS :} 
    \begin{itemize}
        \item Sujet $\to$ \cle{Nominatif} (Wer/Was?).
        \item Objet direct $\to$ \cle{Accusatif} (Wen/Was?).
        \item Objet indirect / destinataire $\to$ \cle{Datif} (Wem?).
        \item Appartenance $\to$ \cle{Génitif} (Wessen?).
    \end{itemize}
    Déclinez articles, adjectifs, pronoms en conséquence.
    \item \textbf{Construire la phrase selon l'ORDRE ALLEMAND :}
    \begin{itemize}
        \item \textbf{Phrase principale :} Verbe conjugué en \textbf{2e position} (Sujet - Verbe - ... ou Complément - Verbe - Sujet - ...).
        \item \textbf{Subordonnée (si, quand, parce que, que) :} Conjonction en 1re position, verbe conjugué à la \textbf{fin}.
        \item \textbf{Règle TMP :} Temps — Manière — Lieu (avant les infinitifs/participe à la fin).
    \end{itemize}
    \item \textbf{Vérifier :} Majuscules à \textbf{TOUS les noms} ? Umlauts (ä, ö, ü) et ß corrects ? Verbe à la bonne place ? Accords sujet/verbe ? Ponctuation (virgule avant subordonnée) ?
\end{enumerate}
\end{methode}

\begin{exemplebox}[Thème guidé : deux phrases types]
\begin{enumerate}
    \item \textbf{Français :} « Si l'économie grandit, il y aura plus d'emplois. » \\
    \textbf{Analyse :} « Si » $\to$ \all{wenn} (subordonnée $\to$ verbe à la fin). « l'économie » (suj., fém.) $\to$ \all{die Wirtschaft}. « grandit » $\to$ \all{wächst} (Présent). Principale : Futur I $\to$ \all{werden} + infinitif. « il y a » $\to$ \all{es gibt} (infinitif \all{geben}). « plus d'emplois » (acc. pl.) $\to$ \all{mehr Arbeitsplätze}.
    \textbf{Allemand :} \all{Wenn die Wirtschaft wächst, wird es mehr Arbeitsplätze geben.}
    
    \item \textbf{Français :} « Ce magasin est moins cher que l'autre. » \\
    \textbf{Analyse :} « Ce magasin » (suj., neutre) $\to$ \all{Dieses Geschäft}. « est » $\to$ \all{ist}. « moins cher » $\to$ \all{billiger} (comparatif de \all{billig}) ou \all{weniger teuer}. « que » $\to$ \all{als}. « l'autre » (neutre, nom., se référant à \all{Geschäft}) $\to$ \all{das andere}.
    \textbf{Allemand :} \all{Dieses Geschäft ist billiger als das andere.} / \all{Dieses Geschäft ist weniger teuer als das andere.}
\end{enumerate}
\end{exemplebox}

\courssub{Pièges fréquents et points de vigilance}

Les francophones commettent des erreurs systématiques. Mémorisez ces \cle{points de vigilance} pour gagner des points à l'examen.

\begin{attention}[Faux ami majeur : \all{bekommen} $\neq$ « devenir »]
\all{bekommen} signifie \textbf{« recevoir / obtenir »} (quelque chose que l'on attend ou qu'on achète). 
\begin{itemize}
    \item \all{Ich bekomme ein Gehalt.} = « Je reçois un salaire. »
    \item \all{Was bekommst du zum Geburtstag?} = « Qu'est-ce que tu reçois / as-tu pour ton anniversaire ? »
\end{itemize}
Pour « devenir » (changement d'état), on utilise \cle{\all{werden}} :
\begin{itemize}
    \item \all{Er wird Arzt.} = « Il devient médecin. »
    \item \all{Die Preise werden höher.} = « Les prix deviennent plus élevés. »
\end{itemize}
\textbf{Astuce mnémotechnique :} \all{bekommen} $\to$ « be-kommen » $\to$ « venir vers soi » $\to$ \textbf{recevoir}.
\end{attention}

\begin{attention}[Majuscule à TOUS les noms — Sans exception !]
En allemand, \textbf{tout nom commun} (substantif) prend une majuscule, où qu'il se trouve dans la phrase.
\begin{itemize}
    \item \all{Ich gehe auf den \textbf{Markt}.} (Je vais au marché.)
    \item \all{Das \textbf{Geld} liegt auf dem \textbf{Tisch}.} (L'argent est sur la table.)
    \item \all{Wir suchen \textbf{Arbeit}.} (Nous cherchons du travail / un emploi.)
\end{itemize}
\textbf{Piège :} Les noms verbalisés (infinitifs substantivés) aussi : \all{das \textbf{Kaufen}}, \all{das \textbf{Verkaufen}}, \all{beim \textbf{Bezahlen}}.
\end{attention}

\begin{attention}[Ordre des mots dans la subordonnée \all{wenn} / \all{als} / \all{weil}]
La conjonction (\all{wenn}, \all{als}, \all{weil}, \all{dass}, \all{ob}) place le \textbf{verbe conjugué à la toute fin} de la proposition.
\begin{center}
\begin{tabular}{|l|l|}
\hline
\textbf{Correct} & \textbf{Incorrect (ordre français)} \\
\hline
\all{Wenn ich \textbf{gehe}, ...} & \all{Wenn \textbf{gehe} ich, ...} \\
\all{..., weil es \textbf{regnet}.} & \all{..., weil \textbf{regnet} es.} \\
\hline
\end{tabular}
\end{center}
Rappel : \all{wenn} = « si » (présent/futur) ou « quand » (répété au passé). \all{als} = « quand » (ponctuel au passé).
\end{attention}

\courssub{Tableau de correspondances grammaticales}

Ce tableau récapitule les correspondances essentielles pour la traduction des textes économiques (comparaison, condition, futur).

\begin{center}
\begin{tabularx}{\linewidth}{|X|X|X|}
\hline
\rowcolor{popblueL} 
\textbf{Français (Fonction)} & \textbf{Deutsch (Structure)} & \textbf{Remarque / Exemple} \\
\hline
« Si » (condition, présent/futur) & \all{wenn} + \textbf{Présent} / \textbf{Futur} & \all{Wenn es regnet, bleibe ich zu Hause.} (Si pleut, reste je.) \\
\hline
« Quand » (passé, une fois) & \all{als} + \textbf{Präteritum} & \all{Als ich jung \textbf{war}, ...} (Quand j'étais jeune...) \\
\hline
« Quand » (passé, répété) & \all{wenn} + \textbf{Präteritum} & \all{Wenn er \textbf{kam}, aß er.} (Chaque fois qu'il venait, il mangeait.) \\
\hline
« Plus ... que » (supériorité) & \textbf{-er} + \all{als} & \all{größer \textbf{als}}, \all{teurer \textbf{als}}, \all{besser \textbf{als}} (irrégulier) \\
\hline
« Moins ... que » (infériorité) & \all{weniger} + adjectif + \all{als} & \all{weniger teuer \textbf{als}}, \all{weniger wichtig \textbf{als}} \\
\hline
« Autant ... que » (égalité) & \all{so} + adjectif + \all{wie} & \all{so groß \textbf{wie}}, \all{so billig \textbf{wie}} \\
\hline
« Je vais + infinitif » (futur proche) & \all{werden} + \textbf{Infinitif} (Futur I) & \all{Ich \textbf{werde} kaufen.} (J'achèterai / Je vais acheter.) \\
\hline
« Il y a / Il y aura » & \all{es gibt} / \all{es wird ... geben} & \all{Es gibt viele Märkte.} / \all{Es wird mehr Märkte \textbf{geben}.} \\
\hline
« On » (sujet indéfini) & \all{man} (conjug. 3e pers. sg.) & \all{Man findet alles.} (On trouve tout / Tout se trouve.) \\
\hline
\end{tabularx}
\end{center}

\courssub{Schéma des étapes de la traduction}

\begin{popfigure}
\begin{tikzpicture}[
    node distance=1.2cm and 2cm,
    box/.style={draw=popblue, thick, rounded corners, fill=popblueL, text width=3.2cm, align=center, minimum height=1.8cm, font=\small},
    arrow/.style={-Stealth, thick, popdark}
]
\node[box, align=center] (lire){\textbf{1. LIRE}\\\emph{Globalement}\\\textit{Sujet, ton, structure}};
\node[box, right=of lire, align=center] (analyser){\textbf{2. ANALYSER}\\\emph{Verbes, cas, subordonnées}\\\textit{Comparatifs, Futur, Wenn/Als}};
\node[box, right=of analyser, align=center] (traduire){\textbf{3. TRADUIRE}\\\emph{Par segments de sens}\\\textit{Idiomatique, vocab. précis}};
\node[box, right=of traduire, align=center] (relire){\textbf{4. RELIRE}\\\emph{Sans l'original}\\\textit{Fluidité, orthographe, cas}};

\draw[arrow] (lire) -- (analyser);
\draw[arrow] (analyser) -- (traduire);
\draw[arrow] (traduire) -- (relire);

% Boucle de retour
\draw[arrow, poporange, dashed] (relire.south) -- ++(0,-0.5) -| node[below, midway, font=\tiny, poporange] {Si doute : retour étape 2} (analyser.south);
\end{tikzpicture}
\legende{Les 4 étapes incontournables de la traduction (Version et Thème). La flèche pointillée rappelle qu'on revient à l'analyse en cas d'incertitude.}
\end{popfigure}

\courssub{Entraînement guidé : Thème (français → allemand)}

Appliquez la méthode en 4 temps. Traduisez les 5 phrases suivantes. Elles réinvestissent le comparatif, \all{wenn}-Sätze et le Futur I.

\begin{exoresolu}[Thème n°1 — Phrase modèle]
\textbf{Français :} « Si les prix baissent, les clients achèteront plus de produits. » \\
\textbf{Résolution pas à pas :}
\begin{enumerate}
    \item \textbf{Verbes/Sujets :} « baissent » (suj. \emph{les prix} $\to$ pl.) = \all{sinken} (Présent, subordonnée $\to$ fin). « achèteront » (suj. \emph{les clients} $\to$ pl.) = Futur I \all{werden kaufen} (principale $\to$ \all{werden} pos. 2, \all{kaufen} fin).
    \item \textbf{Cas :} \all{die Preise} (Nom. pl.), \all{die Kunden} (Nom. pl.), \all{mehr Produkte} (Acc. pl., pas d'article indéfini au pluriel).
    \item \textbf{Ordre :} Subordonnée : \all{Wenn die Preise sinken,} / Principale : \all{werden die Kunden mehr Produkte kaufen} (inversion sujet/verbe car la subordonnée est en position 1).
    \item \textbf{Vérif :} Majuscules \all{Preise, Kunden, Produkte}. Virgule. Umlaut \all{sinken} (ok). \all{werden} pos 2. \all{kaufen} fin.
\end{enumerate}
\textbf{Allemand :} \all{Wenn die Preise sinken, werden die Kunden mehr Produkte kaufen.}
\end{exoresolu}

\begin{exercice}[Thème — À vous de jouer !]
Traduisez en allemand. Soignez l'ordre des mots, les cas et les majuscules.
\begin{enumerate}
    \item Le marché de Mokolo est plus grand que le marché de Bafoussam.
    \item Si j'ai de l'argent, j'achèterai un nouveau téléphone portable. (\textit{Indice : Handy})
    \item Ce vendeur est moins aimable que l'autre vendeur.
    \item Quand on arrive tôt, on trouve les meilleurs prix. (\textit{Utilisez \all{man} et \all{wenn}})
    \item L'année prochaine, l'économie du Cameroun va croître. (\textit{Futur I avec \all{nächstes Jahr}})
\end{enumerate}
\end{exercice}

\begin{corrige}[Exercice Thème]
\begin{enumerate}
    \item \all{Der Markt in Mokolo ist größer als der Markt in Bafoussam.} (\textit{Comparatif \all{größer als}, cas nominatif des deux côtés}).
    \item \all{Wenn ich Geld habe, werde ich ein neues Handy kaufen.} (\textit{\all{wenn} + Présent à la fin ; Futur I : \all{werde} pos 2, \all{kaufen} fin ; \all{ein neues Handy} Acc. sg. n. adjectif fort \all{neues}}).
    \item \all{Dieser Verkäufer ist weniger freundlich als der andere Verkäufer.} (\textit{Comparatif d'infériorité \all{weniger freundlich als} ; \all{der andere} décliné comme \all{dieser}}).
    \item \all{Wenn man früh kommt, findet man die besten Preise.} (\textit{\all{man} = on ; \all{früh} (adv.) ; \all{die besten Preise} Acc. pl. superlatif \all{beste} + faible \all{-n} car article défini \all{die}}).
    \item \all{Nächstes Jahr wird die Wirtschaft Kameruns wachsen.} (\textit{Topicalisation \all{Nächstes Jahr} $\to$ \all{werde} pos 2 ; \all{die Wirtschaft Kameruns} (Génitif \all{Kameruns} ou \all{des Kameruns}) sujet ; \all{wachsen} infinitif fin)}.
\end{enumerate}
\end{corrige}

\courssub{Entraînement guidé : Version (allemand → français)}

Voici un texte original d'environ 180 mots sur le marché de Mokolo à Yaoundé. Il contient du vocabulaire économique, des comparatifs, des \all{wenn}-Sätze et du Futur I.

\begin{textebox}[Texte original : Der Markt von Mokolo — Wirtschaftliches Herz Yaoundés]
Der Mokolo-Markt in Yaoundé ist einer der größten Märkte in Zentralafrika. Er ist viel größer als der Markt in Bafoussam und bietet eine riesige Auswahl an Waren. Man findet dort alles: frisches Obst, Gemüse, Kleidung, Elektronik und traditionelle Stoffe. Die Preise sind oft günstiger als in den Supermärkten, aber man muss gut verhandeln können.

Wenn man früh morgens kommt, findet man die frischesten Produkte und die besten Preise. Wenn die Sonne stark scheint, wird es sehr heiß und voll auf den Gängen. Viele Händler kommen aus den umliegenden Dörfern, um ihre Ernte zu verkaufen. Auch Touristen kaufen gerne Souvenirs dort.

Nächstes Jahr wird die Stadtverwaltung den Markt modernisieren. Es wird dann mehr überdachte Stände und bessere Toiletten geben. Wenn die Bauarbeiten fertig sind, wird der Markt noch attraktiver für Kunden und Händler sein. Die Wirtschaft in Yaoundé profitiert stark von diesem lebendigen Handelsplatz.
\end{textebox}

\begin{definition}[Glossaire — Mots difficiles du texte]
\begin{itemize}
    \item \all{der Mokolo-Markt} (nom propre) : le marché Mokolo
    \item \all{Zentralafrika} : l'Afrique centrale
    \item \all{riesig} : immense, énorme
    \item \all{die Auswahl} (die Auswahlen) : le choix, la sélection
    \item \all{die Ware} (die Waren) : la marchandise, l'article
    \item \all{günstig} : avantageux, bon marché
    \item \all{verhandeln} (verhandelte, hat verhandelt) : marchander, négocier
    \item \all{der Händler} (die Händler) : le commerçant, le marchand
    \item \all{das umliegende Dorf} (die umliegenden Dörfer) : le village environnant
    \item \all{die Ernte} (die Ernten) : la récolte
    \item \all{der Souvenir} (die Souvenirs) : le souvenir (objet)
    \item \all{die Stadtverwaltung} : l'administration municipale, la mairie
    \item \all{modernisieren} : moderniser
    \item \all{überdacht} : couvert, sous abri
    \item \all{der Stand} (die Stände) : l'étal, le stand
    \item \all{die Toilette} (die Toiletten) : les toilettes
    \item \all{die Bauarbeit} (die Bauarbeiten) : les travaux (de construction)
    \item \all{fertig} : fini, terminé
    \item \all{attraktiv} : attractif
    \item \all{profitieren von} (+ Datif) : profiter de, bénéficier de
    \item \all{der Handelsplatz} (die Handelsplätze) : la place de commerce, le marché
\end{itemize}
\end{definition}

\begin{exoresolu}[Version guidée — Premiers paragraphes]
\textbf{Allemand :} \all{Der Mokolo-Markt in Yaoundé ist einer der größten Märkte in Zentralafrika. Er ist viel größer als der Markt in Bafoussam und bietet eine riesige Auswahl an Waren.}
\begin{itemize}
    \item \all{Der Mokolo-Markt ... ist einer der größten Märkte} : « Le marché Mokolo ... est l'un des plus grands marchés ». \all{einer der größten} = superlatif avec article indéfini (l'un des...). \all{in Zentralafrika} = en Afrique centrale.
    \item \all{Er ist viel größer als der Markt in Bafoussam} : « Il est bien plus grand que le marché de Bafoussam. » \all{viel größer als} = comparatif renforcé par \all{viel}.
    \item \all{und bietet eine riesige Auswahl an Waren} : « et offre un choix immense de marchandises. » \all{bietet} (verbe \all{bieten}, 3e pers. sg.) ; \all{eine riesige Auswahl} (Acc. fém. sg.) ; \all{an Waren} (préposition \all{an} + Datif pl. ici figé).
\end{itemize}
\textbf{Français :} « Le marché Mokolo à Yaoundé est l'un des plus grands marchés d'Afrique centrale. Il est bien plus grand que le marché de Bafoussam et offre un choix immense de marchandises. »
\end{exoresolu}

\begin{exercice}[Version — Traduisez le texte complet en français]
Produisez une traduction française \textbf{fluide, naturelle et correcte} (niveau langue soutenue / journalistique). Utilisez la grille d'auto-correction ci-dessous pour vous relire.
\end{exercice}

\begin{corrige}[Exercice Version — Proposition de traduction]
« Le marché Mokolo à Yaoundé est l'un des plus grands marchés d'Afrique centrale. Il est bien plus grand que le marché de Bafoussam et offre un choix immense de marchandises. On y trouve de tout : fruits frais, légumes, vêtements, électronique et tissus traditionnels. Les prix y sont souvent plus avantageux que dans les supermarchés, mais il faut savoir bien marchander.

Quand on arrive tôt le matin, on trouve les produits les plus frais et les meilleurs prix. Quand le soleil tape fort, il fait très chaud et c'est très encombré dans les allées. Beaucoup de commerçants viennent des villages environnants pour vendre leur récolte. Les touristes aussi aiment y acheter des souvenirs.

L'année prochaine, la mairie va moderniser le marché. Il y aura alors plus d'étals couverts et de meilleures toilettes. Quand les travaux seront terminés, le marché sera encore plus attractif pour les clients et les commerçants. L'économie de Yaoundé profite fortement de cette place de commerce vivante. »

\textbf{Notes de traduction :}
\begin{itemize}
    \item \all{Man findet dort alles} $\to$ « On y trouve de tout » (meilleur que « On trouve tout là-bas »).
    \item \all{günstiger als in den Supermärkten} $\to$ « plus avantageux que dans les supermarchés ».
    \item \all{Wenn die Sonne stark scheint} $\to$ « Quand le soleil tape fort / brille fort » (adaptation idiomatique).
    \item \all{voll auf den Gängen} $\to$ « encombré dans les allées / bondé dans les allées ».
    \item \all{Nächstes Jahr wird die Stadtverwaltung ... modernisieren} $\to$ Futur I $\to$ futur simple français « va moderniser / modernisera ».
    \item \all{Es wird dann mehr überdachte Stände ... geben} $\to$ « Il y aura alors plus d'étals couverts... » (\all{es gibt} $\to$ il y a).
    \item \all{profitiert stark von} + Datif \all{diesem ... Handelsplatz} $\to$ « profite fortement de cette... ».
\end{itemize}
\end{corrige}

\courssub{Grille d'auto-correction : Relisez-vous systématiquement !}

Avant de rendre votre copie (thème ou version), cochez chaque case. Une erreur non corrigée = des points perdus.

\begin{aretenir}[Check-list de relecture — 6 points de contrôle]
\begin{enumerate}
    \item \textbf{Ordre des mots (Syntax) :}
    \begin{itemize}
        \item Phrase principale : Verbe conjugué en \textbf{2e position} ? (Sujet-Verbe... OU Complément-Verbe-Sujet...)
        \item Subordonnée (\all{wenn, weil, dass, als, ob}) : Verbe conjugué à la \textbf{fin} ?
        \item Règle TMP (Temps - Manière - Lieu) respectée avant la fin de phrase ?
    \end{itemize}
    \item \textbf{Cas et Déclinaisons (Kasus) :}
    \begin{itemize}
        \item Sujet = Nominatif ? Objet direct = Accusatif ? Objet indirect = Datif ?
        \item Articles (\all{der, die, das, ein, kein, dieser, jener...}) correctement déclinés ?
        \item Adjectifs : déclinaison \textbf{forte} (sans article), \textbf{faible} (article défini), \textbf{mixtes} (article indéfini) ?
        \item Prépositions : \all{durch, für, gegen, ohne, um} + Acc. ? \all{aus, bei, mit, nach, seit, von, zu} + Dat. ? \all{an, auf, hinter, in, neben, über, unter, vor, zwischen} + Acc. (mouvement) / Dat. (position) ?
    \end{itemize}
    \item \textbf{Accords Sujet - Verbe :} Verbe au pluriel si sujet pluriel (\all{die Preise sind}, pas \all{ist}) ?
    \item \textbf{Orthographe allemande (Rechtschreibung) :}
    \begin{itemize}
        \item \textbf{Majuscule à TOUS les noms} (\all{der Markt, die Ware, das Geld, beim Kauf}) ?
        \item Umlauts \all{ä, ö, ü} et \all{ß} (après voyelle longue/diphtongue) corrects ?
        \item Verbes à particule séparable : particule à la fin (\all{Ich kaufe ein}, \all{Er kommt an}) ?
    \end{itemize}
    \item \textbf{Vocabulaire précis (Faux amis) :}
    \begin{itemize}
        \item \all{bekommen} = recevoir (pas devenir) ? \all{werden} = devenir ?
        \item \all{aktuell} = actuel (pas actuel = \all{tatsächlich}) ?
        \item \all{konsequent} = cohérent/conséquent (pas « conséquent » au sens « important ») ?
        \item \all{das Rezept} = l'ordonnance/la recette (pas « le reçu » = \all{der Kassenzettel / die Quittung}) ?
    \end{itemize}
    \item \textbf{Ponctuation et Style (Version) :}
    \begin{itemize}
        \item Virgule \textbf{obligatoire} avant toute subordonnée en allemand (reportée en français si nécessaire).
        \item Français naturel : pas de calque (\all{es gibt} $\to$ « il y a », pas « il donne » ; \all{man} $\to$ « on / passif »).
        \item Temps des verbes français cohérents (Présent/Imparfait/Passé composé/Futur).
    \end{itemize}
\end{enumerate}
\end{aretenir}

\begin{savaistu}[Le marché Mokolo : poumon économique de Yaoundé]
Le marché \all{Mokolo} (quartier \all{Mokolo}, arrondissement de Yaoundé II) est effectivement le plus grand marché de détail du Cameroun. Il s'étend sur plusieurs hectares et emploie directement des milliers de commerçants (\all{Händler}), sans compter les emplois induits (porteurs \all{Träger}, transporteurs, restaurateurs). On y pratique le \cle{bilingue économique} : français, anglais, mais aussi langues locales (Bassa, Ewondo, Bamiléké) et parfois l'allemand pour les produits importés ou les touristes. La modernisation annoncée (couverture des étals, assainissement) est un enjeu majeur pour la \cle{ville de Yaoundé} (\all{Yaoundé-Stadt}), soutenue par la coopération décentralisée avec des villes allemandes (ex. jumelage Yaoundé - \all{Düsseldorf} depuis 1973). Traduire ce réel, c'est aussi traduire cette vitalité !
\end{savaistu}

\newpage

\coursec{Activité d'intégration}

\coursec{Leçon ALA09 : Production et traduction de textes liés à la vie économique}

\courssub{Texte support et compréhension}

\begin{textebox}[Ankündigung : Jobbörse Yaoundé]
„Herzlich willkommen zur Jobbörse Yaoundé! Am Samstag, den 15. Juni, treffen sich deutsche und kamerunische Firmen im Palais des Congrès. Suchen Sie einen Job? Dann bringen Sie Ihren Lebenslauf mit! Viele Firmen suchen junge Mitarbeiter: Verkäufer, Fremdenführer, Büroangestellte und Handwerker. Der Eintritt ist frei. Wenn Sie gut Deutsch sprechen, haben Sie bessere Chancen! Kommen Sie zahlreich — wir freuen uns auf Sie!“
\end{textebox}

\textbf{Glossaire :} die \cle{Jobbörse}, -n (salon de l'emploi) ; die \cle{Firma}, die \cle{Firmen} (entreprise) ; der \cle{Mitarbeiter}, - (employé) ; der \cle{Verkäufer}, - (vendeur) ; der \cle{Fremdenführer}, - (guide touristique) ; der \cle{Büroangestellte}, -n (employé de bureau) ; der \cle{Handwerker}, - (artisan) ; der \cle{Eintritt} (entrée) ; die \cle{Chance}, -n (chance, opportunité) ; \cle{zahlreich} (nombreux) ; sich \cle{freuen auf} + Akkusativ (se réjouir de).

\begin{exercice}[Compréhension du document]
Répondez en français :
\begin{enumerate}
\item Quand et où a lieu le salon ?
\item Quels métiers sont proposés ?
\item Quel avantage ont les candidats qui parlent allemand ?
\item Soulignez dans le texte le comparatif et la phrase introduite par \all{wenn}.
\end{enumerate}
\end{exercice}

\begin{corrige}[Exercice 1]
\begin{enumerate}
\item Le salon a lieu le samedi 15 juin au Palais des Congrès de Yaoundé.
\item Les métiers proposés sont : vendeur, guide touristique, employé de bureau, artisan.
\item Les candidats qui parlent allemand ont de meilleures chances (\all{bessere Chancen}).
\item Comparatif : \all{bessere Chancen} (meilleures chances) ; phrase avec \all{wenn} : \all{Wenn Sie gut Deutsch sprechen, haben Sie bessere Chancen!}
\end{enumerate}
\end{corrige}

\courssub{Lexique thématique : la vie économique}

\begin{aretenir}[Vocabulaire de base : entreprise, emploi, candidature]
\begin{center}
\begin{tabularx}{\linewidth}{|X|X|X|}
\hline
\rowcolor{popblueL} \textbf{Français} & \textbf{Deutsch} & \textbf{Remarque} \\
\hline
le métier & der \cle{Beruf}, -e &  \\
le poste, la place & die \cle{Stelle}, -n &  \\
l'employeur / l'employée & der \cle{Arbeitgeber}, - / die \cle{Arbeitnehmerin}, -nen &  \\
le salaire (ouvrier) & der \cle{Lohn}, die \cle{Löhne} &  \\
le salaire (employé) & das \cle{Gehalt}, die \cle{Gehälter} &  \\
embaucher & \cle{einstellen} & tr., Partizip II: \all{eingestellt} \\
postuler à & sich \cle{bewerben um} + Akk. & réfléchi + préposition \all{um} \\
la candidature & die \cle{Bewerbung}, -en &  \\
le CV & der \cle{Lebenslauf}, die \cle{Lebensläufe} &  \\
l'expérience & die \cle{Erfahrung}, -en &  \\
ponctuel & \cle{pünktlich} &  \\
fiable & \cle{zuverlässig} &  \\
vendre / acheter & \cle{verkaufen} / \cle{kaufen} &  \\
bon marché & \cle{billig} / \cle{preiswert} & \all{preiswert} = bon rapport qualité/prix \\
la qualité & die \cle{Qualität}, -en &  \\
la coopérative & die \cle{Genossenschaft}, -en &  \\
l'artisanat & das \cle{Handwerk} &  \\
l'agroalimentaire & die \cle{Lebensmittelindustrie} &  \\
\hline
\end{tabularx}
\end{center}
\end{aretenir}

\begin{savaistu}[Contexte camerounais]
Le Cameroun est le 5e producteur mondial de cacao. Les coopératives de la région de l'Est (Bertoua, Abong-Mbang) exportent vers l'Allemagne (transformateurs de chocolat à Cologne, Hambourg). Le port de Kribi (nouveau port en eau profonde) facilite l'export. L'artisanat de Foumban (bronzes, tissus) attire le tourisme germanophone. L'allemand est un atout majeur pour ces filières.
\end{savaistu}

\courssub{Grammaire 1 : Révision du comparatif et du superlatif}

\begin{propriete}[Formation du comparatif et du superlatif]
\begin{itemize}
\item \textbf{Comparatif d'inégalité :} adjectif/adverbe + \all{-er} + \all{als} (que).
\item \textbf{Superlatif (prédicatif) :} \all{am} + adjectif + \all{-sten} (après verbe d'état).
\item \textbf{Superlatif (attributif) :} article défini + adjectif + \all{-ste} + terminaison de déclinaison.
\end{itemize}
\end{propriete}

\begin{center}
\begin{tabular}{|l|l|l|l|}
\hline
\rowcolor{popblueL} \textbf{Positif} & \textbf{Comparatif} & \textbf{Superlatif (préd.)} & \textbf{Superlatif (attr.)} \\
\hline
\all{gut} & \all{besser als} & \all{am besten} & \all{der/die/das beste} \\
\all{viel} & \all{mehr als} & \all{am meisten} & \all{der/die/das meiste} \\
\all{gern} & \all{lieber als} & \all{am liebsten} & \all{der/die/das liebste} \\
\all{hoch} & \all{höher als} & \all{am höchsten} & \all{der/die/das höchste} \\
\all{teuer} & \all{teurer als} & \all{am teuersten} & \all{der/die/das teuerste} \\
\all{billig} & \all{billiger als} & \all{am billigsten} & \all{der/die/das billigste} \\
\all{alt} & \all{älter als} & \all{am ältesten} & \all{der/die/das älteste} \\
\hline
\end{tabular}
\end{center}

\begin{attention}[Pièges francophones]
\begin{itemize}
\item Ne pas ajouter \all{mehr} devant le comparatif : \textbf{faux} \all{mehr besser}, \textbf{correct} \all{besser}.
\item Umlaut fréquent au comparatif/superlatif pour les monosyllabes : \all{alt} $\rightarrow$ \all{älter}, \all{am ältesten} ; \all{jung} $\rightarrow$ \all{jünger}, \all{am jüngsten} ; \all{kalt} $\rightarrow$ \all{kälter}, \all{am kältesten} ; \all{groß} $\rightarrow$ \all{größer}, \all{am größten} ; \all{kurz} $\rightarrow$ \all{kürzer}, \all{am kürzesten}.
\item Après \all{als}, le cas est le même qu'avant : \all{Er ist größer als ich (Nominativ)} / \all{Ich kenne ihn besser als sie (Akkusativ)}.
\end{itemize}
\end{attention}

\begin{exemplebox}[Exemples contextuels]
\all{Unser Kakao ist \textbf{besser} und \textbf{preiswerter} als der Kakao im Supermarkt.} (Notre cacao est meilleur et plus avantageux que celui du supermarché.) \\
\all{Die Stelle als Fremdenführer ist \textbf{interessanter} als die im Büro.} (Le poste de guide est plus intéressant que celui de bureau.) \\
\all{Am Samstag war der Andrang \textbf{am größten}.} (Samedi, l'affluence était la plus grande.) \\
\all{Wir suchen den \textbf{besten} Verkäufer.} (Nous cherchons le meilleur vendeur.)
\end{exemplebox}

\courssub{Grammaire 2 : Les phrases avec \all{wenn} et le futur}

\begin{propriete}[La subordonnée avec \all{wenn} (si / quand)]
\begin{itemize}
\item \all{wenn} introduit une subordonnée : le verbe conjugué se place \textbf{à la fin}.
\item La principale suit (ou précède) : son verbe conjugué est en \textbf{deuxième position}.
\item Structure : \all{Wenn} + sujet + ... + \textbf{verbe}, \textbf{verbe} + sujet + ...
\end{itemize}
\end{propriete}

\begin{propriete}[Le futur (Futur I) avec \all{werden}]
\begin{itemize}
\item Auxiliaire \all{werden} (conjugué) + infinitif du verbe lexical \textbf{à la fin} de la phrase.
\item Conjugaison de \all{werden} au Présent : \all{ich werde}, \all{du wirst}, \all{er/sie/es wird}, \all{wir werden}, \all{ihr werdet}, \all{sie/Sie werden}.
\end{itemize}
\end{propriete}

\begin{center}
\begin{tabular}{|l|l|l|}
\hline
\rowcolor{popblueL} \textbf{Personne} & \textbf{werden} & \textbf{Exemple : \all{arbeiten}} \\
\hline
ich & werde & \all{ich werde arbeiten} \\
du & wirst & \all{du wirst arbeiten} \\
er/sie/es & wird & \all{er wird arbeiten} \\
wir & werden & \all{wir werden arbeiten} \\
ihr & werdet & \all{ihr werdet arbeiten} \\
sie/Sie & werden & \all{sie werden arbeiten} \\
\hline
\end{tabular}
\end{center}

\begin{exemplebox}[Combinaison \all{wenn} + Futur]
\all{Wenn Sie mich einstellen, \textbf{werde} ich jeden Tag pünktlich \textbf{kommen}.} \\
\all{\textbf{Werde} ich \textbf{kommen}, wenn Sie mich einstellen?} (Question) \\
\all{Ich \textbf{werde} hart \textbf{arbeiten}, wenn ich die Stelle \textbf{bekomme}.} (Principale d'abord)
\end{exemplebox}

\begin{attention}[Erreurs fréquentes]
\begin{itemize}
\item \textbf{Ordre des mots après \all{wenn}} : \textbf{faux} \all{Wenn Sie einstellen mich,} \textbf{correct} \all{Wenn Sie mich einstellen,}
\item \textbf{Place de l'infinitif au futur} : \textbf{faux} \all{Ich werde kommen pünktlich,} \textbf{correct} \all{Ich werde pünktlich kommen.}
\item \textbf{Verbe à particule séparable au futur} : \all{anrufen} $\rightarrow$ \all{ich werde anrufen} (particule collée à l'infinitif).
\item \textbf{\all{wenn} vs \all{falls}} : \all{falls} = au cas où (hypothèse moins certaine), même structure.
\end{itemize}
\end{attention}

\courssub{Types de textes : Petite annonce, lettre, compte rendu}

\begin{methode}[Rédiger une petite annonce d'emploi / de vente]
\begin{enumerate}
\item \textbf{Titre} court, accrocheur, souvent sans verbe : \all{Verkäufer gesucht!} / \all{Top-Kakao aus Bertoua!}
\item \textbf{Présentation de l'entreprise} : nom, lieu, secteur (\all{Die Genossenschaft Cacao-Bertoua in Yaoundé...}).
\item \textbf{Offre} : poste ou produit, avec \textbf{comparatif} pour valoriser (\all{besser als}, \all{günstiger als}).
\item \textbf{Profil recherché / Qualités} : âge, langues, expérience, adjectifs (\all{pünktlich, zuverlässig, teamfähig}).
\item \textbf{Contact} : téléphone, email, adresse ; impératif poli (\all{Rufen Sie uns an!}, \all{Schreiben Sie uns!}).
\item \textbf{Formule de conclusion} : \all{Wir freuen uns auf Sie!}
\end{enumerate}
\end{methode}

\begin{methode}[Rédiger une lettre de motivation formelle]
\begin{enumerate}
\item \textbf{En-tête} (aligné à gauche) : Expéditeur (nom, rue, code postal, ville), Destinataire (entreprise, rue, ville), Lieu et date (\all{Yaoundé, den 15. Juni 2025}).
\item \textbf{Objet} (optionnel mais recommandé) : \all{Bewerbung als Verkäuferin / Ihre Anzeige vom 15. Juni}.
\item \textbf{Formule d'appel} : \all{Sehr geehrte Damen und Herren,} (si destinataire inconnu) ou \all{Sehr geehrte Frau Müller, / Sehr geehrter Herr Müller,}
\item \textbf{Corps de la lettre} (3 paragraphes types) :
  \begin{itemize}
  \item Présentation + source de l'annonce + poste visé.
  \item Compétences, expériences, qualités (\all{Ich habe Erfahrung im Verkauf...}).
  \item \textbf{Phrase clé avec \all{wenn} + Futur} : \all{Wenn Sie mir die Stelle geben, werde ich mein Bestes geben.}
  \end{itemize}
\item \textbf{Formule de politesse} : \all{Mit freundlichen Grüßen}
\item \textbf{Signature} (manuscrite) + Nom tapé + \all{Anlagen} (CV, certificats).
\end{enumerate}
\end{methode}

\begin{methode}[Rédiger un compte rendu (Bericht) d'événement]
\begin{itemize}
\item \textbf{Temps} : \cle{Perfekt} (oral, récit d'événement passé) ou \cle{Präteritum} (écrit formel, journalistique). Mélange possible.
\item \textbf{Structure} : Introduction (quand, où, quoi, qui) $\rightarrow$ Déroulement (actions principales) $\rightarrow$ Résultat / Bilan.
\item \textbf{Sujets typiques} : \all{Die Firmen, die Besucher, die Schüler, die Organisatoren}.
\item \textbf{Verbes fréquents} : \all{stattfinden} (haben), \all{kommen} (sein), \all{präsentieren} (haben), \all{schreiben} (haben), \all{probieren} (haben), \all{sein} (sein).
\end{itemize}
\end{methode}

\courssub{Traduction : Thème (FR $\rightarrow$ DE) et Version (DE $\rightarrow$ FR)}

\begin{methode}[Stratégies de traduction]
\begin{enumerate}
\item \textbf{Analyse} : identifier le type de texte, le temps, le registre, les marques de relation (connecteurs).
\item \textbf{Transposition grammaticale} :
  \begin{itemize}
  \item FR \all{si} + indicatif $\rightarrow$ DE \all{wenn} + indicatif (réel) / \all{Konjunktiv II} (irréel).
  \item FR futur simple $\rightarrow$ DE \all{werden} + infinitif.
  \item FR passé composé $\rightarrow$ DE Perfekt (oral) ou Präteritum (écrit).
  \item FR \textbf{comparatif} \emph{plus... que} $\rightarrow$ DE \all{...-er als}.
  \end{itemize}
\item \textbf{Recherche lexicale} : attention aux faux amis (\all{die Rente} $\neq$ rendement, \all{der Chef} $\neq$ cuisinier), aux collocations (\all{sich bewerben um}, \all{Erfahrung haben}).
\item \textbf{Rédaction} : produire un texte cible fluide, respecter la syntaxe allemande (place du verbe, majuscules des noms).
\item \textbf{Relecture} : grille d'auto-correction (voir fin de leçon).
\end{enumerate}
\end{methode}

\begin{exercice}[Thème : phrases isolées]
Traduisez en allemand :
\begin{enumerate}
\item Notre coopérative est plus grande que l'autre.
\item Si vous parlez allemand, vous aurez de meilleures chances.
\item Je postule pour le poste de guide touristique.
\item Nous chercherons un vendeur fiable.
\item Le cacao de Bertoua est le meilleur.
\end{enumerate}
\end{exercice}


\begin{corrige}[Exercice 2]
\begin{enumerate}
\item \all{Unsere Genossenschaft ist größer als die andere.}
\item \all{Wenn Sie Deutsch sprechen, werden Sie bessere Chancen haben.} (ou \all{haben Sie bessere Chancen} pour sens général)
\item \all{Ich bewerbe mich um die Stelle als Fremdenführer.}
\item \all{Wir werden einen zuverlässigen Verkäufer suchen.}
\item \all{Der Kakao aus Bertoua ist der beste.} (superlatif attributif)
\end{enumerate}
\end{corrige}

\begin{exercice}[Version : court texte]
Traduisez en français :
\all{„Gestern hat die Jobbörse in Douala stattgefunden. Viele Firmen sind gekommen. Die Schüler haben ihre Anzeigen präsentiert. Ein Reporter hat Fotos gemacht. Alle Teilnehmer waren sehr zufrieden.“}
\end{exercice}


\begin{corrige}[Exercice 3]
« Hier, le salon de l'emploi a eu lieu à Douala. Beaucoup d'entreprises sont venues. Les élèves ont présenté leurs annonces. Un reporter a pris des photos. Tous les participants étaient très satisfaits. »
\end{corrige}

\courssub{Activité intégrée : Simulation « Jobbörse Yaoundé »}

\courssub{Objectifs de l'activité}
À la fin de cette activité, tu seras capable de :
\begin{itemize}
\item rédiger une \cle{petite annonce} d'emploi ou de vente en allemand, claire et correcte, en utilisant au moins un \cle{comparatif} ;
\item écrire une courte \cle{lettre de motivation} en allemand contenant une phrase conditionnelle avec \all{wenn} et un verbe au \cle{futur} ;
\item traduire en français un \cle{compte rendu} rédigé en allemand par un autre groupe (version) ;
\item relire et corriger ta production à l'aide d'une grille d'auto-correction (majuscules des noms, place du verbe, cas, orthographe) ;
\item réinvestir le \cle{lexique de la vie économique} du chapitre dans une situation de communication authentique.
\end{itemize}

\courssub{Situation}
La ville de Yaoundé organise, en partenariat avec l'ambassade d'Allemagne, un salon de l'emploi germano-camerounais : la \emph{Jobbörse Yaoundé}. Des entreprises allemandes et camerounaises (agroalimentaire, tourisme, artisanat, commerce) y cherchent de jeunes collaborateurs et y présentent leurs produits. Ta classe est invitée à participer : chaque groupe joue le rôle d'une entreprise, d'un candidat et d'un journaliste.

\courssub{Consignes}
Travaillez par groupes de 3 ou 4 élèves. Suivez les étapes dans l'ordre.

\begin{enumerate}
\item \textbf{Compréhension du document.} (Déjà faite en section 1).
\item \textbf{Production 1 : la petite annonce.} Votre groupe est une entreprise (ex. : coopérative de cacao de Bertoua, hôtel de Kribi, atelier d'artisanat de Foumban, magasin d'informatique de Douala). Rédigez en allemand une petite annonce d'emploi ou de vente de 6 à 8 lignes. Elle doit contenir : un titre, la présentation de l'entreprise, le poste ou le produit, les qualités demandées, le contact — et \textbf{au moins un comparatif}.
\item \textbf{Production 2 : la lettre de motivation.} Échangez votre annonce avec un autre groupe. Chaque élève écrit individuellement une courte lettre de motivation en allemand (8 à 10 lignes) pour répondre à l'annonce reçue. La lettre doit contenir : l'en-tête (expéditeur, destinataire, lieu et date), la formule d'appel \all{Sehr geehrte Damen und Herren,}, un paragraphe de présentation, \textbf{une phrase avec \all{wenn} et un verbe au futur}, et la formule de politesse \all{Mit freundlichen Grüßen}.
\item \textbf{Traduction (version).} Un troisième groupe joue le rôle du journaliste : il rédige en allemand un court compte rendu du salon (5 à 6 phrases au Perfekt ou au Präteritum : qui était présent, ce qui s'est passé, les résultats). Votre groupe reçoit ce compte rendu et le \textbf{traduit en français}, en respectant le sens et en produisant un français correct.
\item \textbf{Relecture et correction.} Avant l'affichage, relisez chaque production avec la grille d'auto-correction ci-dessous.
\item \textbf{Le mur des annonces.} Affichez toutes les productions. Chaque groupe présente oralement son annonce en deux phrases, puis la classe vote pour l'annonce la plus convaincante.
\end{enumerate}

\begin{aretenir}[Grille d'auto-correction (à garder dans le cahier)]
\begin{center}
\begin{tabularx}{\linewidth}{|X|X|}
\hline
\rowcolor{popblueL} \textbf{Point de vigilance} & \textbf{Question à se poser / Règle} \\
\hline
Majuscules & Tous les noms (substantifs) commencent par une majuscule ? (\all{die Stelle}, \all{der Beruf}) \\
\hline
Place du verbe (phrase principale) & Le verbe conjugué est-il en \textbf{deuxième position} ? (\all{Ich \textbf{suche} einen Job.}) \\
\hline
Place du verbe (\all{wenn}-Satz) & Le verbe conjugué est-il \textbf{à la fin} de la subordonnée ? (\all{Wenn ich \textbf{gehe}...}) \\
\hline
Place de l'infinitif (Futur, modaux) & L'infinitif est-il \textbf{en toute fin} de phrase ? (\all{Ich werde kommen.}) \\
\hline
Cas (prépositions) & Akk. après \all{um, für, gegen, ohne, bis} ? Dat. après \all{aus, bei, mit, nach, seit, von, zu} ? \\
\hline
Particules séparables & La particule est-elle à la fin ? (\all{anrufen} $\rightarrow$ \all{rufen Sie \textbf{an}}) \\
\hline
Orthographe & Umlaute (\all{ä, ö, ü}), \all{ß} (après voyelle longue/diphtongue), doublement consonnes ? \\
\hline
Comparatif & Forme en \all{-er} + \all{als} ? Irréguliers (\all{besser, mehr, lieber}) respectés ? \\
\hline
\end{tabularx}
\end{center}
\end{aretenir}

\courssub{Propositions de productions modèles et corrigé commenté}

\begin{exoresolu}[Productions modèles : annonce, lettre et traduction]
\textbf{1. Petite annonce (groupe « Coopérative Cacao-Bertoua ») :}

\begin{textebox}[Modell-Anzeige]
„Verkäufer gesucht! Die Genossenschaft Cacao-Bertoua sucht einen jungen Verkäufer für unseren Laden in Yaoundé. Unser Kakao ist besser und preiswerter als der Kakao im Supermarkt. Sie sind älter als 18 Jahre und sprechen Deutsch und Französisch? Dann rufen Sie uns an: 6 90 00 00 00. Wir freuen uns auf Sie!“
\end{textebox}

\textbf{2. Lettre de motivation (réponse à cette annonce) :}

\begin{textebox}[Modell-Bewerbung]
Amina Mbarga \\
Quartier Nlongkak \\
Yaoundé, den 16. Juni 2025

Genossenschaft Cacao-Bertoua \\
Yaoundé

Betreff: Bewerbung als Verkäuferin / Ihre Anzeige

Sehr geehrte Damen und Herren,

ich heiße Amina Mbarga und ich bin 19 Jahre alt. Ich habe Ihre Anzeige gelesen und ich bewerbe mich um die Stelle als Verkäuferin. Ich spreche gut Deutsch und Französisch, und ich habe schon Erfahrung im Verkauf. Ich bin pünktlich und zuverlässig. Wenn Sie mich einstellen, werde ich jeden Tag motiviert arbeiten und ich werde viele Kunden finden.

Mit freundlichen Grüßen

Amina Mbarga
\end{textebox}

\textbf{3. Compte rendu du journaliste (allemand) et traduction (français) :}

\begin{textebox}[Bericht des Journalisten]
„Am Samstag hat die Jobbörse in Yaoundé stattgefunden. Viele deutsche und kamerunische Firmen sind gekommen. Die Schüler haben ihre Anzeigen präsentiert und Briefe geschrieben. Viele Besucher haben die Produkte probiert. Alle waren sehr zufrieden.“
\end{textebox}

\textbf{Traduction française :} « Samedi, le salon de l'emploi a eu lieu à Yaoundé. Beaucoup d'entreprises allemandes et camerounaises sont venues. Les élèves ont présenté leurs annonces et ont écrit des lettres. De nombreux visiteurs ont goûté les produits. Tout le monde était très satisfait. »
\tcblower
\textbf{Commentaire des choix de langue :}
\begin{itemize}
\item \textbf{Annonce :} le comparatif \all{besser und preiswerter als} respecte la consigne ; \all{älter als 18 Jahre} montre un second comparatif avec Umlaut (\all{alt} $\rightarrow$ \all{älter}). L'impératif poli \all{rufen Sie uns an} (verbe à particule séparable : \all{an} en fin) est typique du style publicitaire. Le titre elliptique \all{Verkäufer gesucht!} (participe II comme adjectif) est standard.
\item \textbf{Lettre :} l'en-tête respecte la norme DIN 5008 (expéditeur, destinataire, lieu + \all{den} + date accusative). La phrase clé \all{Wenn Sie mich einstellen, werde ich jeden Tag motiviert arbeiten} combine correctement la subordonnée \all{wenn} (verbe \all{einstellen} en fin) et le futur dans la principale (\all{werde} en position 2, infinitif \all{arbeiten} en fin). La construction \all{sich bewerben um} + Akk. évite le faux ami \all{für}.
\item \textbf{Compte rendu :} mélange correct de Perfekt avec \all{haben} (\all{hat stattgefunden}, \all{haben präsentiert}, \all{haben probiert}) et avec \all{sein} pour le verbe de mouvement (\all{sind gekommen}). \all{Alle waren zufrieden} utilise le Präteritum de \all{sein} (style écrit). En traduction, \all{stattgefunden hat} $\rightarrow$ « a eu lieu » ; \all{alle waren zufrieden} $\rightarrow$ « tout le monde était satisfait » (plus naturel que « tous étaient »).
\end{itemize}
\end{exoresolu}

\courssub{Bilan de l'activité}

Cette simulation de salon de l'emploi t'a permis de réinvestir, dans une tâche authentique et motivante, l'ensemble des acquis du chapitre : le \cle{lexique de la vie économique} (métiers, entreprise, salaire, candidature), le \cle{comparatif} et le \cle{superlatif} dans un texte publicitaire, la \cle{subordonnée avec wenn} et le \cle{futur} dans une lettre de motivation, et le \cle{Perfekt} / \cle{Präteritum} dans un compte rendu. Tu as également pratiqué les trois types de textes du programme — petite annonce, lettre formelle, compte rendu — ainsi que la traduction de l'allemand vers le français (version) et du français vers l'allemand (thème). Enfin, la relecture systématique avec la grille d'auto-correction t'a entraîné à vérifier les points critiques pour un francophone : majuscules des noms, place du verbe (principale, subordonnée, infinitifs), cas gouvernés par les prépositions, particules séparables, Umlaute et \all{ß}. Ces compétences seront réévaluées dans le devoir de synthèse du chapitre : garde précieusement ta grille de correction dans ton cahier.

\newpage

\coursec{Méthodes et exercices résolus}

\coursec{Leçon ALA09 : Production et traduction de textes liés à la vie économique}

\courssub{Texte support et compréhension}

\begin{textebox}[Forum économique jeunesse Yaoundé–Berlin — Article de presse simulé]
\all{Im Juni 2024 fand in Yaoundé das „Forum Wirtschaft Jugend“ statt. Zwanzig junge Unternehmer aus Kamerun und zehn Start-up-Gründer aus Deutschland trafen sich im Konferenzzentrum. Die Eröffnungsrede hielt der kamerunische Minister für Klein- und Mittelbetriebe. Er betonte: „Die Wirtschaft wächst nur, wenn junge Leute Chancen bekommen.“ Danach präsentierten die Teilnehmer ihre Projekte: eine App für Kakaobauern, ein Recycling-Unternehmen für Plastikmüll und eine Solarenergie-Firma für Dörfer ohne Strom. Ein deutscher Investor sagte: „Diese Ideen sind besser als viele klassische Projekte, weil sie lokale Probleme lösen.“ Am Nachmittag gab es Workshops zu Themen wie „Businessplan schreiben“ und „Finanzierung finden“. Schließlich unterzeichneten zwei kamerunische Firmen Kooperationsverträge mit deutschen Partnern. Der Besuch war sehr nützlich für den Austausch von Erfahrungen.}
\end{textebox}

\begin{definition}[Glossaire du texte (mots difficiles)]
\begin{center}
\begin{tabularx}{\linewidth}{|l|X|X|}
\hline
\rowcolor{popblueL} \textbf{Deutsch} & \textbf{Français} & \textbf{Remarque} \\
\hline
\all{das Forum, die Foren} & le forum &  \\
\all{der Unternehmer, die Unternehmer} & l'entrepreneur &  \\
\all{der Start-up-Gründer, die -Gründer} & le fondateur de start-up & nom composé \\
\all{das Konferenzzentrum, die -zentren} & le centre de conférences &  \\
\all{die Eröffnungsrede, die -n} & le discours d'ouverture & \\
\all{betonen} & insister sur, souligner & verbe régulier \\
\all{die Chance, die -n} & la chance, l'opportunité &  \\
\all{der Kakaobauer, die -n} & le planteur de cacao & \\
\all{das Recycling-Unternehmen, die -n} & l'entreprise de recyclage & \\
\all{der Plastikmüll} & les déchets plastiques & indénombrable \\
\all{die Solarenergie} & l'énergie solaire & \\
\all{der Investor, die -en} & l'investisseur & \\
\all{lokal} & local & adjectif \\
\all{der Workshop, die -s} & l'atelier (de travail) & anglicisme intégré \\
\all{der Businessplan, die -s} & le plan d'affaires & \\
\all{die Finanzierung, die -en} & le financement & \\
\all{unterzeichnen} & signer (un contrat) & verbe fort : unterzeichnet, hat unterzeichnet \\
\all{der Kooperationsvertrag, die -\"{e}} & le contrat de coopération & \\
\all{der Austausch, die -e} & l'échange &  \\
\all{die Erfahrung, die -en} & l'expérience &  \\
\hline
\end{tabularx}
\end{center}
\end{definition}

\begin{exercice}[Compréhension du texte]
Répondez en allemand par des phrases complètes.
\begin{enumerate}
\item Wann und wo fand das Forum statt?
\item Wer hielt die Eröffnungsrede?
\item Nennen Sie drei Projekte, die präsentiert wurden.
\item Was sagte der deutsche Investor über die Ideen?
\item Was geschah am Ende des Forums?
\end{enumerate}
\end{exercice}

\begin{corrige}[Exercice 1]
\begin{enumerate}
\item \all{Das Forum fand im Juni 2024 in Yaoundé statt.}
\item \all{Der kamerunische Minister für Klein- und Mittelbetriebe hielt die Eröffnungsrede.}
\item \all{Eine App für Kakaobauern, ein Recycling-Unternehmen für Plastikmüll und eine Solarenergie-Firma.}
\item \all{Er sagte: „Diese Ideen sind besser als viele klassische Projekte, weil sie lokale Probleme lösen.“}
\item \all{Zwei kamerunische Firmen unterzeichneten Kooperationsverträge mit deutschen Partnern.}
\end{enumerate}
\end{corrige}

\courssub{Lexique thématique : la vie économique}

\begin{aretenir}[Vocabulaire de base — Vie économique (Niveau A1+/A2)]
\textbf{Les acteurs et les lieux}
\begin{center}
\begin{tabular}{|l|l|l|}
\hline
\rowcolor{popblueL} \textbf{Deutsch (Article + Pluriel)} & \textbf{Français} & \textbf{Exemple} \\
\hline
\all{der Markt, die Märkte} & le marché & \all{Der Markt in Mokolo ist groß.} \\
\all{die Firma, die Firmen / das Unternehmen, die Unternehmen} & l'entreprise & \all{Sie arbeitet in einer großen Firma.} \\
\all{die Bank, die Banken} & la banque & \all{Er geht zur Bank.} \\
\all{der Unternehmer, die Unternehmer} & l'entrepreneur & \all{Ein erfolgreicher Unternehmer.} \\
\all{der Arbeitgeber, die Arbeitgeber} & l'employeur & \all{Der Arbeitgeber zahlt das Gehalt.} \\
\all{der Arbeitnehmer, die Arbeitnehmer} & l'employé & \all{Die Arbeitnehmer haben Rechte.} \\
\all{der Bewerber, die Bewerber / die Bewerberin, die -innen} & le candidat & \all{Der Bewerber hat Erfahrung.} \\
\all{der Kollege, die Kollegen / die Kollegin, die -innen} & le collègue & \all{Meine Kollegen sind nett.} \\
\all{der Kunde, die Kunden / die Kundin, die -innen} & le client & \all{Der Kunde ist König.} \\
\hline
\end{tabular}
\end{center}

\textbf{Les actions et le contrat}
\begin{center}
\begin{tabular}{|l|l|l|}
\hline
\rowcolor{popblueL} \textbf{Deutsch (Infinitif)} & \textbf{Français} & \textbf{Exemple (Präsens)} \\
\hline
\all{produzieren} & produire & \all{Die Firma produziert Kakao.} \\
\all{verkaufen} & vendre & \all{Wir verkaufen Früchte.} \\
\all{kaufen} & acheter & \all{Er kauft ein Auto.} \\
\all{investieren} & investir & \all{Sie investieren in Solarenergie.} \\
\all{einstellen} & embaucher & \all{Die Bank stellt neue Leute ein.} \\
\all{kündigen} & licencier / démissionner & \all{Er kündigt den Vertrag.} \\
\all{bewerben (sich \textbf{um} + Akk)} & postuler & \all{Ich bewerbe mich um die Stelle.} \\
\all{arbeiten} & travailler & \all{Wir arbeiten gut zusammen.} \\
\all{verdienen} & gagner (de l'argent) & \all{Er verdient gut.} \\
\hline
\end{tabular}
\end{center}

\textbf{L'argent et les indicateurs}
\begin{center}
\begin{tabular}{|l|l|lcccccccccccc|}
\hline
\rowcolor{popblueL} \textbf{Deutsch} & \textbf{Français} & \textbf{Remarque} \\
\hline
\all{das Geld} & l'argent (indén.) & pas de pluriel usuel \\
\all{der Preis, die Preise} & le prix &  \\
\all{das Gehalt, die Gehälter} & le salaire (mensuel) &  \\
\all{der Lohn, die Löhne} & le salaire (horaire/hebdo) &  \\
\all{die Steuer, die Steuern} & l'impôt, la taxe & \\
\all{der Gewinn, die Gewinne} & le bénéfice & \\
\all{der Verlust, die Verluste} & la perte & \\
\all{der Vertrag, die Verträge} & le contrat & \\
\all{die Rechnung, die Rechnungen} & la facture & \\
\all{das Angebot, die Angebote} & l'offre / le devis & \\
\all{die Nachfrage} & la demande & indénombrable \\
\hline
\end{tabular}
\end{center}

\textbf{Adjectifs utiles \& Connecteurs logiques}
\begin{center}
\begin{tabular}{|l|l|l|}
\hline
\rowcolor{popblueL} \textbf{Deutsch} & \textbf{Français} & \textbf{Exemple} \\
\hline
\all{wirtschaftlich} & économique & \all{das wirtschaftliche Wachstum} \\
\all{rentabel} & rentable & \all{ein rentables Geschäft} \\
\all{wettbewerbsfähig} & compétitif & \all{die Firma ist wettbewerbsfähig} \\
\all{verhandelbar} & négociable & \all{Der Preis ist verhandelbar.} \\
\all{arbeitslos} & au chômage & \all{Er ist arbeitslos.} \\
\hline
\all{zuerst / zuerst einmal} & d'abord / tout d'abord &  \\
\all{dann / danach} & puis / ensuite &  \\
\all{schließlich / zum Schluss} & enfin / finalement &  \\
\all{deshalb / daher} & c'est pourquoi / donc & \textbf{verbe position 2 !} \\
\all{jedoch / aber} & cependant / mais &  \\
\all{außerdem} & de plus / en outre &  \\
\hline
\end{tabular}
\end{center}
\end{aretenir}



\courssub{Grammaire 1 : Révision du comparatif et du superlatif}

\begin{propriete}[Comparatif et superlatif — Formes et emploi]
\textbf{A. Comparatif d'inégalité (\all{...-er als})}
\begin{itemize}
\item Règle générale : adjectif/adverbe + \all{-er} + \all{als} + \textbf{Nominatif}.
\item \all{schnell} $\rightarrow$ \all{schneller als} ; \all{teuer} $\rightarrow$ \all{teurer als} ; \all{interessant} $\rightarrow$ \all{interessanter als}.
\item \textbf{Umlaut} pour les monosyllabes à voyelle \all{a, o, u} : \all{groß $\rightarrow$ größer}, \all{alt $\rightarrow$ älter}, \all{kalt $\rightarrow$ kälter}, \all{jung $\rightarrow$ jünger}, \all{klug $\rightarrow$ klüger}, \all{stark $\rightarrow$ stärker}.
\item \textbf{Irréguliers majeurs} : \all{gut $\rightarrow$ besser}, \all{viel $\rightarrow$ mehr}, \all{gerne $\rightarrow$ lieber}, \all{hoch $\rightarrow$ höher}, \all{nah $\rightarrow$ näher}.
\end{itemize}

\textbf{B. Comparatif d'égalité (\all{so ... wie})}
\begin{itemize}
\item \all{so} + adjectif (forme de base) + \all{wie} + Nominatif.
\item \all{Der Markt in Douala ist so groß wie der Markt in Yaoundé.}
\end{itemize}

\textbf{C. Superlatif}
\begin{itemize}
\item \textbf{Prédicatif (après verbe) :} \all{am} + adjectif + \all{-sten} : \all{am größten, am besten, am teuersten, am interessantesten}.
\item \textbf{Attributif (devant nom) :} article décliné + adjectif + \all{-ste} + déclinaison adjectivale.
\item \all{der größte Markt, die größte Firma, das größte Auto, die größten Märkte}.
\item \textbf{Irréguliers :} \all{am besten, am meisten, am liebsten, am höchsten, am nächsten}.
\end{itemize}

\textbf{D. Cas particulier : \all{als} vs \all{wie}}
\begin{center}
\begin{tabularx}{\linewidth}{|X|X|}
\hline
\rowcolor{popblueL} \all{größer \textbf{als}} & \all{so groß \textbf{wie}} \\
\hline
Comparatif de supériorité & Comparatif d'égalité \\
\hline
\end{tabularx}
\end{center}
\end{propriete}

\begin{exemplebox}[Exemples économiques]
\begin{enumerate}
\item \all{Die Wirtschaft in Kamerun wächst \textbf{schneller als} in manchen europäischen Ländern.}
\item \all{Ein Praktikum in einer Bank ist \textbf{besser als} gar keine Erfahrung.}
\item \all{Die Preise auf dem Markt sind \textbf{so hoch wie} im Supermarkt.}
\item \all{Das ist \textbf{das rentabelste} Projekt des Jahres. (Superlatif attributif)}
\item \all{Diese Idee ist \textbf{am innovativsten}. (Superlatif prédicatif)}
\end{enumerate}
\end{exemplebox}

\begin{attention}[Pièges comparatif — Spécifiques francophones]
\begin{enumerate}
\item \textbf{Oubli de l'Umlaut} : \all{*großer} au lieu de \all{größer}, \all{*alter} au lieu de \all{älter}. Vérifiez la voyelle radicale (a, o, u).
\item \textbf{Mauvais cas après \all{als}} : \all{*als mich} (accusatif) au lieu de \all{als ich} (nominatif). \all{Er ist älter als ich.}
\item \textbf{Confusion \all{mehr} / \all{viel}} : \all{viel} = beaucoup (quantité), \all{mehr} = plus (comparatif). \all{Ich habe viel Geld. Ich brauche mehr Geld.}
\item \textbf{Superlatif prédicatif sans \all{am}} : \all{*Das ist besten} $\rightarrow$ \all{Das ist \textbf{am} besten}.
\end{enumerate}
\end{attention}

\courssub{Grammaire 2 : Les phrases conditionnelles avec \all{si} (\all{wenn}) et le futur}

\begin{propriete}[\all{wenn}-Sätze (Conditionnelles réalisables — Présent / Futur)]
\textbf{Structure :} Subordonnée \all{wenn} + \textbf{Présent} (verbe à la fin) — Principale + \textbf{Futur I} (ou Présent + indicateur de temps).
\begin{center}
\begin{tabularx}{\linewidth}{|X|X|}
\hline
\rowcolor{popblueL} \textbf{Position de la subordonnée} & \textbf{Ordre des mots (Principale)} \\
\hline
\all{Wenn ich Geld \textbf{habe},} \all{\textbf{werde} ich einen Computer \textbf{kaufen}.} & Verbe \all{werde} en \textbf{position 2} (inversion). \\
\hline
\all{Ich \textbf{werde} einen Computer \textbf{kaufen},} \all{\textbf{wenn} ich Geld \textbf{habe}.} & Verbe \all{werde} en \textbf{position 2} (sujet d'abord). \\
\hline
\end{tabularx}
\end{center}

\textbf{Règles clés :}
\begin{enumerate}
\item \all{wenn} = « si » (condition). Le verbe conjugué va \textbf{à la fin} de la subordonnée.
\item Si la subordonnée commence la phrase, le verbe de la principale saute en \textbf{position 2} (avant le sujet) : \all{..., verb ...}.
\item \textbf{Futur I ( \all{werden} + Infinitif )} : exprime l'avenir, une intention ferme ou une supposition.
\item \textbf{Conjugaison de \all{werden} (Präsens) :} \all{ich werde, du wirst, er/sie/es wird, wir werden, ihr werdet, sie/Sie werden}. (Verbe fort : voyelle \all{e} $\rightarrow$ \all{i} aux 2\textsuperscript{e} et 3\textsuperscript{e} pers. sg).
\item \textbf{Infinitif à la fin de la phrase principale} : \all{Ich werde morgen \textbf{arbeiten}.}
\end{enumerate}

\textbf{Distinction \all{wenn} / \all{falls} / \all{ob} :}
\begin{itemize}
\item \all{wenn} : condition générale, temporelle (« quand » au passé) ou future réelle.
\item \all{falls} : « au cas où », possibilité plus faible, plus formel.
\item \all{ob} : « si » dans les questions indirectes (\all{Ich weiß nicht, ob er kommt.}), \textbf{jamais} pour condition.
\end{itemize}
\end{propriete}

\begin{exemplebox}[Exemples \all{wenn} + Futur — Contexte économique]
\begin{enumerate}
\item \all{Wenn die Wirtschaft wächst, \textbf{werden} mehr Jugendliche Arbeit \textbf{finden}.} (Subordonnée en tête $\rightarrow$ inversion).
\item \all{Wir \textbf{werden} die Fabrik \textbf{eröffnen}, \textbf{falls} der Kredit genehmigt \textbf{wird}.} (Passif \all{wird genehmigt} à la fin).
\item \all{Wenn du \textbf{suchst}, \textbf{findest} du eine Stelle.} (Présent dans les deux : loi générale / conseil).
\item \all{Ich \textbf{werde} mich \textbf{bewerben}, \textbf{wenn} ich das Zeugnis \textbf{habe}.}
\end{enumerate}
\end{exemplebox}

\begin{attention}[Erreurs fréquentes \all{wenn} / Futur]
\begin{enumerate}
\item \textbf{Verbe pas à la fin après \all{wenn}} : \all{*Wenn ich habe Geld...} $\rightarrow$ \all{Wenn ich Geld \textbf{habe}...}
\item \textbf{Pas d'inversion après subordonnée initiale} : \all{*Wenn ich Geld habe, ich werde kaufen...} $\rightarrow$ \all{..., \textbf{werde} ich kaufen.}
\item \textbf{Oubli de l'infinitif final} : \all{*Ich werde einen Computer kaufen.} (Correct) mais \all{*Ich werde kaufen einen Computer.} (Faux : infinitif \textbf{toujours} à la fin).
\item \textbf{Confusion \all{werden} (futur) / \all{werden} (devenir)} : \all{Er wird Arzt (devenir). Er wird arbeiten (futur).}
\item \textbf{\all{ob} au lieu de \all{wenn}} : \all{*Ob ich Zeit habe, komme ich.} $\rightarrow$ \all{\textbf{Wenn} ich Zeit habe, komme ich.}
\end{enumerate}
\end{attention}

\courssub{Types de textes économiques : petite annonce, lettre, compte rendu}

\textbf{La petite annonce (\all{die Kleinanzeige})}

\begin{methode}[Rédiger une petite annonce (\all{die Kleinanzeige})]
La petite annonce est le texte économique le plus fréquent au Bac (offre d'emploi, vente, recherche). Voici la démarche en 5 étapes :
\begin{enumerate}
\item \textbf{Identifier le type d'annonce} : offre (\all{Ich biete...}) ou demande (\all{Ich suche...}). Le verbe est toujours en position 2.
\item \textbf{Rédiger un titre court} : nom + majuscule, sans article si possible : \all{Verkaufe Fahrrad} / \all{Suche Praktikum}.
\item \textbf{Décrire l'objet ou le poste} en 2-3 phrases simples : état (\all{neu, gebraucht, in gutem Zustand}), prix (\all{Preis: 25\,000 FCFA, verhandelbar} = négociable), qualités recherchées (\all{motiviert, pünktlich, freundlich}).
\item \textbf{Utiliser les comparatifs} si besoin : \all{besser als neu}, \all{günstiger als im Geschäft}.
\item \textbf{Terminer par les coordonnées} : \all{Bei Interesse: Tel. 6\,90\,00\,00\,00 (Yaoundé).} Relire : majuscules aux noms, verbe en position 2, orthographe.
\end{enumerate}
\textbf{Structures à retenir :} \all{Ich biete/suche...} ; \all{Der Preis ist verhandelbar.} ; \all{Bei Interesse rufen Sie mich an.} ; \all{Ab sofort} (dès maintenant).
\end{methode}

\begin{exoresolu}[Petite annonce : vente d'un téléphone]
\textbf{Énoncé (type Bac) :} Vous vendez votre téléphone portable. Rédigez en allemand une petite annonce de 4 à 6 lignes : description, prix (30\,000 FCFA, négociable), contact à Douala.
\tcblower
\textbf{Raisonnement :} il s'agit d'une \emph{offre} : on commence par \all{Verkaufe...} ou \all{Ich biete...}. On place le verbe en position 2, on met la majuscule à tous les noms (\all{Handy, Preis, Zustand}), et on termine par le contact.

\textbf{Solution rédigée :}

\all{Verkaufe Handy Samsung}

\all{Ich biete ein Handy Samsung, nur sechs Monate alt und in sehr gutem Zustand. Es funktioniert besser als viele neue Handys und hat zwei SIM-Karten. Der Preis ist 30\,000 FCFA, verhandelbar. Bei Interesse: Tel. 6\,99\,12\,34\,56 (Douala).}

« Vends téléphone Samsung. Je propose un téléphone Samsung, âgé de six mois seulement et en très bon état. Il fonctionne mieux que beaucoup de téléphones neufs et a deux cartes SIM. Prix : 30\,000 FCFA, négociable. Si intéressé : tél. 6\,99\,12\,34\,56 (Douala). »

\textbf{Vérifications :} majuscules à \all{Handy, Zustand, SIM-Karten, Preis, Interesse} ; comparatif \all{besser als} + nominatif ; \all{verhandelbar} sans article ; verbe en position 2 dans chaque phrase.
\end{exoresolu}

\begin{exercice}[Petite annonce : offre de stage]
Rédigez une petite annonce (5 lignes) : Une entreprise de Yaoundé cherche un stagiaire en marketing (allemand obligatoire, motivé, travail d'équipe). Contact : email.
\end{exercice}

\textbf{La lettre formelle (\all{der formelle Brief} / \all{die Bewerbung})}

\begin{methode}[Rédiger une lettre formelle (\all{der Brief})]
\begin{enumerate}
\item \textbf{L'en-tête} : lieu et date à droite : \all{Yaoundé, den 15. März 2025} (attention : \all{den} + ordinal + point).
\item \textbf{La formule d'appel} : \all{Sehr geehrte Damen und Herren,} (inconnu) ou \all{Sehr geehrter Herr Mbarga,} (connu, masculin) / \all{Sehr geehrte Frau Ngono,} (féminin). Virgule obligatoire, puis ligne suivante en \emph{minuscule}.
\item \textbf{Le corps} : 3 paragraphes — motif (\all{Ich schreibe Ihnen, weil...}), arguments (formation, expérience, qualités), demande polie (\all{Ich würde mich über eine positive Antwort freuen.}).
\item \textbf{La formule de politesse} : \all{Mit freundlichen Grüßen} (sans virgule !), puis la signature.
\item \textbf{La relecture} : \all{Sie} avec majuscule partout (vouvoiement), verbe en position 2, \all{weil} + verbe à la fin.
\end{enumerate}
\textbf{Structures à retenir :} \all{Ich möchte mich als Verkäufer bewerben.} ; \all{Ich habe Erfahrung im Verkauf.} ; \all{Ich bin motiviert und pünktlich.} ; \all{Über eine Einladung zum Gespräch freue ich mich.}
\end{methode}

\begin{exoresolu}[Lettre de candidature à un stage]
\textbf{Énoncé (type Bac) :} Vous voulez faire un stage dans une banque à Yaoundé pendant les vacances. Rédigez en allemand une lettre courte (8-10 lignes) : appel, motif, vos qualités, formule de politesse.
\tcblower
\textbf{Raisonnement :} destinataire inconnu $\rightarrow$ \all{Sehr geehrte Damen und Herren,}. Corps en 3 paragraphes. Vouvoiement : \all{Sie, Ihnen, Ihre} avec majuscule. Subordonnées avec \all{weil} : verbe conjugué à la fin.

\textbf{Solution rédigée :}

\all{Yaoundé, den 15. März 2025}

\all{Sehr geehrte Damen und Herren,}

\all{ich bin Schülerin der Klasse Première A am Gymnasium Yaoundé und ich möchte mich für ein Praktikum in Ihrer Bank bewerben.}

\all{Ich lerne seit vier Jahren Deutsch und ich spreche gut Französisch und Englisch. Ich bin motiviert, pünktlich und ich arbeite gern mit Kunden. Ich möchte später in der Wirtschaft arbeiten, deshalb ist dieses Praktikum sehr wichtig für mich.}

\all{Ich würde mich sehr über eine positive Antwort freuen.}

\all{Mit freundlichen Grüßen}

\all{Aminatou Bello}

\textbf{Vérifications :} minuscule après la virgule de l'appel (\all{ich bin...}) ; \all{Ihrer, Ihnen} avec majuscule ; \all{deshalb} + verbe en position 2 (\all{deshalb ist}) ; pas de virgule après \all{Mit freundlichen Grüßen}.
\end{exoresolu}

\begin{exercice}[Lettre de réclamation]
Vous avez acheté un ordinateur portable qui ne fonctionne pas. Écrivez une lettre formelle au vendeur (Magasin « TechWorld », Douala) : description du problème, demande de réparation ou remboursement, délai.
\end{exercice}

\textbf{Le compte rendu (\all{der Bericht})}

\begin{methode}[Rédiger un compte rendu (\all{der Bericht})]
Le compte rendu rapporte des faits (visite d'entreprise, marché, foire économique) de façon \emph{objective}.
\begin{enumerate}
\item \textbf{Titre + phrase d'introduction} : qui ? quoi ? quand ? où ? \all{Am Freitag besuchte unsere Klasse den Markt von Mokolo.}
\item \textbf{Développement au Präteritum} (ou Perfekt à l'oral) : faits dans l'ordre chronologique, connecteurs \all{zuerst, dann, danach, schließlich}.
\item \textbf{Chiffres et détails} : \all{Es gab viele Händler. Die Preise waren höher als letztes Jahr.}
\item \textbf{Conclusion brève} : bilan ou impression générale : \all{Der Besuch war sehr interessant und nützlich.}
\item \textbf{Relecture} : concordance des temps (passé !), verbe en position 2, pas d'opinion personnelle excessive.
\end{enumerate}
\textbf{Präteritum utile (verbes forts/irréguliers) :} \all{war (sein), hatte (haben), gab (geben), besuchte (besuchen), sah (sehen), kaufte (kaufen), kostete (kosten), fand (finden), ging (gehen), sprach (sprechen), unterschrieb (unterschreiben)}.
\end{methode}

\begin{exoresolu}[Compte rendu d'une visite au marché]
\textbf{Énoncé (type Bac) :} Votre classe a visité le marché de Mokolo à Yaoundé. Rédigez en allemand un compte rendu de 6 à 8 lignes au passé.
\tcblower
\textbf{Raisonnement :} introduction avec les 4 questions (qui/quoi/quand/où), faits au Präteritum/Perfekt, connecteurs chronologiques, comparatif possible (\all{höher als}).

\textbf{Solution rédigée :}

\all{Besuch am Markt Mokolo}

\all{Am Freitag, den 10. Mai, besuchte unsere Klasse mit unserer Deutschlehrerin den Markt Mokolo in Yaoundé. Zuerst sahen wir die Obst- und Gemüseabteilung. Es gab sehr viele Händler und Kunden. Dann sprachen wir mit einer Händlerin: Sie verkauft seit zwanzig Jahren Tomaten und Zwiebeln. Die Preise waren höher als letztes Jahr, weil der Transport teurer wurde. Danach kauften wir Früchte für das Klassenfest. Schließlich machten wir Fotos für unsere Schülerzeitung. Der Besuch war sehr interessant und nützlich für unsere Wirtschaftsarbeit.}

\textbf{Vérifications :} Präteritum correct : \all{besuchte, sahen (sehen), gab, waren, wurde, kauften, machten, war} ; \all{weil} + verbe à la fin (\all{teurer wurde}) ; comparatif \all{höher als}, \all{teurer} ; connecteurs en début de phrase + verbe en position 2 (\all{Zuerst sahen wir...}).
\end{exoresolu}

\begin{exercice}[Compte rendu : Foire de l'emploi]
Votre classe a participé à la « Foire de l'Emploi Jeunes » à Douala. Rédigez un compte rendu (Präteritum) : date, lieu, stands visités, entretien avec un recruteur, impression finale.
\end{exercice}

\courssub{Traduction : thème et version}

\textbf{Le thème (Français $\rightarrow$ Allemand)}

\begin{methode}[Réussir le thème en 6 étapes]
\begin{enumerate}
\item \textbf{Lire toute la phrase} et repérer le temps (présent, passé composé $\rightarrow$ Perfekt, futur $\rightarrow$ Futur I ou présent + indication de temps).
\item \textbf{Repérer la structure} : phrase principale (verbe en position 2) ou subordonnée (\all{weil, wenn, dass} $\rightarrow$ verbe à la fin).
\item \textbf{Traduire les noms avec leur article} : \emph{le marché} = \all{der Markt}, \emph{l'économie} = \all{die Wirtschaft}, \emph{l'argent} = \all{das Geld}.
\item \textbf{Attention aux cas} : après \all{für, durch, ohne, gegen} $\rightarrow$ accusatif ; après \all{mit, von, zu, bei, seit, in} (position) $\rightarrow$ datif. Ex: \all{für die Wirtschaft} (Akk), \all{mit dem Geld} (Dat).
\item \textbf{Construire la phrase allemande} : Sujet + Verbe (pos 2) + Compléments (Règle TeKaMoLo : Temps — Causalité — Manière — Lieu). Ex: \all{Ich arbeite \textbf{morgen} (T) \textbf{in Douala} (O).}
\item \textbf{Relire} : majuscules aux noms, Umlaute, accord du verbe, place du verbe.
\end{enumerate}
\textbf{Réflexe Futur :} « je travaillerai » = \all{ich werde arbeiten} (\all{werden} conjugué en pos 2 + infinitif à la fin).
\textbf{Réflexe « Depuis » :} « depuis 2 ans » = \all{seit zwei Jahren} + \textbf{Présent} (\all{Ich lerne seit zwei Jahren Deutsch}).
\end{methode}

\begin{exoresolu}[Thème : phrases économiques]
\textbf{Énoncé (type Bac) :} Traduisez en allemand :
\begin{enumerate}
\item Le marché de Douala est plus grand que le marché de Bafoussam.
\item Si j'ai de l'argent, j'achèterai un ordinateur pour mes études.
\item Ma sœur travaille depuis deux ans dans une banque parce qu'elle aime l'économie.
\end{enumerate}
\tcblower
\textbf{1.} Comparatif : \all{größer als} + nominatif. Noms avec articles : \all{der Markt}. $\rightarrow$ \all{Der Markt von Douala ist größer als der Markt von Bafoussam.} (Justification : \all{groß – größer – am größten}, comparatif avec Umlaut.)

\textbf{2.} Phrase conditionnelle : \all{wenn} + verbe à la fin ; futur : \all{werde kaufen} avec infinitif final. $\rightarrow$ \all{Wenn ich Geld habe, werde ich einen Computer für mein Studium kaufen.} (Justification : quand la subordonnée \all{wenn} est en tête, le verbe de la principale suit immédiatement : \all{..., werde ich...} ; \all{einen Computer} = accusatif masculin ; \all{für mein Studium} = accusatif.)

\textbf{3.} \all{seit} + datif + présent (action qui continue !) ; \all{weil} + verbe à la fin. $\rightarrow$ \all{Meine Schwester arbeitet seit zwei Jahren in einer Bank, weil sie die Wirtschaft liebt.} (Justification : « depuis deux ans » = \all{seit zwei Jahren}, datif pluriel en \emph{-n} ; en allemand on garde le \emph{présent}, pas le passé ; \all{in einer Bank} = datif féminin.)
\end{exoresolu}

\begin{exercice}[Thème : phrases variées]
Traduisez en allemand :
\begin{enumerate}
\item Notre entreprise est plus rentable que l'année dernière.
\item Si vous investissez dans l'énergie solaire, vous gagnerez de l'argent.
\item Je cherche un emploi depuis trois mois, mais je n'ai pas encore trouvé.
\item Le patron veut embaucher de nouveaux employés pour la nouvelle usine.
\end{enumerate}
\end{exercice}

\begin{corrige}[Exercice Thème]
\begin{enumerate}
\item \all{Unsere Firma ist rentabler als letztes Jahr.} (\all{rentabel – rentabler}, \all{als} + Nominativ \all{letztes Jahr}).
\item \all{Wenn Sie in Solarenergie investieren, werden Sie Geld verdienen.} (\all{wenn} + verbe fin \all{investieren}, Futur \all{werden ... verdienen}, \all{in} + Datif \all{Solarenergie}).
\item \all{Ich suche seit drei Monaten einen Job, aber ich habe noch keinen gefunden.} (\all{seit} + Datif pluriel \all{Monaten} + Présent \all{suche}, \all{aber} + pos 2 \all{habe}, \all{keinen} accusatif masc).
\item \all{Der Chef will neue Mitarbeiter für die neue Fabrik einstellen.} (\all{wollen} + infinitif \all{einstellen} à la fin, \all{für} + Akk \all{die neue Fabrik}).
\end{enumerate}
\end{corrige}

\textbf{La version (Allemand $\rightarrow$ Français)}

\begin{methode}[Réussir la version en 5 étapes]
\begin{enumerate}
\item \textbf{Lire le texte entier} une première fois pour saisir le sens global (sujet, époque, lieu).
\item \textbf{Repérer le verbe conjugué} de chaque phrase (position 2 ou finale) et son temps : Perfekt $\rightarrow$ passé composé ; Präteritum $\rightarrow$ passé simple / imparfait ; Futur I $\rightarrow$ futur simple.
\item \textbf{Traduire les mots dans leur contexte} : \all{die Bank} = la banque \emph{ou} le banc ; \all{bekommen} = recevoir / obtenir (pas « devenir » !) ; \all{die Fabrik} = l'usine ; \all{der Chef} = le patron (pas le chef cuisinier).
\item \textbf{Rédiger en bon français} : une version n'est pas un mot à mot ; reformulez naturellement : \all{Es gibt viele Arbeitsplätze} $\rightarrow$ « Il y a beaucoup d'emplois / de postes ».
\item \textbf{Relire} : rien d'oublié ? chaque groupe de mots allemand doit avoir son équivalent français. Attention aux faux amis (\all{aktuel} = actuel / courant, \all{eventuell} = éventuel).
\end{enumerate}
\end{methode}

\begin{exoresolu}[Version : un texte sur l'économie camerounaise]
\textbf{Énoncé (type Bac) :} Traduisez en français le texte suivant :

\all{Die Wirtschaft in Kamerun wächst langsam. Viele junge Leute suchen eine Arbeit, aber es gibt nicht genug Arbeitsplätze. Letztes Jahr hat eine deutsche Firma in Douala eine neue Fabrik eröffnet. Sie hat zweihundert Arbeiter bekommen. Wenn die Wirtschaft stärker wird, werden mehr Jugendliche eine Stelle finden.}
\tcblower
\textbf{Analyse grammaticale :}
\begin{itemize}
\item \all{wächst} (Présent, \all{wachsen}) = croît.
\item \all{suchen} (Présent) = cherchent.
\item \all{hat ... eröffnet} (Perfekt) = a ouvert (passé composé).
\item \all{hat ... bekommen} (Perfekt) = a recruté / a obtenu (faux ami : \all{bekommen} $\neq$ devenir).
\item \all{Wenn die Wirtschaft stärker \textbf{wird}} : Subordonnée \all{wenn} + verbe \all{wird} (Présent de \all{werden}) à la fin $\rightarrow$ « Si l'économie devient plus forte » (Présent en français dans clause \emph{si} + futur dans principale).
\item \all{werden ... finden} (Futur I) = trouveront.
\end{itemize}

\textbf{Traduction proposée :}
« L'économie au Cameroun croît lentement. Beaucoup de jeunes cherchent un emploi, mais il n'y a pas assez de postes. L'année dernière, une entreprise allemande a ouvert une nouvelle usine à Douala. Elle a recruté deux cents ouvriers. Si l'économie devient plus forte, davantage de jeunes trouveront un emploi. »

\textbf{Vérifications :} \all{die Firma} = l'entreprise ; \all{die Stelle} = le poste/l'emploi ; \all{die Jugendlichen} = les jeunes (adjectif nominalisé, majuscule) ; futur allemand rendu par futur français.
\end{exoresolu}

\begin{exercice}[Version : texte court]
Traduisez en français :
\all{Gestern hat unser Lehrer uns eine Firma in Douala gezeigt. Der Chef hat uns erklärt, wie man einen Businessplan schreibt. Das war sehr interessant, weil wir viel Neues gelernt haben. Wenn ich groß bin, werde ich auch Unternehmer.}
\end{exercice}

\begin{corrige}[Exercice Version]
« Hier, notre professeur nous a montré une entreprise à Douala. Le patron nous a expliqué comment on rédige un plan d'affaires. C'était très intéressant, car nous avons appris beaucoup de choses nouvelles. Quand je serai grand, je deviendrai aussi entrepreneur. »
\textit{Notes :} \all{hat ... gezeigt} (Perfekt), \all{hat ... erklärt} (Perfekt), \all{wie man ... schreibt} (subordonnée indirecte, verbe à la fin), \all{war} (Präteritum), \all{haben ... gelernt} (Perfekt), \all{werde ... werden} (Futur I, \all{werden} comme verbe plein « devenir »).
\end{corrige}

\courssub{Relire et corriger sa production}

\begin{methode}[La grille de relecture en 5 points]
Avant de rendre toute production ou traduction, vérifiez dans l'ordre :
\begin{enumerate}
\item \textbf{Majuscules} : tous les noms allemands prennent la majuscule (\all{das Geld, die Arbeit, der Markt, die Wirtschaft}) ; \all{Sie/Ihnen/Ihr} de vouvoiement aussi.
\item \textbf{Place du verbe} : position 2 en principale ; finale après \all{weil, wenn, dass, obwohl, bevor} ; infinitif ou participe à la fin (\all{Ich werde morgen arbeiten.} / \all{Ich habe gestern gearbeitet.}).
\item \textbf{Cas} : accusatif après \all{für, durch, ohne, gegen, um...zu} ; datif après \all{mit, von, zu, bei, seit, in} (position), \all{aus, außer, gegenüber}.
\item \textbf{Orthographe} : Umlaute (\all{größer, schöner, Müller}), \all{ß} après voyelle longue/diphtongue (\all{Straße, groß, weiß}) vs \all{ss} après voyelle brève (\all{Fluss, muss, Kasse}), \all{schließlich} (Allemagne) et non \emph{schliesslich} (Suisse).
\item \textbf{Concordance des temps} : un récit au passé reste au passé (Präteritum/Perfekt) ; « depuis » + présent avec \all{seit} (\all{Ich lerne seit zwei Jahren Deutsch}).
\end{enumerate}
\end{methode}

\begin{exoresolu}[Correction d'un texte fautif]
\textbf{Énoncé (type Bac) :} Le texte suivant contient 6 erreurs. Trouvez-les et corrigez-les :

\all{Gestern ich habe der markt besucht. Ich habe viele früchte gekauft, weil sie billig waren. Die preise sind besser als im geschäft. Wenn ich geld habe, ich kaufe mehr.}
\tcblower
\textbf{Corrections détaillées :}
\begin{enumerate}
\item \all{Gestern ich habe} $\rightarrow$ \all{Gestern habe ich} : le verbe conjugué reste en position 2, même après un complément de temps en tête.
\item \all{der markt} $\rightarrow$ \all{den Markt} : majuscule au nom + accusatif masculin après \all{besuchen} (\all{der $\rightarrow$ den}).
\item \all{früchte} $\rightarrow$ \all{Früchte} : majuscule obligatoire à tous les noms.
\item \all{preise} $\rightarrow$ \all{Preise} : majuscule.
\item \all{geschäft} $\rightarrow$ \all{Geschäft} : majuscule.
\item \all{Wenn ich geld habe, ich kaufe mehr} $\rightarrow$ \all{Wenn ich Geld habe, kaufe ich mehr} : majuscule à \all{Geld} et, après une subordonnée en tête, le verbe de la principale vient immédiatement (inversion sujet-verbe).
\end{enumerate}
\textbf{Texte corrigé :} \all{Gestern habe ich den Markt besucht. Ich habe viele Früchte gekauft, weil sie billig waren. Die Preise sind besser als im Geschäft. Wenn ich Geld habe, kaufe ich mehr.}
\end{exoresolu}

\begin{attention}[Les 5 erreurs qui coûtent des points au Bac]
\begin{enumerate}
\item \textbf{L'oubli des majuscules aux noms} : \all{der markt, das geld} sont fautifs. En allemand, \emph{tout} nom prend la majuscule : \all{der Markt, das Geld, die Wirtschaft}. C'est la faute la plus sanctionnée.
\item \textbf{Le verbe mal placé} : jamais \all{Gestern ich habe...} ni \all{..., weil sie sind billig}. Principale : verbe en position 2 ; après \all{weil, wenn, dass} : verbe à la fin.
\item \textbf{Les faux amis} : \all{bekommen} = recevoir (devenir = \all{werden}) ; \all{die Bank} = la banque \emph{et} le banc ; \all{aktual} = actuel / courant (et non « actuel » au sens de vrai/réel) ; \all{der Chef} = le patron (et non « le chef » de cuisine) ; \all{eventuell} = éventuel (pas « éventuellement »).
\item \textbf{Les calques du français} : « depuis deux ans » se traduit \all{seit zwei Jahren} + \emph{présent} (\all{Ich lerne seit zwei Jahren Deutsch}), jamais avec un passé composé ; « j'ai 16 ans » = \all{Ich bin 16 Jahre alt} (être, pas avoir) ; « je m'appelle » = \all{Ich heiße} (pas \all{Ich rufe}).
\item \textbf{Les cas après prépositions} : \all{für} + accusatif (\all{für die Wirtschaft}), \all{mit, seit, von, zu, bei, nach} + datif (\all{mit dem Geld, seit zwei Jahren, bei der Bank}). Un mauvais cas dans une lettre ou une annonce fait perdre des points de grammaire.
\end{enumerate}
\end{attention}

\newpage

\coursec{Exercices}

\courssub{Je vérifie mes connaissances}

\begin{exercice}[QCM : grammaire et vie économique]
Pour chaque question, entoure la bonne réponse.
\begin{enumerate}
\item \all{Der neue Markt in Yaoundé ist \ldots{} als der alte.}
\begin{itemize}
\item[a)] \all{größer}
\item[b)] \all{mehr groß}
\item[c)] \all{am größten}
\end{itemize}
\item \all{Wenn ich mehr Geld \ldots{}, kaufe ich einen Computer.}
\begin{itemize}
\item[a)] \all{habe}
\item[b)] \all{hätte}
\item[c)] \all{hat}
\end{itemize}
\item Une petite annonce allemande commence souvent par :
\begin{itemize}
\item[a)] \all{Liebe Frau Mballa,}
\item[b)] \all{Zu verkaufen:}
\item[c)] \all{Sehr geehrte Damen und Herren,}
\end{itemize}
\item \all{Nächstes Jahr \ldots{} ich in Douala arbeiten.}
\begin{itemize}
\item[a)] \all{werde}
\item[b)] \all{will}
\item[c)] \all{würde}
\end{itemize}
\end{enumerate}
\end{exercice}

\begin{exercice}[Vrai ou faux ? Justifie ta réponse]
Réponds par vrai ou faux et justifie en une phrase en français.
\begin{enumerate}
\item Dans une subordonnée introduite par \all{wenn}, le verbe conjugué se place à la fin de la phrase.
\item Le superlatif de \all{gut} est \all{gutest}.
\item Une lettre de motivation allemande se termine toujours par \all{Mit freundlichen Grüßen}.
\item Au futur, l'infinitif se place juste après le sujet.
\item Le mot \all{die Firma} signifie « la ferme ».
\end{enumerate}
\end{exercice}

\begin{exercice}[Vocabulaire : la vie économique]
Complète chaque phrase avec le mot allemand qui convient, choisi dans la liste : \all{der Arbeitsplatz, die Ware, der Kunde, das Geschäft, der Preis, die Rechnung}.
\begin{enumerate}
\item Die \ldots{} auf dem Markt sind heute sehr frisch : Mangos, Bananen und Ananas.
\item Mein Bruder hat einen neuen \ldots{} in einer Bank in Douala gefunden.
\item Der \ldots{} für Kakao ist dieses Jahr gestiegen.
\item Der \ldots{} bezahlt an der Kasse.
\item Nach dem Essen im Restaurant bekommen wir die \ldots{}.
\item Das \ldots{} öffnet jeden Morgen um acht Uhr.
\end{enumerate}
\end{exercice}

\begin{exercice}[Conjugaison à compléter]
Complète le tableau avec les formes correctes du verbe entre parenthèses (Präsens et Futur I).
\begin{center}
\begin{tabular}{|l|l|l|}
\hline
\rowcolor{popblueL} Personne & Präsens & Futur I \\ \hline
ich (arbeiten) & \ldots{} & \ldots{} \\ \hline
du (verdienen) & \ldots{} & \ldots{} \\ \hline
er/sie (kaufen) & \ldots{} & \ldots{} \\ \hline
wir (verkaufen) & \ldots{} & \ldots{} \\ \hline
sie/Sie (gründen) & \ldots{} & \ldots{} \\ \hline
\end{tabular}
\end{center}
\end{exercice}

\courssub{Je m'entraîne}

\begin{exercice}[Thème : phrases à traduire]
Traduis les phrases suivantes en allemand. Attention au comparatif, à la place du verbe et aux majuscules.
\begin{enumerate}
\item Le marché de Douala est plus grand que le marché de Bafoussam.
\item Si j'ai de l'argent, j'achèterai un vélo.
\item Ma sœur travaille dans une banque à Yaoundé.
\item Le cacao est le produit le plus important de notre région.
\item L'année prochaine, nous visiterons une entreprise allemande.
\end{enumerate}
\end{exercice}

\begin{exercice}[Transformation : du présent au futur]
Réécris chaque phrase au Futur I avec \all{werden} + infinitif. Attention à la place de l'infinitif.
\begin{enumerate}
\item \all{Ich lerne Deutsch an der Schule.}
\item \all{Mein Freund sucht eine Stelle in einer Fabrik.}
\item \all{Wir verkaufen unseren Kakao an eine Firma.}
\item \all{Die Wirtschaft in Kamerun wächst.}
\item \all{Du machst ein Praktikum in Douala.}
\end{enumerate}
\end{exercice}

\begin{exercice}[Texte à trous : un compte rendu économique]
Complète le texte avec la forme correcte des mots entre parenthèses (comparatif, superlatif ou verbe conjugué).

\begin{textebox}[Compte rendu à compléter]
Der Markt in Mvog-Mbi ist heute \ldots{} (voll) als gestern. Viele Händler verkaufen ihre Waren \ldots{} (billig) als letzte Woche. Wenn das Wetter gut \ldots{} (sein), kommen mehr Kunden. Die Preise für Tomaten sind am Morgen am \ldots{} (niedrig). Wenn die Händler mehr verdienen \ldots{} (wollen), müssen sie früh aufstehen. Nächstes Jahr \ldots{} (werden) die Stadt einen neuen Markt bauen.
\end{textebox}
\end{exercice}

\begin{exercice}[Appariements : textes et situations]
Associe chaque début de texte allemand (1-4) au type de texte qui convient (a-d).
\begin{enumerate}
\item \all{Zu verkaufen: Fahrrad, fast neu, Preis 25 000 FCFA. Tel.: 699 00 00 00.}
\item \all{Sehr geehrte Damen und Herren, ich bewerbe mich um ein Praktikum in Ihrer Firma.}
\item \all{Am Montag fand in Douala eine Konferenz über die Wirtschaft in Zentralafrika statt.}
\item \all{Lieber Herr Ngo, vielen Dank für Ihre Bestellung vom 5. Mai.}
\end{enumerate}
\begin{itemize}
\item[a)] une petite annonce
\item[b)] une lettre de motivation
\item[c)] un compte rendu
\item[d)] une lettre commerciale
\end{itemize}
\end{exercice}

\courssub{Je me prépare au Bac}

\begin{exercice}[Compréhension de l'écrit et questions de langue]
Lis le texte suivant, puis réponds aux questions.

\begin{textebox}[Eine deutsche Firma in Kamerun]
Seit fünf Jahren arbeitet die deutsche Firma Holzmann in Douala. Sie kauft Holz und Kakao von kamerunischen Bauern und exportiert diese Produkte nach Europa. Die Firma ist größer als viele andere Unternehmen in der Region und beschäftigt heute mehr als zweihundert Arbeiter. „Wenn wir mehr junge Leute ausbilden, wird die Wirtschaft hier stärker“, sagt der Direktor, Herr Keller. Deshalb bietet die Firma jedes Jahr Praktika für Schüler und Studenten an. Viele Jugendliche lernen dort, wie man Maschinen repariert und wie man ein Geschäft führt. Nächstes Jahr wird die Firma eine neue Fabrik in Yaoundé eröffnen. Dort wird sie noch mehr Arbeitsplätze schaffen. Die Bauern sind zufrieden, denn sie bekommen bessere Preise als früher. Wenn die Zusammenarbeit so weitergeht, werden beide Seiten profitieren.
\end{textebox}

\textbf{I. Questions de compréhension} — Réponds en allemand par des phrases complètes.
\begin{enumerate}
\item Was kauft die Firma Holzmann in Kamerun ?
\item Wie viele Arbeiter beschäftigt die Firma ?
\item Was bietet die Firma den Schülern und Studenten an ?
\item Was wird die Firma nächstes Jahr machen ?
\item Warum sind die Bauern zufrieden ?
\end{enumerate}

\textbf{II. Questions de langue}
\begin{enumerate}
\item Relève dans le texte deux comparatifs et donne leur adjectif de base.
\item Relève une subordonnée avec \all{wenn} et explique la place du verbe conjugué.
\item Relève deux phrases au futur et souligne l'auxiliaire et l'infinitif.
\item Trouve dans le texte le contraire de : \all{kleiner, schlechter, früher}.
\end{enumerate}
\end{exercice}

\begin{exercice}[Thème et version type Baccalauréat]
\textbf{A. Thème} — Traduis les cinq phrases suivantes en allemand. (Chaque phrase mobilise le comparatif, la subordonnée avec \all{wenn} ou le futur.)
\begin{enumerate}
\item Si je trouve un emploi, je gagnerai plus d'argent que mon frère.
\item Les produits camerounais sont meilleurs que les produits importés.
\item Quand j'aurai fini mes études, je travaillerai dans une entreprise à Douala.
\item Si les prix montent, les clients achèteront moins de marchandises.
\item Le football est le sport le plus populaire au Cameroun.
\end{enumerate}

\textbf{B. Version} — Traduis le texte suivant en français.

\begin{textebox}[Texte à traduire]
Mein Onkel hat ein kleines Geschäft in Bafoussam. Er verkauft Handys und Computer. Sein Geschäft ist kleiner als die großen Läden in Douala, aber er hat viele Kunden. Wenn die Schule anfängt, kommen viele Eltern und kaufen Sachen für ihre Kinder. Nächstes Jahr wird mein Onkel ein zweites Geschäft eröffnen. Dann wird er mehr Geld verdienen und zwei junge Leute aus dem Viertel einstellen. Ich bin stolz auf ihn, denn er hat alles allein aufgebaut.
\end{textebox}
\end{exercice}

\begin{exercice}[Sujet de rédaction type Baccalauréat]
Choisis \textbf{un} des deux sujets suivants et rédige ton texte en allemand.

\textbf{Sujet 1 : Petite annonce.} Tu veux vendre ton téléphone portable (ou ton vélo, ou tes livres scolaires). Rédige une petite annonce en allemand (8 à 10 lignes) en respectant la structure et le style télégraphique : objet, état, prix en FCFA, lieu, moyen de contact. Utilise au moins un comparatif.

\textbf{Sujet 2 : Lettre de motivation.} L'entreprise allemande Holzmann à Douala offre un stage (\all{ein Praktikum}) pendant les grandes vacances. Rédige une lettre de motivation courte en allemand (10 à 12 lignes) : en-tête avec \all{Sehr geehrte Damen und Herren}, présentation, motivation, qualités, formule de politesse finale \all{Mit freundlichen Grüßen}. Utilise au moins une phrase au futur et une subordonnée avec \all{wenn}.
\end{exercice}

\newpage

\coursec{Corrigés des exercices}

\begin{corrige}[Exercice 1 — QCM : grammaire et vie économique]
\begin{enumerate}
\item \textbf{Bonne réponse : a)} \all{größer}. Le comparatif de \all{groß} se forme avec le suffixe \all{-er} et l'Umlaut : \all{groß → größer}. \all{Mehr groß} est un calque du français « plus grand » ; \all{am größten} est le superlatif, or la phrase contient \all{als}, qui exige le comparatif.
\item \textbf{a)} \all{habe}. La phrase est au conditionnel réel (présent + futur implicite) : \all{Wenn ich mehr Geld habe, kaufe ich einen Computer.} \all{Hätte} (Konjunktiv II) exigerait \all{würde ich kaufen} ; \all{hat} ne s'accorde pas avec \all{ich}.
\item \textbf{b)} \all{Zu verkaufen:} (« à vendre »). \all{Liebe Frau Mballa} ouvre une lettre personnelle, \all{Sehr geehrte Damen und Herren} une lettre formelle.
\item \textbf{a)} \all{werde}. Futur I : \all{werden} + infinitif en fin de phrase : \all{Nächstes Jahr werde ich in Douala arbeiten.} \all{Will} exprime une volonté, pas un futur pur.
\end{enumerate}
\end{corrige}

\begin{corrige}[Exercice 2 — Vrai ou faux ?]
\begin{enumerate}
\item \textbf{Vrai.} \all{Wenn} est une conjonction de subordination : le verbe conjugué se rejette en dernière position de la subordonnée : \all{Wenn ich Zeit habe, \ldots}
\item \textbf{Faux.} Le superlatif de \all{gut} est irrégulier : \all{gut – besser – am besten}. \all{Gutest} n'existe pas.
\item \textbf{Vrai} (dans le cadre scolaire). La lettre formelle allemande se conclut par la formule de politesse \all{Mit freundlichen Grüßen} (abréviation \all{MfG}), suivie de la signature. (Nuance : on peut ajouter un nom de ville, mais la formule reste la norme.)
\item \textbf{Faux.} Au Futur I, l'auxiliaire \all{werden} occupe la deuxième position et l'infinitif se place \textbf{à la fin} de la phrase : \all{Ich werde in Douala arbeiten.}
\item \textbf{Faux.} \all{Die Firma} (pluriel \all{die Firmen}) signifie « l'entreprise, la société ». « La ferme » se dit \all{der Bauernhof} : c'est un faux ami classique.
\end{enumerate}
\end{corrige}

\begin{corrige}[Exercice 3 — Vocabulaire : la vie économique]
\begin{enumerate}
\item Die \textbf{Waren} auf dem Markt sind heute sehr frisch. (Attention : sujet pluriel \all{sind} → pluriel de \all{die Ware}.)
\item Mein Bruder hat einen neuen \textbf{Arbeitsplatz} in einer Bank in Douala gefunden. (Accusatif masculin : \all{einen neuen Arbeitsplatz}.)
\item Der \textbf{Preis} für Kakao ist dieses Jahr gestiegen.
\item Der \textbf{Kunde} bezahlt an der Kasse.
\item Nach dem Essen im Restaurant bekommen wir die \textbf{Rechnung}.
\item Das \textbf{Geschäft} öffnet jeden Morgen um acht Uhr.
\end{enumerate}
\textbf{Points de vigilance :} majuscule à tous les noms ; bien vérifier l'accord du verbe avec le sujet choisi (singulier/pluriel).
\end{corrige}

\begin{corrige}[Exercice 4 — Conjugaison à compléter]
\begin{center}
\begin{tabular}{|l|l|l|}
\hline
\rowcolor{popblueL} Personne & Präsens & Futur I \\ \hline
ich (arbeiten) & arbeite & werde arbeiten \\ \hline
du (verdienen) & verdienst & wirst verdienen \\ \hline
er/sie (kaufen) & kauft & wird kaufen \\ \hline
wir (verkaufen) & verkaufen & werden verkaufen \\ \hline
sie/Sie (gründen) & gründen & werden gründen \\ \hline
\end{tabular}
\end{center}
\textbf{Rappels :} \all{arbeiten} prend un \all{-e-} de liaison à la 2\up{e} et 3\up{e} personne du singulier (\all{du arbeitest, er arbeitet}) ; au futur, l'infinitif se place en fin de phrase, sans \all{zu}.
\end{corrige}

\begin{corrige}[Exercice 5 — Thème : phrases à traduire]
\begin{enumerate}
\item \all{Der Markt in Douala ist größer als der Markt in Bafoussam.} (Comparatif + \all{als} ; \all{groß → größer} avec Umlaut.)
\item \all{Wenn ich Geld habe, werde ich ein Fahrrad kaufen.} (Subordonnée \all{wenn} : verbe conjugué à la fin ; principale inversée : \all{werde ich} ; infinitif \all{kaufen} en fin de phrase.)
\item \all{Meine Schwester arbeitet in einer Bank in Yaoundé.} (Datif féminin : \all{in einer Bank}.)
\item \all{Kakao ist das wichtigste Produkt unserer Region.} (Superlatif attribut : \all{das wichtigste Produkt} ; génitif : \all{unserer Region}.)
\item \all{Nächstes Jahr werden wir eine deutsche Firma besuchen.} (Futur : \all{werden} en 2\up{e} position, infinitif \all{besuchen} à la fin.)
\end{enumerate}
\textbf{Points de vigilance :} ne jamais traduire « plus \ldots{} que » par \all{mehr} + adjectif simple ; « que » de comparaison = \all{als}, jamais \all{wie}.
\end{corrige}

\begin{corrige}[Exercice 6 — Transformation : du présent au futur]
\begin{enumerate}
\item \all{Ich werde Deutsch an der Schule lernen.}
\item \all{Mein Freund wird eine Stelle in einer Fabrik suchen.}
\item \all{Wir werden unseren Kakao an eine Firma verkaufen.}
\item \all{Die Wirtschaft in Kamerun wird wachsen.}
\item \all{Du wirst ein Praktikum in Douala machen.}
\end{enumerate}
\textbf{Méthode :} on remplace le verbe conjugué par la forme de \all{werden} correctement accordée (2\up{e} position), et le verbe lexical passe à l'infinitif, renvoyé en fin de phrase. Le reste de la phrase ne change pas.
\end{corrige}

\begin{corrige}[Exercice 7 — Texte à trous : un compte rendu économique]
Texte complété :

\all{Der Markt in Mvog-Mbi ist heute \textbf{voller} als gestern. Viele Händler verkaufen ihre Waren \textbf{billiger} als letzte Woche. Wenn das Wetter gut \textbf{ist}, kommen mehr Kunden. Die Preise für Tomaten sind am Morgen am \textbf{niedrigsten}. Wenn die Händler mehr verdienen \textbf{wollen}, müssen sie früh aufstehen. Nächstes Jahr \textbf{wird} die Stadt einen neuen Markt bauen.}

\textbf{Justifications :}
\begin{itemize}
\item \all{voller}, \all{billiger} : comparatifs devant \all{als} (\all{voll} et \all{billig} ne prennent pas d'Umlaut).
\item \all{ist} : dans la subordonnée \all{wenn}, le verbe conjugué va en fin de proposition.
\item \all{am niedrigsten} : superlatif adverbial avec \all{am \ldots{} -sten}. Le radical \all{niedrig} se termine en \all{-g} : on ajoute simplement \all{-sten}, sans \all{-e-} de liaison (celui-ci n'apparaît qu'après \all{-s, -ß, -z, -d, -t} : \all{am heißesten, am kürzesten}).
\item \all{wollen} : verbe conjugué en fin de subordonnée \all{wenn} ; la principale suit avec inversion : \all{müssen sie}.
\item \all{wird} : futur avec sujet singulier \all{die Stadt} ; infinitif \all{bauen} en fin de phrase.
\end{itemize}
\end{corrige}

\begin{corrige}[Exercice 8 — Appariements : textes et situations]
\begin{center}
\begin{tabularx}{\linewidth}{|c|X|}
\hline
\rowcolor{popblueL} Texte & Type de texte \\ \hline
1 & a) petite annonce : style télégraphique, objet + prix + téléphone. \\ \hline
2 & b) lettre de motivation : \all{ich bewerbe mich um \ldots} (« je pose ma candidature »). \\ \hline
3 & c) compte rendu : fait rapporté au Präteritum (\all{fand \ldots{} statt}), ton informatif. \\ \hline
4 & d) lettre commerciale : remerciement pour une commande (\all{Bestellung}). \\ \hline
\end{tabularx}
\end{center}
\textbf{Astuce de reconnaissance :} \all{Zu verkaufen} → annonce ; \all{sich bewerben um} → candidature ; \all{fand statt} → compte rendu ; \all{Bestellung} → correspondance commerciale.
\end{corrige}

\begin{corrige}[Exercice 9 — Compréhension de l'écrit et questions de langue]
\textbf{I. Questions de compréhension} (réponses modèles)
\begin{enumerate}
\item \all{Die Firma Holzmann kauft Holz und Kakao von kamerunischen Bauern.}
\item \all{Sie beschäftigt heute mehr als zweihundert Arbeiter.}
\item \all{Sie bietet ihnen jedes Jahr Praktika an.} (Verbe à particule : \all{anbieten} → \all{bietet \ldots{} an}.)
\item \all{Nächstes Jahr wird sie eine neue Fabrik in Yaoundé eröffnen.}
\item \all{Die Bauern sind zufrieden, weil sie bessere Preise als früher bekommen.} (ou avec \all{denn} : \all{denn sie bekommen bessere Preise als früher}.)
\end{enumerate}

\textbf{II. Questions de langue}
\begin{enumerate}
\item Comparatifs : \all{größer} (base : \all{groß}) ; \all{bessere} (base : \all{gut}, comparatif irrégulier \all{besser}).
\item Subordonnée : \all{Wenn wir mehr junge Leute ausbilden, wird die Wirtschaft hier stärker.} Le verbe conjugué \all{ausbilden} se place à la fin de la subordonnée ; comme il est conjugué en position finale, la particule \all{aus-} reste soudée au verbe. La principale commence par le verbe \all{wird}.
\item Phrases au futur : \all{Nächstes Jahr \textbf{wird} die Firma eine neue Fabrik in Yaoundé \emph{eröffnen}.} ; \all{Dort \textbf{wird} sie noch mehr Arbeitsplätze \emph{schaffen}.} (Auxiliaire \all{wird} en 2\up{e} position, infinitif en fin.)
\item Contraires : \all{kleiner} → \all{größer} ; \all{schlechter} → \all{besser} ; \all{früher} → \all{später} (dans le texte, \all{als früher} oppose passé et présent ; le contraire lexical est \all{später}).
\end{enumerate}
\textbf{Barème indicatif (20 pts) :} compréhension 5 × 2 pts ; questions de langue 4 × 2,5 pts.
\end{corrige}

\begin{corrige}[Exercice 10 — Thème et version type Baccalauréat]
\textbf{A. Thème} (traductions modèles)
\begin{enumerate}
\item \all{Wenn ich eine Stelle finde, werde ich mehr Geld verdienen als mein Bruder.} (« un emploi » = \all{eine Stelle} ou \all{einen Arbeitsplatz} ; comparatif \all{mehr \ldots{} als}.)
\item \all{Die kamerunischen Produkte sind besser als die importierten Produkte.} (Comparatif irrégulier \all{gut → besser}.)
\item \all{Wenn ich mein Studium beendet habe, werde ich in einer Firma in Douala arbeiten.} (En allemand, le futur antérieur français se rend par le Perfekt dans la subordonnée : \all{beendet habe} à la fin.)
\item \all{Wenn die Preise steigen, werden die Kunden weniger Waren kaufen.} (Verbe conjugué \all{steigen} en fin de subordonnée ; principale au futur avec inversion.)
\item \all{Fußball ist der beliebteste Sport in Kamerun.} (Superlatif attribut : \all{der beliebteste Sport} ; \all{der Fußball} avec article est également correct.)
\end{enumerate}

\textbf{B. Version} (traduction modèle)

« Mon oncle a un petit magasin à Bafoussam. Il vend des téléphones portables et des ordinateurs. Son magasin est plus petit que les grandes boutiques de Douala, mais il a beaucoup de clients. À la rentrée scolaire, beaucoup de parents viennent et achètent des affaires pour leurs enfants. L'année prochaine, mon oncle ouvrira un deuxième magasin. Alors il gagnera plus d'argent et embauchera deux jeunes gens du quartier. Je suis fier de lui, car il a tout construit tout seul. »

\textbf{Points de vigilance :} \all{die Sachen} = « les affaires » (pas « les choses » au sens abstrait) ; \all{einstellen} = « embaucher » ; \all{stolz auf jemanden sein} = « être fier de quelqu'un » ; \all{aufbauen} = « construire, monter ».

\textbf{Barème indicatif (20 pts) :} thème 5 × 2 pts (grammaire 1 pt + vocabulaire 1 pt) ; version 10 pts (fidélité 6 pts, qualité du français 4 pts).
\end{corrige}

\begin{corrige}[Exercice 11 — Sujet de rédaction type Baccalauréat]
\textbf{Sujet 1 — Petite annonce (production modèle)}

\all{Zu verkaufen: Handy Samsung, fast neu, nur sechs Monate alt. Das Handy ist billiger als ein neues Modell, aber es funktioniert sehr gut. Farbe: schwarz, mit Ladegerät und Hülle. Preis: 35 000 FCFA (verhandelbar). Ort: Yaoundé, Quartier Mvog-Mbi. Bei Interesse bitte anrufen: 699 12 34 56 (Jean). Auch per WhatsApp erreichbar.}

Traduction : « À vendre : téléphone Samsung, presque neuf, seulement six mois d'usage. Le téléphone est moins cher qu'un modèle neuf, mais il fonctionne très bien. Couleur : noir, avec chargeur et coque. Prix : 35 000 FCFA (négociable). Lieu : Yaoundé, quartier Mvog-Mbi. Si intéressé, appeler : 699 12 34 56 (Jean). Joignable aussi par WhatsApp. »

\textbf{Sujet 2 — Lettre de motivation (production modèle)}

\all{Yaoundé, den 15. Juni

Sehr geehrte Damen und Herren,

ich heiße Amina Etoundi und ich bin Schülerin der Klasse Première am Gymnasium Yaoundé. Ich bewerbe mich um ein Praktikum in Ihrer Firma während der großen Ferien. Ich lerne seit drei Jahren Deutsch und ich interessiere mich sehr für die Wirtschaft. Wenn ich bei Ihnen arbeite, werde ich pünktlich und motiviert sein. Ich bin fleißig, freundlich und ich arbeite gern im Team. Nächstes Jahr werde ich mein Abitur machen, und dieses Praktikum ist sehr wichtig für meine Zukunft.

Über eine positive Antwort freue ich mich sehr.

Mit freundlichen Grüßen

Amina Etoundi}

\textbf{Barème indicatif (20 pts) :} respect du type de texte et de la structure 4 pts ; correction grammaticale (comparatif, \all{wenn}-Satz, futur, ordre des mots) 6 pts ; richesse du vocabulaire économique 4 pts ; orthographe et majuscules 3 pts ; présentation et longueur 3 pts.

\textbf{Points de vigilance :}
\begin{itemize}
\item Petite annonce : style télégraphique admis (phrases nominales), mais le comparatif exigé doit être correct (\all{billiger als}, jamais \all{mehr billig}).
\item Lettre : \all{Sehr geehrte Damen und Herren,} avec virgule, puis ligne suivante en minuscule ; formule finale \all{Mit freundlichen Grüßen} sans virgule.
\item Vérifier la place du verbe dans chaque subordonnée \all{wenn} et l'infinitif final au futur.
\item Majuscules à tous les noms : \all{Praktikum, Ferien, Firma, Abitur, Zukunft}.
\end{itemize}
\end{corrige}

\newpage

\coursec{Fiche bilan}

\begin{popfigure}
\begin{tikzpicture}[mindmap, grow cyclic, every node/.style={concept}, root concept/.append style={concept color=popblue, font=\large\bfseries}, level 1/.append style={level distance=3.4cm, sibling angle=60}, level 2/.append style={level distance=2.2cm, sibling angle=45}]
\node[root concept, align=center]{Vie économique\\ ALA09}
  child[concept color=poporange]{ node[align=center]{Lexique\\ Wirtschaft}
    child { node[concept color=poporangeL, font=\small, align=center]{der Markt\\ die Ware} }
    child { node[concept color=poporangeL, font=\small, align=center]{kaufen /\\ verkaufen} }
  }
  child[concept color=popgreen]{ node[align=center]{Comparatif\\ Superlatif}
    child { node[concept color=popgreenL, font=\small, align=center]{billiger als\\ am billigsten} }
  }
  child[concept color=poppurple]{ node {wenn-Sätze}
    child { node[concept color=poppurpleL, font=\small, align=center]{verbe\\ à la fin} }
  }
  child[concept color=poppink]{ node {Futur}
    child { node[concept color=poppinkL, font=\small, align=center]{werden +\\ Infinitiv} }
  }
  child[concept color=popgold]{ node[align=center]{Types de\\ textes}
    child { node[concept color=popgoldL, font=\small, align=center]{Annonce\\ Brief\\ Bericht} }
  }
  child[concept color=popblue!60]{ node[align=center]{Traduction\\ thème / version}
    child { node[concept color=popblueL, font=\small, align=center]{ordre des\\ mots, cas} }
  };
\end{tikzpicture}
\legende{Carte mentale de la leçon : production et traduction de textes sur la vie économique.}
\end{popfigure}

\begin{bilan}
\textbf{1. Lexique clé de la vie économique}
\begin{itemize}
  \item \all{die Wirtschaft, -en} : l'économie ; \all{der Markt, die Märkte} : le marché ; \all{das Geschäft, -e} : le magasin, l'affaire.
  \item \all{der Preis, -e} : le prix ; \all{die Ware, -n} : la marchandise ; \all{das Geld} : l'argent.
  \item \all{kaufen / verkaufen} : acheter / vendre ; \all{bezahlen} : payer ; \all{kosten} : coûter ; \all{verdienen} : gagner (de l'argent).
  \item \all{billig / teuer} : bon marché / cher ; \all{der Kunde, -n / die Kundin, -nen} : le client / la cliente.
  \item \all{die Rechnung, -en} : la facture ; \all{der Rabatt, -e} : la remise ; \all{die Anzeige, -n} : l'annonce.
\end{itemize}

\textbf{2. Grammaire à connaître par cœur}
\begin{center}
\begin{tabularx}{\linewidth}{|l|X|X|}
\hline
\rowcolor{popblueL} \textbf{Point} & \textbf{Règle} & \textbf{Exemple} \\
\hline
Comparatif & adjectif + \all{-er} + \all{als} & \all{Kakao ist teurer als Kaffee.} \\
\hline
Superlatif & \all{am} + adjectif + \all{-sten} & \all{Gold ist am teuersten.} \\
\hline
Irréguliers & \all{gut $\to$ besser $\to$ am besten} ; \all{viel $\to$ mehr $\to$ am meisten} ; \all{gern $\to$ lieber $\to$ am liebsten} & \all{Ich kaufe lieber Obst.} \\
\hline
\all{wenn}-Satz & subordonnée : verbe conjugué à la fin ; si elle ouvre la phrase, le verbe principal suit en position 2 & \all{Wenn ich Geld habe, kaufe ich Reis.} \\
\hline
Futur & \all{werden} (conjugué) + infinitif à la fin & \all{Ich werde ein Geschäft eröffnen.} \\
\hline
\end{tabularx}
\end{center}

\textbf{3. Méthodes express}
\begin{itemize}
  \item \textbf{Petite annonce} : titre court, produit + prix, contact ; phrases brèves (\all{Verkaufe Fahrrad, sehr gut, 50 Euro. Rufen Sie an!}).
  \item \textbf{Lettre} : lieu et date en haut à droite, formule d'appel (\all{Sehr geehrte Damen und Herren,}), paragraphes clairs, salutation (\all{Mit freundlichen Grüßen}).
  \item \textbf{Compte rendu} : qui ? quoi ? quand ? où ? ; au passé (\all{Perfekt} ou \all{Präteritum} de \all{sein/haben}) ; ton objectif.
  \item \textbf{Thème (fr $\to$ all)} : repérer le verbe, le placer en position 2, mettre les compléments dans l'ordre \emph{temps -- cause -- manière -- lieu}, vérifier les cas et les majuscules.
  \item \textbf{Version (all $\to$ fr)} : identifier sujet / verbe / compléments, traduire les subordonnées en respectant l'ordre français, rendre un français naturel.
  \item \textbf{Relecture} : majuscule à tous les noms, verbe en position 2, particule séparable en fin de phrase, accords et cas.
\end{itemize}
\end{bilan}

\begin{aretenir}[Checklist avant le Bac]
Avant de rendre ma copie, je vérifie :
\begin{itemize}
  \item[$\square$] Je connais le lexique économique avec article et pluriel (\all{der Markt, die Märkte}).
  \item[$\square$] Je forme le comparatif avec \all{-er als} et le superlatif avec \all{am \ldots -sten}.
  \item[$\square$] Je cite les irréguliers : \all{gut -- besser -- am besten}, \all{viel -- mehr -- am meisten}.
  \item[$\square$] Dans une subordonnée avec \all{wenn}, mon verbe conjugué est à la fin.
  \item[$\square$] Si la subordonnée est en tête, le verbe principal vient juste après la virgule.
  \item[$\square$] Je conjugue le futur : \all{werden} + infinitif en fin de phrase.
  \item[$\square$] Je sais rédiger une petite annonce : produit, prix, contact.
  \item[$\square$] Je respecte la structure de la lettre : date, appel, corps, \all{Mit freundlichen Grüßen}.
  \item[$\square$] En thème, je place le verbe en position 2 et les infinitifs/participes en fin.
  \item[$\square$] Je mets une majuscule à tous les noms allemands.
  \item[$\square$] Je vérifie les cas : accusatif après \all{für}, datif après \all{mit, von, zu}.
  \item[$\square$] Je relis ma traduction : elle doit sonner naturel en français correct.
\end{itemize}
\end{aretenir}

\findecours

\end{document}
